\documentclass[a4paper,10pt,openany]{book}
\usepackage[utf8]{inputenc}
\usepackage[czech]{babel}
\usepackage{graphicx}
\usepackage{wrapfig}
\usepackage{longtable}
\usepackage{needspace}
\usepackage{titlesec}
\title{ImproWiki}
\author{wiki.improliga.cz pod licencí Creative Commons Uveďte autora}
\makeatletter
\renewcommand\thesection{}
\renewcommand\thesubsection{}
\makeatother
\usepackage[margin=1.5cm]{geometry}
\usepackage{tgpagella}
\usepackage{multicol}

\usepackage[T1]{fontenc}
\usepackage{fancyhdr}
\pagestyle{fancy}
\fancyhf{}
\fancyhead[LO,LE]{\today}
\fancyhead[CE,CO]{wiki.improliga.cz}
\fancyhead[RE,RO]{\thepage}



\newcommand{\todo}[1]{

TODO: \textbf{#1} 

}

\setlength{\intextsep}{0pt}


\newcommand{\katabox}[3]{
	\p{#1}  
\p{#2}
\p{#3} 


}

\newcommand{\faulbox}[3]{

\begin{wrapfigure}{L}{0}
#1
\end{wrapfigure}

\textbf{Gesto:}\p{#2} 


\textbf{Trestné body:}\p{#3} 


}


\newcommand{\obrazek}[2]{
\includegraphics[width=4cm]{#1}

\textbf{#2}}


\newcommand{\obrazekmaly}[1]{
\includegraphics[width=1.8cm]{#1}}


\newcommand{\odkaz}[2]{\textbf{#1}\textsuperscript{s.~\pageref{#2}}} 

%\footnote{\textbf{#2}, strana \pageref{#2} }

\clubpenalty 10000
\widowpenalty 10000


\titlespacing*{\section}{-5mm}{4mm}{3mm}

\begin{document}
\hyphenation{pros-tře-dí-mi dvo-ji-ce-mi}
\begin{titlepage}
\begin{center}

% Upper part of the page. The '~' is needed because \\
% only works if a paragraph has started.
\includegraphics[width=0.15\textwidth]{./logo}~\\[2cm]

\textsc{\Large wiki.improliga.cz pod licencí Creative Commons Uveďte autora }\\[0.5cm]

% Title

{ \huge \bfseries ImproWiki \\[0.4cm] }




\vfill

% Bottom of the page
{\large Vygenerováno \today}

\end{center}




\end{titlepage}
 
\chapter{Úvod}\label{úvod}

\documentclass[main.tex]{subfiles}\begin{document}
\needspace{5cm} \label{improwiki} \textbf{ImproWiki}{} si klade za cíl zejména popsat běžně hrané \odkaz{improligové}{čili} \odkaz{kategorie}{kategorie} tak, aby \odkaz{hráči}{hráč} z různých \odkaz{týmů}{tým} mohli na \odkaz{zápasech}{zápas} hrát podle stejných pravidel. 
 
Tištěná forma ImproWiki je ke stažení na titulní stránce wiki. 

\end{document} 

Oproti verzi dostupné na internetu zde chybí odkazy na videa a další internetové zdroje.

Aktuálnost a užitečnost tištěné verze (ani wiki) není zaručena, v případě nejasností konzultujte se zkušenějším hráčem či s improwiki.

Tučně vysázená \odkaz{hesla}{úvod} říkají, na které straně se o daném termínu dočtete víc.

Některé termíny ještě nejsou popsané, přihlašte se na http://wiki.improliga.cz/ a dopište je.

\textbf{Varování:} Nikdo se ještě nenaučil improvizovat z knížky. Popis kategorií je určen spíše zkušeným hráčům pro připomenutí, než nováčkům na naučení se.  

\LaTeX{} sazbu připravuje  \odkaz{Václav Černý}{uživatel:vatoz}.

Chyby či náměty na vylepšení konverze hlaste na https://github.com/vatoz/improwikibook/issues

Případné faktické chyby a překlepy můžete opravit přímo na improwiki.
\chapter{Zápas}
\needspace{5cm} \label{zápas} Zápas v divadelní improvizaci je unikátní divadelní forma, ve které se dva \odkaz{týmy}{tým} utkají v improvizačních \odkaz{kategoriích}{kategorie}.  
Forma je částečně inspirována hokejem. 
 
\section{ Celkový průběh z pohledu diváka }  
\subsection{ Před zahájením } \begin{itemize}
\item  Diváci při příchodu mají možnost napsat na lístečky \odkaz{témata}{askfor}
\item  Pokud není zařízeno jinak, obvykle diváci dostávají \odkaz{hlasovací kartičky}{hlasovací kartička}
\item  Diváci dostávají \odkaz{papuče}{papuče} nebo \odkaz{míčky}{papuče} na vyjádření nesouhlasu s rozhodčím
\end{itemize}
 
\subsection{ Úvod } \odkaz{Konferenciér}{konferenciér}  
\begin{enumerate}
\item  přivítá diváky, trochu je \odkaz{rozehřeje}{rozehřívání diváků} a vysvětlí jim na co vlastně přišli a princip improvizačního zápasu.
\item  zve \odkaz{hráče}{hráč} na \odkaz{veřejnou rozcvičku}{veřejná rozcvička}, hráči jsou ještě v nediferencovaných trikotech zpravidla černých, proběhne několik cvičení, hráči následně odbíhají.
\item  dál informuje diváky, např. o \odkaz{hlasování}{hlasování}, možnosti vyjádření nesouhlasu papučí či míčkem a \odkaz{faulech}{faul}
\end{enumerate}
 
\subsection{ Začátek zápasu } \odkaz{Konferenciér}{konferenciér}  
\begin{enumerate}
\item  zve jednotlivé \odkaz{hráče}{hráč} v \odkaz{týmech}{tým} a představí je, případné zmíní, že některý improvizátor hraje poprvé zápas
\item  zve \odkaz{pomocňáky}{pomocný rozhodčí} (pomocné rozhodčí)
\item  zve \odkaz{hlavního rozhodčího}{rozhodčí} a předává mu hlavní slovo
\end{enumerate}
 
\subsection{ První polovina } \begin{enumerate}
\item  Pokřiky týmů (jediná nacvičená věc)
\item  \odkaz{Odehraje}{zápas}  se několik \odkaz{kategorií}{kategorie}
\item  Přestávka
\end{enumerate}
 
\subsection{ Druhá polovina } \begin{enumerate}
\item  Výměna stran týmů
\item  Voliteně pokřiky a veřejná rozcvička na druhou půli
\item  \odkaz{Hraje}{zápas} se několik dalších \odkaz{kategorií}{kategorie}
\end{enumerate}
 
\subsection{ Závěr } \begin{enumerate}
\item  Ukončení zápasu a vyhlášení vítěze podle skóre
\item  Představení \odkaz{hráčů}{hráč}, \odkaz{pomocných rozhodčích}{pomocný rozhodčí}, \odkaz{hlavního rozhodčího}{rozhodčí}, \odkaz{hudebníka}{hudebník}, \odkaz{konferenciéra}{konferenciér} a ostatních
\item  Pozvánka na další zápasy
\item  Klanění
\item  Volitelně přídavek
\end{enumerate}
 
\section{ Průběh z pohledu účastníků } \begin{itemize}
\item  \odkaz{Předzápasový trénink}{předzápasový trénink}
\item  Domluva s \odkaz{rozhodčím}{rozhodčí}, které \odkaz{kategorie}{kategorie} se budou hrát, případně která varianta.
\item  \odkaz{Zápas (divácká část)}{zápas}
\item  Volitelně loučení s odcházejícími diváky
\item  \odkaz{pozápasová reflexe}{reflexe}
\item  Volitelně osobní neřízená reflexe v restauračním zřízení
\end{itemize}
 
\section{ Průběh kategorie } \begin{enumerate}
\item  \odkaz{Rozhodčí}{rozhodčí} si vyžádá \odkaz{askfor}{askfor} ke kategorii a určí která kategorie se bude hrát, určí čas na přípravu a čas trvání kategorie
\item  \odkaz{Konferenciér}{konferenciér} stručně vysvětlí pravidla dané kategorie
\item  Hraje se \odkaz{kategorie}{kategorie}. \odkaz{Rozhodčí}{rozhodčí} má možnost v průběhu pískat \odkaz{fauly}{faul}. \odkaz{Pomocný rozhodčí}{pomocný rozhodčí} ukazuje hráčům gesty zbývající čas kategorie.
\item  Kategorie je ukončena. \odkaz{Rozhodčí}{rozhodčí} přiděluje \odkaz{trestné body}{trestný bod} za \odkaz{fauly}{faul}. Může si pozvat kapitány a nebo přímo \odkaz{hráče}{hráč}, kterým byl \odkaz{faul}{faul} písknut, aby odůvodnili svůj faul odůvodnili, případně se z něj vykecali. Diváci mají možnost vyjádřit nesouhlas s rozhodčím, např hodem \odkaz{papučí}{papuče}, míčkem či bučením.
\item  Probíhá \odkaz{hlasování}{hlasování} diváků a přidělení vítězných bodů.
\end{enumerate}
 
\section{ Kam dál } \begin{itemize}
\item  \odkaz{Zápasové kategorie}{:kategorie:zápasové kategorie}
\item  \odkaz{Příprava zápasu}{příprava zápasu}
\end{itemize}
 
 
 
 







\chapter{Postavy}
\documentclass[main.tex]{subfiles}\begin{document}
\needspace{5cm} \section{Hráč} \label{hráč} Označení pro jedince na scéně, kteří jsou zodpovědní za tvorbu obsahu improvizace. Mezi hráče nepatří \odkaz{MC}{mc} a \odkaz{hudebníci}{hudebník}. V rámci improvizačních představení se preferuje výraz "\textbf{hráč}{}"{} nad výrazem "\textbf{herec}{}"{} i když mají téměř stejný význam. Všichni jedinci na scéně jsou herci, ale hráči jsou jedinci zastupující vždy svůj tým (zápas) či jsou vyčleněni od \odkaz{MC}{mc} (improshow).  
 
\subsection{ Zápasová forma } Na \odkaz{zápasech}{zápas} jsou hráči rozděleni do \odkaz{týmů}{tým}. 
 
\subsection{ Zápasový oděv } Na veřejnou rozcvičku jsou hráči oděni čistě v černém (bez výšivek, potisků apod),  
na samotný zápas si pak berou trika v barvách svých týmů. 
Trikot tedy zahrnuje 
\begin{itemize}
\item ponožky a kalhoty   případně leginy/punčocháče (40 DEN a víc) a sukni v barvě černé
\item triko v barvě černé
\item triko v barvě týmu.
\end{itemize}
V případě dlouhých vlasů je takřka nutnost gumička do vlasů, jinak nebývá skrz bujnou hřívu vidět do obličeje. 
 
 
 
 
 
\needspace{5cm} \section{Hudebník} \label{hudebník} Hudebník je skvělým oživením každého zápasu či improshow. 
Jeho role je proměnlivá, při \odkaz{hlasování}{hlasování} na zápasech hraje výrazně až freneticky. 
Ve všech \odkaz{zpívaných kategoriích}{:kategorie:zpívané kategorie} je určovatelem/hybatelem charakteru písní. 
A v dalších kategoriích může podkreslovat. 
 
Skvělé je mít několik nástrojů (či klávesy s přepínáním rejstříků) a volit dle charakteru děje/ hlasu hráčů, či jen tak pro zpestření. 
Je vhodné, když, zejména pro méně hudebně zkušené hráče, zvýrazní první dobu a případně i ukáže kývnutím. 
 
\subsection{Role hudebníka} Na hudebníka může být mnoho různých požadavků. Záleží na oraganizátorovi představení, formátu představení, zkušenostech hudebníka,  zkušenostech herců a pod. Jeho role může být: 
\begin{itemize}
\item \textbf{Hudební podkres mimo improvizaci}{} - Před začátkem představení, během úvodu, vyhlášení pauzy, hlasování, ukončení přestavení. Žánr záleží na cílené atmosféře celého představení, ale obvykle by měla být hudba pozitivní, příjemná, veselá až radostná. Pokud jde o jasné oznámení začátku nebo konce, je vhodná i fanfára.
\item \textbf{Podkresová hudba}{} - Slouží pro doplnění nálady a atmosféry během improvizace. Neměla by brát \odkaz{fokus}{fokus} hercům.
\item \textbf{Krátké znělky (jingly)}{} - Mezi větami či souslovími, například v kategoriích \odkaz{1000 způsobů jak}{1000 způsobů jak}, \odkaz{Věty, které by neměly zaznít}{věty, které by neměly zaznít} nebo v rozcvičce \odkaz{Jsem, beru}{jsem, beru}.
\item \textbf{Doprovod písně}{} - Píseň by měla začít jasnou předehrou, ve které hudebník deklaruje tóninu, (jednoduchý) rytmus, tempo a ideálně i (jednoduchý) akordový postup, kterého se po zbytek písně více méně drží. Možnosti variací jsou velmi závislé na schopnostech herce a na jeho souhře s hudebníkem. Experimentování se nedoporučuje, pokud si hudebník není opravdu jistý hercovými schopnostmi.
\item \textbf{Zvukové efekty}{} - Klepání na dveře, vyzvánení telefonu nebo třeba pípající přístroj v nemocnici improvizaci okoření.
\item \textbf{Ukončení improvizace}{} - Improvizaci může ukončit i hudebník. Zvláště, pokud byla právě řečena skvělá pointa na závěr, je to chvíle, kdy dát hercům případně \odkaz{MC}{mc} hudbou najevo, že je vhodný okamžik improvizaci ukončit. V písních má hudebník na starosti ukončování ve většině případů.
\item \textbf{Nový impulz do improvizace}{} - Ve vhodné chvíli, může hudebník výrazně změnit styl doprovodu a přinést tak do improvizace impulz, kterého se herci chytí. Je důležité tuto změnu udělat tak výrazně, aby si toho herci všimli. Vhodná situace je například ve chvíli, kdy se v příběhu dlouho nic neděje.
\item \textbf{Umět nehrát}{} - Do některých kategorií se hudba hodí více, do některých méně. Vždy je ale důležité nezapomínat na využití ticha a na přehlcení hudbou.
\end{itemize}
 
\subsection{Schopnosti a dovednosti} \begin{itemize}
\item \textbf{Poznat, kdy hrát a kdy nehrát}{}
\item '''Hrát v \odkaz{náladě}{hudební nálady}, která se hodí''' nebo aspoň neodporuje
\item \textbf{Umět improvizovat.}{} Naučený repertoir má své limity.
\item \textbf{Umět zahrát v mnoha náladách}{}, nejen \textif{veselý - smutný}{}
\item '''Umět zahrát v mnoha hudebních žánrech''' a doprovodit mnoho filmových a literárních žánrů
\item \textbf{Umět zahrát improvizovanou píseň}{}, která je čitelná pro herce
\item '''Umět pomoci hercům dodržet rytmus, tóninu''', vrátit se zpátky do rytmu, tóniny, nebo se podle herců přizpůsobit.
\item \textbf{Neplést herce složitostmi}{}
\end{itemize}
 
\subsection{Kdy hrát a kdy nehrát} Některé kategorie hudbu přímo vyžadují, v některých by hudba překážela. Je dobré si hudební doprovod dvakrát rozmyslet v kategoriích, které jsou založené na zvucích odjinud, jako \odkaz{Zvuky}{zvuky} nebo \odkaz{Dabing filmu}{dabing filmu}. Dále se hudba příliš nehodí do kategorií s rychlým spádem (\odkaz{Schizofrenie}{schizofrenie}) a obecně kategorií náročných na pozornost (\odkaz{Sportovní komentátor}{sportovní komentátor}). Do kategorií založených na gagování (\odkaz{1000 způsobů jak}{1000 způsobů jak}, \odkaz{Metafory}{metafory}) se hodí pouze oddělování znělkou (jinglem), do zpívaných kategorií je zase vhodné nehrát mezi písničkami, aby nebylo hudby příliš. Hudba se naopak hodí pro kategorie, které vybízejí ke dlouhému příběhu, ke kategoriím, kde se střídají prostředí (\odkaz{Poslední věta}{poslední věta}) nebo úhly pohledů (\odkaz{Spoon river}{spoon river}) pro jasné odlišení. Hudební podkres se dá použít ke gradaci v kategoriích, které jinak hudbu příliš nevyžadují (\odkaz{Polovina času}{polovina času}, \odkaz{Smrt v 1 minutě}{smrt v 1 minutě}). Hudbou se dá vyplnit prázdný prostor nebo naopak podtrhnout silný okamžik, je ale dobré rozvrhnout si množství hudby do celého představení, aby hudba nehrála pořád. V \odkaz{Longformách}{longforma} má hudba na starosti úvod, závěr, předěly a má podkreslovat důležité okamžiky, ale zároveň má nechat prostor na většinu dialogů pouze hercům.  
 
{{Todo| 
\begin{itemize}
\item  Doplnit navyklé způsoby komunikace mezi hudebníkem a hráčem
\end{itemize}
}} 
 
\needspace{5cm} \section{Konferenciér} \label{konferenciér} Konferenciér řídí průběh \odkaz{zápasu}{zápas} či \odkaz{improshow}{improshow}. 
 
\subsection{Běžné chyby} \begin{itemize}
\item rozvleklý úvod
\item chybný či nepřehledný popis kategorií
\end{itemize}
 
 
 
\needspace{5cm} \section{MC} \label{mc} \textbf{MC}{} z anglického \textbf{Master of Ceremony}{}. \odkaz{Konferenciér}{konferenciér}/moderátor/\odkaz{rozhodčí}{rozhodčí}. Prostě ten, kdo má pod pantoflem celý průběh \odkaz{zápasu}{zápas}/\odkaz{improshow}{improshow}. 
Na improshow někdy bývá "kolující MC"{} - \odkaz{hráči}{hráč} se v řízení průběhu střídají. 
 
 
\needspace{5cm} \section{Pomocný rozhodčí} \label{pomocný rozhodčí} \textbf{Pomocný rozhodčí}{}, též \textbf{pomocňák}{} je typicky k vidění na \odkaz{zápasech v improvizaci}{zápas}. 
Bývají dva, jejich úkolem je \odkaz{měřit čas}{ukazování času} kategorií, sčítání \odkaz{hlasů}{hlasování} (\odkaz{kartičky}{hlasovací kartička}), aktualizace skóre a faulů, podávání košíčku s tématy. Jejich pomoc \odkaz{hlavnímu rozhodčímu}{rozhodčí} někdy může zahrnovat zapisování a sledování \odkaz{faulů}{faul}, které rozhodčí pískne.  
 
Pomoci vám může i \odkaz{Manuál pomocného rozhodčího}{manuál pomocného rozhodčího}. 
\subsection{Běžné chyby} \begin{itemize}
\item Pomocňáci si nedohodnou, který sčítá body kterému týmu a sčítání (typicky první na zápase, či první po přestávce) se musí opakovat.
\end{itemize}
 
 
 
\needspace{5cm} \section{Rozhodčí} \label{rozhodčí} \textbf{Rozhodčí}{} je klíčovou postavou improvizačních \odkaz{zápasů}{zápas}, obvykle tou s nejvyším \odkaz{statusem}{status} na scéně. 
Během zápasu zejména vybírá \odkaz{kategorie}{kategorie}, píská \odkaz{fauly}{faul}, uděluje  \odkaz{trestné body}{trestný bod} a  
dále vede průběh zápasu ve spolupráci s \odkaz{konferenciérem}{konferenciér}. 
Pro některá svá nepopulární rozhodnutí bývá terčem hodu \odkaz{papučí}{papuče} či míčků od diváků. 
 
 
Způsoby, kterými může rozhodčí ovlivňovat běh zápasu 
\begin{itemize}
\item Pískání \odkaz{faulů}{faul} a přidělování \odkaz{trestných bodů}{trestný bod}
\item Vyloučení z následující kategorie/kategorií (např. za opakovanou chybu)
\item Výběr témat (\odkaz{Askfor}{askfor})
\item Výběr \odkaz{kategorií}{kategorie}
\item Zadávání dodatečných úkolů \odkaz{hráčům}{hráč}
\item Zastavení improvizace a např. hlasování diváků o dalším postupu
\item Úkolování \odkaz{hudebníka}{hudebník}, \odkaz{pomocných rozhodčí}{pomocný rozhodčí} a \odkaz{konferenciéra}{konferenciér}
\item Další verbální a nonverbální výlevy
\end{itemize}
 
 
\needspace{5cm} \section{Trenér} \label{trenér} Trenérem je myšlen jedinec, soustavně pracující s \odkaz{týmem}{tým} na jeho zdokonalování v improvizaci na pravidelných 
\odkaz{tréninzích}{trénink}. 
Ve starším pojetí \odkaz{zápasů}{zápas} byl rovněž v průběhu zápasu přítomen na scéně a koučoval hráče. 
  

\end{document}


\chapter{Zápasové kategorie}
\label{zápasové kategorie}
\label{:kategorie:zápasové kategorie}
\needspace{5cm} \label{kategorie} Slovo kategorie je nešikovný překlad anglického slovíčka "play"{} (hra). Tímto slovem se rozumí ohraničení začátku a konce nějakého určitého formátu improvizace. Neplést se slovem \odkaz{forma}{forma}.  
 
 
\needspace{5cm} \label{kategorie:zápasové kategorie} Zde se nacházejí kategorie, které se objevují na \odkaz{Improligových}{čili} \odkaz{zápasech}{zápas}. Některé kategorie se hrají častěji, některé jsou již trochu zapomenuté. Před zápasem je vhodné seznam kategorií projít a některé si upřesnit, protože pravidla kategorií se často liší. Zejména se jedná o kategorie: \odkaz{Zpívaná}{zpívaná}, \odkaz{Pyramida}{pyramida}, \odkaz{S rekvizitou}{s rekvizitou}. 
 
U zápasových kategorií je důležité, aby nabízely stejné možnosti oběma týmům. V případě, že se v kategorii vyskytuje speciální jedna role, doporučuje se v polovině času vyměnit hráče. Příkladem může být např. \odkaz{Překlad do znakové řeči}{překlad do znakové řeči}.  
 
 
Některé kategorie najdete také v  \odkaz{kategoriích na improshow}{:kategorie:kategorie na improshow}, či v \odkaz{rozcvičkách}{:kategorie:rozcvičky} 
 
 
 

\needspace{5cm}

\section{Barman song} 
\label{barman song} 



V kategorii \textbf{Barmanský song}{} se hráči vyzpívají ze zadaných problémů a barman jim přívětivě též písní poradí, jak tento problém řešit.  
 
\subsection{ Průběh } \odkaz{MC}{mc} si do diváků vyžádá několik problémů či fóbií a rozdá je hráčům. Improvizace začíná příchodem barmana do baru, který rozehrává zejména pantomimou. Po chvíli přichází do baru první hráč a může i ve svém projevu či jednání zohlednit zadaný problém či fóbii. Proběhne krátká konverzaci na jejímž konci barman podává hráči mikrofon s tím, ať se ze svého problému vyzpívá. Poté, co hráč dozpívá ihned reaguje barman svou radou, kterou též zpívá. Hudba u rady by měla být veselejšího rázu než u problému. Hráč děkuje barmanovi za radu, zaplatí za drink a odchází. Chvilku na to přichází další hráč a situace se opakuje. 
Dá se říci, že \textbf{Barman Song}{} je specializovanou variantou \odkaz{muzikálu}{muzikál}. 
 
\subsection{ Varianty } \begin{itemize}
\item Hráči si problém vymýšlí sami.
\end{itemize}
 
 
\subsection{ Běžné chyby} \begin{itemize}
\item  Hráči je podán vypnutý mikrofon
\item  Hráč zapomene zaplatit útratu
\item  Rada: utop svůj žal v alkoholu.
\end{itemize}
 
 
 
  
  
 
\needspace{5cm} \section{Báseň od publika} \label{báseň od publika} \katabox{Několik(4) slov pro každého hráče 
}{2 
}{cca 6 min} 
 
 
Uslyšíte báseň, na které se sami budete částečně podílet. 
Známé též jako \textbf{Valounská nedělní}{}. 
 
\subsection{Průběh} Dva hráči jdou za dveře či si zacpou uši a \odkaz{MC}{mc} si vyžádá od publika několik slov, která napíše na papíry a položí textem dolů. 
Hráč přijde, oznámí svou báseň (např. Sbírka:Autor:Báseň) a následně jsou mu papíry otočeny a on musí slova co nejrychleji a přitom logicky zapracovat do své básně , 
která by měla být v nějaké logické veršovací struktuře (např \odkaz{ABAB}{abab} či AABB), a zároveň by měla vyhovovat názvu, který hráč anoncoval. 
 
 
\subsection{ Varianty } \begin{itemize}
\item Píseň od publika
\item Karty se otáčí průběžně
\end{itemize}
 
 
 
 
\needspace{5cm} \section{Co tím chceš jako říct} \label{co tím chceš jako říct} \katabox{Téma, prostředí 
}{2+ 
}{volně, více jak 3 minuty} 
 
V kategorii \textbf{Co tím chceš jako říct}{} uvidíte spoustu nedorozumění. 
 
 
\subsection{ Průběh } Volně se rozehrává příběh. \odkaz{Hráči}{hráč} jsou ale v silně výbušných postavách. V momentě, kdy zazní nějaká narážka na citlivé téma od jiné postavy, začínají se slovy "Co tím chceš jako říct?"{} svůj agresivní výlev, ve kterém by na sebe měli, podobně jako při zpěvu v kategorii \odkaz{Mikrofon}{mikrofon}, něco prozradit. Po výlevu scénka pokračuje dál. 
 
\subsubsection{ Příklad výlevu } A: Ta barva zdi se mi moc líbí. 
B: Co tím chceš jako říct? 
B: Že je to špatně natřený? 
B: Tím chceš říct, že bys to natřel líp? 
B: Že jsem jako komunista, když mám červenej pokoj? 
B: Tím chceš říct, že si nevážíš práce soudruhů? 
... 
A: Ne, tím jsem chtěl říct jenom, že je fakt pěkná. 
 
\subsection{ Časté chyby } \begin{itemize}
\item  Výlev není neutrálně ukončen
\end{itemize}
 
 
 
 
 
\needspace{5cm} \section{Dabing ABC} \label{dabing abc} \katabox{Téma 
}{3 
}{neomezený} 
 
V kategorii \textbf{Dabing ABC}{} vám 3 hráči zahrají na Vaše téma, ale tak, že za každého hráče mluví jeho kolega. 
 
\subsection{Průběh} Připraví se tři hráči. Improvizaci rozehrává jeden hráč (A), ale pokud mluví, pouze otevírá pusu a jeho promluvy dabuje druhý hráč (B), který je zatím mimo scénu. Po nějaké chvíli se do děje zapojí hráč B, který je ale dabován třetím hráčem (C). Hráč B stále dabuje hráče A. Nakonec se na scéně objeví i třetí hráč (C), kterého dabuje hráč A. Pro lepší iluzi je dobré, aby hráč, který je právě dabován otočil hlavu do diváků, zatímco jeho dabér mluví do jeviště. Mnohdy vznikají zmatené situace, kdy hráč začne mluvit a jeho kolega si neuvědomí, že by měl otevírat pusu, ale také naopak. Dabovaný hráč může otevřít pusu a jeho dabér musí začít mluvit. Vtipnost této kategorie také spočívá v tom, že dabér může svého svěřence tím, co řekne do jisté míry ovládat. 
 
\subsection{ Varianty } \begin{itemize}
\item Dabing ABCD - hrají 4 hráči.
\item Ostatní hráči mohou ztvárňovat předměty.
\end{itemize}
 
\subsection{ Běžné chyby } \begin{itemize}
\item Hráči nedávají prostor svým kolegům a stále mluví.
\item Hráči neotvírají pusu když jejich dabér mluví.
\end{itemize}
 
 
 
 
 
\needspace{5cm} \section{Dabing filmu} \label{dabing filmu} \katabox{Téma 
}{2-3 
}{omezen ukázkou} 
 
V kategorii \textbf{Dabing filmu}{} je promítána filmová ukázka bez zvuku, kterou hráči předtím neviděli a naživo ji před Vámi dabují. 
 
\subsection{ Průběh } Scéna se připraví na promítání (příprava plátna, projektoru). \odkaz{MC}{mc} určí počet hráčů a jejich pohlaví, podle vybrané filmové ukázky. Hráči dostanou mikrofon, v sále se zhasne a je spuštěna projekce ukázky. Hráči se musí co nejpřesněji strefovat do úst své postavy na plátně a současně vytvářet smysluplný dialog. Kategorie končí s koncem ukázky. Ne ve všech záběrech některá z postav mluví, pak může mluvit libovolná z nich. 
 
Po skončení kategorie je vhodné nechat před hlasováním hráče, aby některou ze svých \odkaz{replik}{replika} představili své postavy. 
 
\subsection{ Varianty } \begin{itemize}
\item Dabing filmu může být pokračováním předešlé kategorie, pokud hráči dojdou do prostředí, o němž \odkaz{MC}{mc} ví, že se v něm odehrává připravená filmová ukázka.
\end{itemize}
 
\subsection{ Běžné chyby } \begin{itemize}
\item Postava na plátně mluví a dabující hráč nemluví a naopak.
\item Hráči dostanou vypnutý mikrofon.
\item Hráči dabují oba stejnou postavu.
\item Promítaná ukázka obsahuje záběry, na nichž většinou není vidět postava, která právě mluví.
\end{itemize}
 
 
 
 
 
\needspace{5cm} \section{Dabovaná} \label{dabovaná} \katabox{Téma 
}{4 
}{2 minuty} 
 
 
V kategorii \textbf{Dabovaná}{} vám hráči na scéně zahrají na vaše téma, ale pouze pantomimou. Slova jim do úst vkládají dabéři, kteří jsou mimo scénu. 
 
 
\subsection{Průběh} Připraví se sudý počet hráčů, ideálně dva z každého týmu. Jeden se stává dabovaným, který poté hraje a druhý se stává dabérem, který za hrajícího kolegu mluví (oba jsou ze stejného týmu).  
Dabovaní poté na scéně rozehrají improvizaci na zadané téma, ale pouze pantomimou. Jejich promluvy pronášejí dabéři mimo scénu zpravidla před forbínou, aby na své kolegy dobře viděli. Hráči na scéně musí dávat pozor, co bylo řečeno a podle toho také jednat a naopak dabéři musí zohledňovat to, co jejich kolegové na scéně hrají. Pro správné fungování kategorie je důležité, aby dabovaným hráčům bylo vždy dobře vidět na obličej. Pokud se hráč otočí zády, jeho dabér ho nemůže dobře dabovat. 
 
Na konci je vhodné, když se jednotlivé hlasy ještě jednou představí. 
 
\subsection{ Varianty } \begin{itemize}
\item Hraje 6 nebo 8 hráčů (3, 4 dabovaní a 3, 4 dabéři).
\item Dabéři mluví do mikrofonů.
\end{itemize}
 
\subsection{ Běžné chyby } \begin{itemize}
\item Dabovanému hráči není vidět do obličeje.
\item Dabér nepostřehne, že jeho kolega otevírá pusu a nemluví za něj.
\item Jeden dabér stále mluví a nedá tak prostor svému kolegovi.
\item Dabéři mluví přes sebe a není jim rozumět.
\end{itemize}
 
 
 
 
 
\needspace{5cm} \section{Diapozitivy} \label{diapozitivy} \katabox{Téma 
}{6+ 
}{3 minuty} 
 
\odkaz{Hráči}{hráč} vám předvedou prezentaci s diapozitivy na dané téma. 
 
\subsection{ Průběh } Jeden \odkaz{hráč}{hráč} z každého týmu se připraví dopředu na prezentaci, kterou na začátku krátce uvedou. Scéna je rozdělena na přední část, kde stojí moderátoři prezentace a zadní část, kde se objevují diapozitivy. Slovem "Cvak"{} prezentující přepínají obraz. Při přepnutí obrazu naběhne neurčený počet hráčů do zadní části a zastaví se ve \odkaz{štronzu}{štronzo}. Moderátoři poté vysvětlují, co je na obrazu a inspirují se pozicí, ve které jsou diapozitivy. 
 
\subsection{ Technické výrazivo } \begin{itemize}
\item \textbf{cvak}{} - pro přepnutí snímku  - hráči změní pozici, vyskočí ze snímku či do něj vskočí a strnou ve \odkaz{štronzu}{štronzo}
\item \textbf{nepatřící snímek}{}
\item \textbf{dva snímky přes sebe}{} - část hráčů na následující cvaknutí odbíhá, je možné že se snímek s jejch pozicí ještě vrátí.
\item \textbf{smítko/prach/moucha}{} - moderátor odfoukne jednoho z hráčů ze snímku pryč
\item \textbf{stranové převrácení}{} - hráči se přeorganizují dle svislé osy
\item \textbf{převracení vzhůru nohama}{} - hráči se převrátí dle vodorovné osy, obvykle je nutný nějaký stoj na rukách apod.
\item \textbf{zaostření}{} - hráči se přiblíží k sobě i k divákům, zvýrazní mimiku obličeje
\end{itemize}
 
\subsection{ Varianty } \begin{itemize}
\item Scénka začíná s  "úvodním slajdem"{} spuštěným
\item PowerPointová prezentace - v prezentaci se mohou objevit i videa a ozvučené obrázky
\item \odkaz{Kramářská píseň}{kramářská píseň} - "zpívaná verze"
\end{itemize}
 
\subsection{ Běžné chyby } \begin{itemize}
\item Úvod trvá zbytečně dlouho
\item \odkaz{Hráči}{hráč} chodí v obrazu diapozitivů nebo do něj šahají
\item Popletené technické výrazivo (např. hráči \textbf{zoomují}{}, což diaprojektor neumožňuje)
\end{itemize}
 
 
 
 
 
\needspace{5cm} \section{Dopis} \label{dopis} \katabox{Obvykle dvě postavy(kdo a komu píše) 
}{2, následně neomezeno 
}{30 vteřin, 3 minuty 
} 
Dva hráči na střídačku přečtou \textbf{Dopis}{}, dle kterého je následně sehrána celá scéna. 
\subsection{Průběh} Před začátkem kategorie, vyzve \odkaz{MC}{mc} jednoho hráče z každého týmu ať předstoupí a společně přečtou dopis na téma od diváků (či s vyžádaným adresátem a pisatelem). 
Střídají se po cca 1/2 věty.  Většinou na to bývá 20-30 vteřin.  
 
Po přečtení dopisu MC zahájí normální kategorii s časovým omezením, která se může jakkoli inspirovat dopisem (děj se může odehrávat před dopisem, po dopisu, týkat jednoho charakteru z dopisu, týkat se události z dopisu, týkat se zmíněné věci z dopisu atd. ). 
\subsection{Běžné chyby} \begin{itemize}
\item Ignorování dopisu v druhé části
\end{itemize}
 
 
 
 
\needspace{5cm} \section{Duet} \label{duet} \katabox{jednoslovné téma nebo předmět pro každou dvojici 
}{dvojice 
}{volný, do 30 s 
} 
 
V duetech spolu každá dvojice zahraje krátkou scénku na dané téma. Ukončuje se pohledem do diváků. 
 
\subsection{ Průběh } Dvojice stojí na okrajích pódia zády k sobě. \odkaz{MC}{mc} si vyžádá téma/předmět pro každou dvojici. Kreativita není omezena, ale nejčastěji se hraje \odkaz{pantomima}{pantomima} doplněná zvuky. Scénka končí v momentě, kdy se oba podívají nejdříve na sebe a poté do publika. 
Jeden \odkaz{hráč}{hráč} může hrát \odkaz{předmět}{předmět} zadaný od diváků a druhý s ním interaguje, ale možnosti jsou samozřejmě širší. 
 
\subsection{ Varianty } \begin{itemize}
\item \odkaz{Kvartet}{kvartet}
\end{itemize}
 
\subsection{ Běžné chyby } \begin{itemize}
\item Zanedbání očního kontaktu
\end{itemize}
 
 
 
 
\needspace{5cm} \section{Duál} \label{duál} \katabox{Téma}{2 a více 
}{2 min. 
} 
 
V kategorii \textbf{Duál}{} hráči hrají volnou improvizaci ve dvou jazycích, mezi kterými přepínají diváci. 
 
\subsection{Průběh} Hráči rozehrají volnou improvizaci inspirovanou tématem a mluví přitom v češtině. Diváci mají možnost kdykoli během improvizace nahlas vykřiknout slovo "Duál!", což je pro hráče znamení, že musí okamžitě přejít do druhého jazyka - šumlované či zpátky do češtiny. Šumlovaná je smyšlený jazyk, též známý jako švojština či \odkaz{gibberish}{gibberish}. Hráči ani diváci tak přímo nerozumí tomu, co říká právě mluvící hráč. Hráči by se tak měli soustředit také na pantomimický doprovod své mluvy. Diváci mohou vykřikovat dle své libosti, počet přepnutí není omezen.  
 
\subsection{Běžné chyby} \begin{itemize}
\item Divák křičí, ale hráč ho neslyší. \odkaz{MC}{mc} může povel hráčům zopakovat.
\item Diváci nechávají hráče příliš dlouho v jednom jazyce.
\item Opožděná reakce na přepnutí (doříkávání věty)
\item Nerealistické začínání nové věty od začátku po přepnutí
\end{itemize}
 
\subsection{Varianty} \begin{itemize}
\item Přepínání provádí \odkaz{MC}{mc}
\end{itemize}
 
\subsection{Podobné kategorie} \begin{itemize}
\item  \odkaz{Šumlovaná}{šumlovaná}
\item  \odkaz{Šumlovaná se simultánním překladem}{šumlovaná se simultánním překladem}
\end{itemize}
 
 
 
 
\needspace{5cm} \section{Dva na pět} \label{dva na pět} \katabox{prostředí 
}{dva 
}{neomezený} 
 
V kategorii \textbf{Dva na pět}{} dva hráči ztvární alespoň sedm postav. 
 
\subsection{Průběh} \odkaz{MC}{mc} si od diváků vyžádá prostředí, kde se bude improvizace odehrávat. V improvizaci hrají pouze dva hráči, kteří se připraví a poté rozehrají za první dvě postavy. Kdykoli poté je jedním hráčem pojmenována či zavolána nová postava, druhý hráč musí postavu nadále ztvárňovat, ale současně si zachovává i postavy, které hrál předtím. Jeden hráč tak v improvizaci hraje za 1-6 postav, ideální samozřejmě je, když jsou počty postav na hráče vyrovnané. 
 
\subsection{ Běžné chyby } \begin{itemize}
\item  Hráč zapomene na některou ze svých postav.
\item  Hráč neurčí druhému další postavy.
\item  Hráč využívá k posunu děje jen jednu ze svých postav.
\end{itemize}
 
 
 
 
 
 
 
 
\needspace{5cm} \section{Dvojitá stíhačka} \label{dvojitá stíhačka} \katabox{Téma 
}{Dva týmy}{4x40s 
} 
 
Příběh začnou hrát hráči pouze jednoho týmu, jsou pak přesně vystřídáni týmem druhým  
 
\subsection{ Průběh } Průběh je podobný jako v kategorii \odkaz{Stíhačka}{stíhačka}, pouze rozšířen o dvě výměny hráčů. 
 
\subsection{ Rozšíření } Narozdíl od obyčejné \odkaz{stíhačky}{stíhačka}, nemusí \odkaz{hráči}{hráč} prvního týmu hledat problém. Čas navíc by měli využít k \odkaz{expozici}{expozice}. Po první výměně hráči druhého týmu nachází problém. Po druhé výměně \odkaz{hráči}{hráč} prvního týmu problém vyvrcholí a po poslední výměně \odkaz{hráči}{hráč} druhého týmu problém vyřeší. Obvykle téma zazní až po poslední výměně. 
 
\subsection{ Ukazování času } \odkaz{Pomocný rozhodčí}{pomocný rozhodčí} \odkaz{ukazuje}{ukazování času} v prvních třech blocích pouze konec, jen v posledním bloku všechny časové intervaly. 
 
 
 
 
 
\needspace{5cm} \section{Dům a zahrada} \label{dům a zahrada} \katabox{Téma, prostředí 
}{libovolný (ideálně 3 na každé straně) 
}{neomezený} 
 
 
V kategorii \textbf{Dům a Zahrada}{} vám hráči zahrají improvizaci v dvou těsně sousedících prostředích. 
 
 
\subsection{Průběh} Jeviště se ve středu pomysleně rozdělí na dvě poloviny a každé polovině je zadáno prostředí (dům, zahrada). Co je hlavní – prostředí spolu nesousedí středem jeviště, ale jeho portály, kraji. Každý tým hraje jen na své polovině. Na obou stranách hraje stejný počet hráčů, a to tak, že dva soupeřící hráči hrají jednu postavu. Na začátku kategorie nastoupí hráči, kteří hrají libovolně do domu nebo do zahrady, tak, aby byly všechny postavy zastoupené. Rozehraje se improvizace ve které se hráči snaží co nejvíce pracovat s logickým přecházením z domu do zahrady a naopak. Hráči si musí hlídat svého partnera, se kterým hrají stejnou postavu a přebírat jeho akce. 
 
\subsection{ Varianty } \begin{itemize}
\item Vyberou se jiná dvě prostředí, která spolu sousedí (bazén, sauna/restaurace, kuchyně)
\end{itemize}
 
\subsection{ Běžné chyby } \begin{itemize}
\item Hráč zapomene sledovat svého kolegu a zapomene nastoupit nebo nepřevezme akci svého kolegy.
\item Hráči na sebe do sousedních prostředí volají a ohlížejí se přes střed jeviště.
\end{itemize}
\subsection{Podobné kategorie} \begin{itemize}
\item \odkaz{Schizofrenie}{schizofrenie} kde také dva hráči hrají stejnou postavu
\end{itemize}
 
 
 
 
\needspace{5cm} \section{En face} \label{en face} \katabox{Téma 
}{4-6 
}{3 min 
} 
 
Hráči Vám přednesou příběh, budou však zachovávat naprosto strnulé tváře. 
 
\subsection{ Průběh } Hráči se postaví do řady (podobně jako v \odkaz{Spoon river}{spoon river}) střídavě aby hráči stejného týmu nebyli vedle sebe. Během improvizace se hráči nesmějí hýbat a nesmějí mimikou vyjadřovat emoce. Jediným výrazovým prostředkem zůstává slovo a zvuky. Kategorie běžně začíná nejprve zvuky prostředí, kde se bude improvizace odehrávat poté se objeví i postavy. Hráči si po celou dobu zachovávají jednu danou postavu. 
 
\subsection{ Běžné chyby } \begin{itemize}
\item  Hráči se sami ztratí v ději
\item  Hráči mění výraz tváře (smějí se)
\item  Hráč se vůbec nezapojí do příběhu
\item  Hráč nepustí ke slovu ostatní hráče
\end{itemize}
 
\subsection{ Viz také } Má blízko ke kategoriím \odkaz{Spoon river}{spoon river}, \odkaz{Rozhlasová hra}{rozhlasová hra} či \odkaz{Vysílání FM}{vysílání fm}. 
 
 
 
 
\needspace{5cm} \section{Filmové žánry} \label{filmové žánry} \katabox{Téma, několik existujících filmových a divadelních žánrů 
}{2 + 2n 
}{neomezený, 5 - 15 minut} 
 
 
\odkaz{Hráči}{hráč} v této kategorii předvedou krátké scénky podle zvoleného \odkaz{žánru}{:kategorie:žánry}. 
 
 
\subsection{Průběh} Začíná se volně - bez \odkaz{žánru}{:kategorie:žánry}, dvojice spolu začne scénku na dané téma. \odkaz{MC}{mc} přepíná žánry ze seznamu, který má zapamatovaný nebo napsaný. V každém žánru se zahraje krátká scénka a vše co se odehrálo se přelévá do dalšího žánru. Ostatní \odkaz{hráči}{hráč} mohou naskakovat jako \odkaz{rekvizity}{rekvizita} a vedlejší postavy. Vedlejší postavy by měly po změně žánru opustit scénu.  
 
Na \odkaz{zápasech}{zápas} se obvykle vystřídají všichni hráči v dvojicích. \odkaz{MC}{mc} dává povel ke střídání. 
 
\subsection{ Varianty } \begin{itemize}
\item Není žádná výchozí dvojice a postavy volně přicházejí a odcházejí
\end{itemize}
 
\subsection{ Běžné chyby } \begin{itemize}
\item Běžné chyby, které se vyskytují v \odkaz{žánrech}{:kategorie:žánry}
\item Vedlejší postavy upozadí hlavní dvojici
\item MC nerespektuje žánry navržené diváky
\end{itemize}
 
 
 
 
 
\needspace{5cm} \section{Kam zmizel ten starý song} \label{kam zmizel ten starý song} \katabox{Téma, povolání/specializace 
}{libovolný (3 a více) 
}{cca 3 minuty na každou píseň} 
 
 
V kategorii \textbf{Kam zmizel ten starý song}{} moderátor s hostem vzpomínají na písničky, které hosta formovaly. 
 
 
\subsection{Průběh} Moderátor a host spolu krátce hovoří o specializaci/životě hosta. 
Vždy je přivedena řeč na nějakou písničku, kterou moderátor "pustí"{} - nechá hráče, aby ji zazpívali. 
Píseň bývá pojmenována stylem, názvem, případně může být i nějaká citace z ní, která by pak měla také zaznít. 
Zpívajícímu hráči/hráčům ostatní dělají \odkaz{chór}{chór}. 
 
\subsection{ Varianty } \begin{itemize}
\item Zpívají dvojice, každý hráč z jiného týmu.
\item Každý hráč zpívá sám, na zápasech je pak vhodné nechat každý tým shodný počet písní.
\end{itemize}
  
 
\subsection{ Běžné chyby } \begin{itemize}
\item Rozkecávání moderátora s hostem nad únosnou mez.
\item Moderátor s hostem určí obskurní žánry
\end{itemize}
 
 
 
 
 
\needspace{5cm} \section{Komiks} \label{komiks} \katabox{Téma 
}{4 
}{cca 1 minuta} 
 
Hráči vám ukáží a přečtou \textbf{Komiksový}{} strip. 
 
\subsection{Průběh} Dva hráči zaujmou pózy prvního obrázku stripu. 
Druzí dva hráči "přečtou"{} text komiksového okénka - ta postava z obrázku, která má otevřená ústa, mluví, pokud má ústa zavřená, jedná se o myšlenkovou bublinu. 
Následně se totéž opakuje pro druhý a třetí obrázek stripu (či další, pokud \odkaz{MC}{mc} určil jinak. 
Je záhodno používat i komiksové výrazivo typu BOOM apod. 
 
 
 
 
\needspace{5cm} \section{Kvartet} \label{kvartet} \katabox{jednoslovné téma nebo předmět 
}{4 
}{volný 
} 
 
 
Čtyři hráči zahrají krátkou improvizaci na zadané téma a zakončí ji společným pohledem do publika. 
 
 
\subsection{Průběh}  
Podobné jako \odkaz{duet}{duet}, avšak ve čtyřech hráčích. Každý hráč začíná zády ke středu v jednom rohu jeviště. Vše platí jako u duetu - na povel \odkaz{MC}{mc} se všichni hráči najednou otočí a rozehrají společně improvizaci, ve které by měla být více zastoupena pantomima než řeč. Kategorie je pro hráče náročná i co se vnímání ostatních hráčů týče. Doporučuje se více klást důraz na kontrast, ať v tělovém napětí, zvucích, statusu, či v čemkoli jiném. Improvizace končí pohledem všech hráčů na sebe a poté společně do diváků. 
 
\subsection{ Běžné chyby } \begin{itemize}
\item Improvizace se příliš protáhne a nedaří se najít její závěr
\end{itemize}
 
 
 
\needspace{5cm} \section{Loutky} \label{loutky} {{todo|zpřehlednit průběh, youtube}} 
 
\katabox{Téma 
}{sudý počet, 4 a více 
}{3 minuty} 
 
V kategorii \textbf{Loutky}{}  uvidíte krátké loutkové představení. 
 
\subsection{Příprava scény} Na scéně se postaví židličky vedle sebe. Na každé židli stojí hráč - loutkovodič, pod každou židli si zaleze klečící shrbený hráč - loutka. 
 
\subsection{Průběh} Loutkovodič zdánlivě vede loutku, veškeré pohyby ale přitom určuje loutka, kterou loutkovodič stínuje, a dabuje. 
 
\textbf{Loutka}{} po celou dobu drží strnulý výraz tváře (škleb) a pozici prstů a dlaně, může hýbat (předstírat že je vedena) pouze rukama, či v případě klidové pozice rukou podél těla otáčet celým trupem. 
 
\textbf{Loutkovodič}{} drží loutku buď  za "drátek", s ním manipuluje, když chce loutkou otočit na stranu či ji položit zpět. 
Jinak manipuluje pouze "provázky"{} připevněnými k zápěstí loutky.   
  
\subsection{ Běžné chyby } \begin{itemize}
\item  Zamotané provázky - loutkovodiči pohybují s loutkami tak, že by se jim provázky bývaly byly zamotaly.
\end{itemize}
 
 
\subsection{Varianty} Loutku může Loutkovodič v průběhu improvizace "zandat"{} a znovu vytáhnout. Pokud mezitím změní obličej a pozici prstů a dlaně, jde o novou postavu. Pokud ji nezmění, pokračuje příběh na novém místě.  
 
 
 
\needspace{5cm} \section{Love song} \label{love song} \katabox{Dobrovolník z publika, židle 
}{2 
}{neomezený, 3 - 5 minut (podle počtu hráčů)} 
 
 
V kategorii \textbf{Love song}{} je dobrovolníkovi z publika na místě složena a zazpívána zamilovaná píseň. Kategorie je \odkaz{porovnávací}{:kategorie:porovnávací kategorie}. 
 
 
\subsection{Průběh} \odkaz{MC}{mc} vybere jednoho člověka z publika a zeptá se ho na jméno. Hráči se střídavě zeptají na určený počet otázek (3-5). Odpovědi na jejich otázky slouží jako inspirace pro text písně. Poté v doprovodu \odkaz{hudebníka}{hudebník} každý z hráčů zazpívá svoji milostnou/zamilovanou píseň. Dobrovolník následně vybírá, která píseň byla lepší. 
 
\subsection{ Varianty } \begin{itemize}
\item \odkaz{Break-up song}{break-up song}
\end{itemize}
 
\subsection{ Běžné chyby } \begin{itemize}
\item Zapomenutí jména, komu se zpívá
\item Nevyužití informací od člověka
\item Divné či ano/ne otázky
\end{itemize}
 
 
 
 
 
 
 
\needspace{5cm} \section{Metafory} \label{metafory} \katabox{Dva názvy předmětů 
}{libovolný, obvykle se zapojí všichni 
}{volný, dokud nedojdou nápady} 
 
Hráči spojují podobnosti dvou různých předmětů nebo činností - vyvářejí metafory. Zadání může být '''"Děti jsou jako řidičák"'''. 
 
\subsection{Průběh}  
\odkaz{MC}{mc} vybere od publika více zadání metafor. Poté se hráči postaví do půlkruhu a střídají se. Každý hráč, který vyběhne, řekne svojí metaforu. Metafory si drží formu ''{Něco} je jako {něco}, {vysvětlení}''. Obvykle MC po chvíli činnost vymění, zadá jinou. 
 
\subsubsection{Příklad metafory}  
''Děti jsou jako řidičák. Tak dlouho se na ně těšíš a když je máš, tak zjistíš, že nemáš na auto.'' 
 
\subsection{ Varianty } \begin{itemize}
\item  \odkaz{Tisíc způsobů jak}{tisíc způsobů jak}
\item  \odkaz{Věty, které by neměly zaznít}{věty, které by neměly zaznít}
\end{itemize}
 
\subsection{ Běžné chyby } \begin{itemize}
\item  Sklouzávání k záchodovému humoru
\item  Sklouzávání k záporům "Něco není jako něco"
\end{itemize}
 
 
 
 
 
\needspace{5cm} \section{Mikrofon} \label{mikrofon} \katabox{Téma 
}{libovolný 
}{neomezený,volný} 
 
 
V kategorii \textbf{Mikrofon}{} hráči rozehrají improvizaci a na divácký pokyn, vykřiknutí slova "Mikrofon!"{} ten který jako poslední promluvil zazpívá písničku, ve které představí své vnitřní touhy, bolesti a záměry. 
 
 
\subsection{Průběh} Hráči rozehrají volnou improvizaci na dané téma. Diváci mají možnost kdykoli během improvizace zvolat slovo "\textbf{Mikrofon}{}". \odkaz{Hráč}{hráč}, který mluvil naposled nebo hráč, kterého určí \odkaz{MC}{mc} dostane mikrofon a zazpívá píseň, ve které by mělo zaznít něco, co o jeho postavě ještě nevíme - její touhy, bolesti, plány, atd. Po krátké písni se vracíme do děje, před písní a příběh pokračuje.  
 
\subsection{ Varianty } \begin{itemize}
\item Zpívají dva hráči současně.
\item Heslo "Mikrofon!"{} mohou diváci vykřiknout kdykoli během večera a daný hráč, zazpívá píseň bez ohledu na právě hranou kategorii.
\end{itemize}
 
\subsection{ Běžné chyby } \begin{itemize}
\item Hráči je dán vypnutý mikrofon.
\item Hráč zpívá o tom, co o jeho postavě už víme.
\item Diváci vykřiknou mikrofon se začátkem improvizace.
\end{itemize}
 
 
\subsection{ Podobné kategorie } \begin{itemize}
\item \odkaz{Monolog}{monolog}
\end{itemize}
 
 
 
 
 
\needspace{5cm} \section{Monolog} \label{monolog} \katabox{Téma 
}{libovolný 
}{neomezený,volný} 
 
 
V kategorii \textbf{Monolog}{} hráči rozehrají improvizaci a na Váš pokyn, vykřiknutí slova "Monolog!"{} Vám jeden z nich v monologu představí své vnitřní pocity. 
 
 
\subsection{Průběh} Hráči rozehrají volnou improvizaci na dané téma. Diváci mají možnost kdykoli během improvizace zvolat slovo "\textbf{Monolog}{}". \odkaz{Hráč}{hráč}, který mluvil naposled nebo hráč, kterého určí \odkaz{MC}{mc} dostane mikrofon a řekne monolog, ve kterém by mělo zaznít něco, co o jeho postavě ještě nevíme - její touhy, bolesti, plány, atd. Po monologu se vracíme do děje a příběh pokračuje.  
 
\subsection{ Běžné chyby } \begin{itemize}
\item \odkaz{Hráči}{hráč} je dán vypnutý mikrofon.
\item \odkaz{Hráč}{hráč} svým monologem nerozvíjí příběh.
\end{itemize}
 
 
\subsection{ Podobné kategorie } \begin{itemize}
\item \odkaz{Mikrofon}{mikrofon}
\end{itemize}
 
 
 
 
\needspace{5cm} \section{Nátlak (kategorie)} \label{nátlak (kategorie)} \katabox{úkol pro jednoho hráče 
}{2 
}{3 minuty, či do splnění úkolu} 
 
V kategorii \textbf{Nátlak}{} má jeden hráč za úkol dotlačit druhého k provedení úkolu. 
 
\subsection{Průběh} \odkaz{MC}{mc} přidělí jednomu hráčovi činnost, ke které  má tento dohnat druhého, který mezitím neposlouchal. Dál se rozjíždí děj stejně jako ve \odkaz{volné}{volná}, ale první hráč má stále na mysli svůj úkol.  
 
\subsection{Příklad} Má-li hráč donutit druhého ke zpěvu, může např. hrát dítě ukládající se ke spánku nebo hudebníka na rockovém koncertě. 
  
\subsection{ Běžné chyby } \begin{itemize}
\item  "tlačící hráč"{} přímo řekne tomu druhému.
\end{itemize}
 
 
 
 
 
 
 
 
\needspace{5cm} \section{Poetická} \label{poetická} \katabox{Téma 
}{2 
}{3 minuty} 
 
V kategorii \textbf{Poetická}{} smí \odkaz{hráči}{hráč} mluvit pouze veršovaně. 
 
\subsection{ Průběh } V podstatě se jedná o kategorii \odkaz{volná}{volná}. Jediným omezením \odkaz{hráčů}{hráč} je povinnost skládat věty do veršů. Hráči by měli průběhu kategorie zahrát krátký příběh. 
Do kategorie mohou nabíhat ostatní hráči jako \odkaz{postavy}{postava} nebo \odkaz{rekvizity}{rekvizita}. 
 
Krásné je, když hráč nerýmuje jen "sám na sebe"{} ale spolupracuje s druhým hráčem a doříkává za něj verše. 
 
\subsection{ Běžné chyby } \begin{itemize}
\item Nikomu se nechce, je opožděný začátek kategorie a \odkaz{rozhodčí}{rozhodčí} je nucen \odkaz{písknout}{pískání} \odkaz{plynulost hry}{plynulost hry}.
\item \odkaz{Hráči}{hráč} se blokují tím, že se přestanou hýbat
\item Hráči zapomenou veršovat
\end{itemize}
 
\subsection{Varianty} \begin{itemize}
\item \odkaz{Divadlo poezie}{divadlo poezie}
\end{itemize}
 
 
 
 
 
\needspace{5cm} \section{Polovina času} \label{polovina času} \katabox{téma}{2 a více 
}{Např 60, 30, 15, 7.5, 3.75, 1.8   
} 
 
V \textbf{Polovině času}{} hráči zahrají krátký příběh, který následně zahrají rychleji -  za polovinu času , následně ještě rychleji… 
 
\subsection{ Průběh } Dva hráči rozehrávají příběh, který by měl být hodně akční. Ostatní hráči se zapojují hlavně jako rekvizity ale mohou i jako další postavy. Po uplynutí časového limitu se celý příběh přehrává znovu, s tím, že by mělo zaznít a proběhnout vše, co bylo v původní, nejdelší verzi. Toto se opakuje až do nejkratšího časového úseku. Proto je nutné zjednodušovat, zkracovat dialogy, odebírat slovní vatu. Dobré je také spíše jednat, než mluvit, akce se zkrátit dají rozvité souvětí už tolik ne. 
 
\subsection{ Varianty } \begin{itemize}
\item Po předvedení na 5 vteřin se dá hlasovat, zda pokračovat se zkracováním, či scénu sehrát v deseti vteřinách pozpátku.
\item Když začnete na 64 vteřinách, je to mnohonásobně jednodušší počítat.
\item Nejkratší časový úsek se nesnaží do sebe vměstnat obsah celého příběhu, ale jen jeho zakončení.
\end{itemize}
 
\subsection{ Běžné chyby } \begin{itemize}
\item V kratších časech se všichni chtějí prosadit a vzniká nevzhledná kakofonie.
\item Hráči v nejdelším časovém limitu jen stojí a mluví.
\end{itemize}
 
 
 
 
 
\needspace{5cm} \section{Poslední věta} \label{poslední věta} \katabox{Téma a prostředí 
}{dvojice (2 a více) 
}{do nalezení pointy, 5-10 minut} 
 
 
Každá dvojice sehraje scénku ve svém \odkaz{prostředí}{prostředí}, \odkaz{MC}{mc} mezi prostředími/dvojicemi přepíná. 
 
 
\subsection{Průběh} \odkaz{MC}{mc} od diváků vybere prostředí pro každou dvojici. Rozehrává se scénka, dokud nedojde k přepnutí. Scénka v dalším prostředí začíná poslední větou z předchozího. Po přepnutí prostředí by poslední věta měla zaznít od \odkaz{MC}{mc} spolu s prostředím. 
 
\subsection{ Varianty } \begin{itemize}
\item Závěr se hledá pouze v jednom prostředí
\item Všechny dvojice si najdou závěr
\end{itemize}
 
\subsection{ Běžné chyby } \begin{itemize}
\item Zapomenutí první věty
\item \odkaz{MC}{mc} přepíná na úkor příběhu ve scénkách
\end{itemize}
 
 
 
 
 
\needspace{5cm} \section{Počet slov} \label{počet slov} \katabox{Téma,počet pro každého hráče (ten může být přidělený) 
}{libovolný (3 a více) 
}{3 minuty} 
 
 
V kategorii \textbf{Počet slov}{} známé též jako \textbf{7 slov}{} hráči  budou přesně dodržovat počet slov ve větě. 
\subsection{Příprava} \odkaz{MC}{mc} určí každému hráči, kolik slov smí používat - typicky 1,3,5,7.  
 
\subsection{Průběh} Sehrává se scéna, každý hráč musí v každé své větě/replice použít přesně tolik slov (včetně předložek a spojek) , jako mu bylo přiděleno. 
Často se některý divák pověří, aby kontroloval zejména ty delší varianty. 
 
 
\subsection{ Běžné chyby } \begin{itemize}
\item  Přepočítání se
\item  Nekompletní věty
\end{itemize}
 
 
 
\needspace{5cm} \section{Prskavka} \label{prskavka} \katabox{jednoslovné téma pro každého hráče 
}{libovolný, každý sám za sebe 
}{30 vteřin 
} 
V kategorii \textbf{Prskavka}{} má každý hráč sólový výstup na 30 vteřin a co nejlépe ztvárňuje Vaše jednoslovné téma. 
 
\subsection{Průběh} Prskavka se zaměřuje na tvorbu jednotlivce. 
Hráči stojí zády k publiku, \odkaz{MC}{mc} si vyžádá jednoslovné téma (nejčastěji se vybírají předměty), poté vybere hráče, který bude na toto téma hrát a zahájí improvizaci. Hráč má pak právě 30 vteřin (doba hoření idealizované prskavky) na ztvárnění tématu, bez omezení - může zpívat, veršovat, hrát více postav nebo (nejčastěji) hrát pantomimu. 
 
\subsection{ Varianty } \begin{itemize}
\item \odkaz{Prskavka ala šanson}{prskavka ala šanson}
\item Během představeních v okolí vánočních svátků se dá jako \odkaz{časomíra}{ukazování času} použít místo stopek reálná prskavka.
\end{itemize}
 
 
 
 
 
 
 
 
 
\needspace{5cm} \section{Pyramida} \label{pyramida} \katabox{Prostředí, Téma nebo činnost 
}{4-6 
}{volný} 
 
V \textbf{pyramidě}{} se hraje serie příběhů s narůstajícím počtem hráčů, pokaždé v novém prostředí. Následně se začne počet hráčů snižovat a prostředí se vrací. 
 
 
\subsection{ Průběh } Kategorie je podobná cvičení \odkaz{Roztočený brankář}{roztočený brankář}. Začíná jeden hráč (špička pyramidy), který volně improvizuje na zadané \odkaz{askfor}{askfor}. \odkaz{MC}{mc} ho zastaví ve \odkaz{štronzu}{štronzo} písknutím a přichází druhý hráč. Spolu rozehrají novou scénu v jiném prostředí, je dobré, když vycházejí z pozice, ve které jsou. Tímto způsobem se pokračuje dokud není na scéně maximum hráčů (obvykle 6). 
\odkaz{MC}{mc} scénu se šesti hráči po chvíli zastaví, \odkaz{ posune}{posun ve štronzu} např. na dvě fáze   a začíná se bourání pyramidy. Začíná se v poslední rozehrané scénce a hráč, který ji vytvořil z ní musí \odkaz{motivovaně odejít}{motivovaný odchod}. \odkaz{MC}{mc} ohlásí odchod písknutím nebo tlesknutím a pokračuje se ve scénce, která byla rozehraná předtím. Tak to jde dál, až je na scéně pouze první hráč. Pyramidu ukončuje \odkaz{MC}{mc} 
 
\subsection{ Varianty } \begin{itemize}
\item Větší autonomie - \odkaz{hráči}{hráč} si tleskají sami, kategorii ukončují sami
\item Před bouráním pyramidy dojde k časovému nebo pohybovému posunu
\end{itemize}
 
\subsection{ Běžné chyby } \begin{itemize}
\item Nikdo neví, do kterého prostředí se vrací
\item \odkaz{Hráč}{hráč} zapomene, že má odejít
\end{itemize}
 
 
 
 
 
\needspace{5cm} \section{Překlad do znakové řeči} \label{překlad do znakové řeči} \katabox{Téma,odbornost experta 
}{libovolný (3 a více) 
}{3 min} 
 
Diskuze s expertem je \textbf{překládána do znakové řeči}{}. 
 
\subsection{Průběh} Moderátor a expert (často s odborností přidělenou od \odkaz{MC}{mc}) spolu diskutují, obvykle v sedě na židličkách. Třetí hráč  - tlumočník  - převádí celou diskuzi do znakové řeči. 
 
\subsection{Varianty} Na zápase se tlumočník v polovině času vymění za hráče druhého týmu. 
 
\subsection{ Běžné chyby } \begin{itemize}
\item  Expert a moderátor svou gestikulací na sebe strhávají pozornost.
\item  Tlumočník neopakuje gestikulaci u opakovaných výrazů.
\end{itemize}
 
 
 
 
\needspace{5cm} \section{Reklamace} \label{reklamace} \katabox{Co se bude reklamovat 
}{4 
}{neomezený, volný} 
 
 
Uvidíte reklamaci očima servisního technika. 
 
 
\subsection{ Průběh } \odkaz{Hráči}{hráč} se rozdělí na dvojice. Jedna bude reklamovat, druhá bude přijímat reklamaci. Dvojice hráčů, která bude reklamovat odejde ze sálu a nebo je jiným způsobem zbavena možnosti postřehnout, co se bude reklamovat. Následně \odkaz{MC}{mc} vybírá od diváků věc, která se bude reklamovat. Scéna začíná příchodem reklamující dvojice k přijímající dvojici a snaží se vysvětlit, jaké různé problémy má s věcí, kterou přišel reklamovat. Podle reakcí pak reklamující dvojice přichází na to, jaký předmět vlastně reklamuje, čímž \odkaz{kategorie}{kategorie} končí. 
 
\subsection{ Varianty } \begin{itemize}
\item  Hráči nejsou po dvojicích
\item  Do obchodu přijde více lidí se stejným problémem
\end{itemize}
 
\subsection{ Běžné chyby } \begin{itemize}
\item  Reklamující nedávají dost nabídek
\end{itemize}
 
 
 
 
 
\needspace{5cm} \section{Rozhlasová hra} \label{rozhlasová hra} \katabox{Téma 
}{všichni 
}{neomezený} 
 
 
V kategorii \textbf{Rozhlasová hra}{} od hráčů uslyšíte část drama, tak jak byste jej mohli slyšet z rádia. 
 
\subsection{Průběh} Hráči se usadí na jevišti do půlkruhu a v sále se zhasne. Hráči pak podobně jako v kategorii \odkaz{Vysílání FM}{vysílání fm} rozehrají vysílání rádia, ale tentokrát jde jen o jeden pořad a ucelený příběh. Zpravidla začíná Vypravěč rekapitulací toho, co jsme mohli slyšet v minulém díle a poté následuje pokračování, ve kterém vždy jeden hráč ztvárňuje jednu postavu. Tato kategorie se dá přirovnat k \odkaz{En face}{en face} potmě. 
 
\subsection{ Varianty } \begin{itemize}
\item \odkaz{Vysílání FM}{vysílání fm}
\end{itemize}
 
\subsection{ Běžné chyby } \begin{itemize}
\item Nezjistí se dopředu, zda jde v sále úplně zhasnout.
\end{itemize}
 
 
 
 
 
 
\needspace{5cm} \section{Ryba (ve čtyřech)} \label{ryba (ve čtyřech)} \katabox{Prostředí 
}{4 
}{3 minuty 30 vteřin} 
 
 
V kategorii \textbf{Ryba (ve čtyřech)}{} vám 4 hráči zahrají situaci, na jejímž konci budou všichni mrtví. 
 
\subsection{Průběh} Připraví se 4 hráči zády k divákům jako na \odkaz{Prskavku}{prskavka} a je jim zadáno prostředí. Na písknutí \odkaz{MC}{mc} se všichni čtyři otočí, rozehrají situaci a současně musí dávat pozor na časomíru. Funguje zde podobný princip jako v kategorii \odkaz{Smrt v 1 minutě}{smrt v 1 minutě}. V 60 vteřině minutě musí jeden z hráčů motivovaně zemřít, v 120 vteřině umírá druhý, ve 180 vteřině třetí a poslední hráč má dalších 30 vteřin, na jejichž konci musí také zemřít. Mrtvoly v této kategorii zůstávají na scéně do konce improvizace a zbylí hráči s nimi mohou pracovat, ale mrtvoly se už aktivně nezapojují. 
 
\subsection{ Běžné chyby } \begin{itemize}
\item Hráči nesledují čas.
\end{itemize}
 
 
 
 
\needspace{5cm} \section{S rekvizitou} \label{s rekvizitou} \katabox{Reálný předmět , Prostředí 
}{2 
}{2 minuty} 
 
V kategorii \textbf{S rekvizitou}{} vám 2 hráči předvedou, kolika různými způsoby se dá daná rekvizita použít jako předmět. 
 
\subsection{Průběh} Hrají dva hráči, kteří do začátku kategorie rekvizitu nesmí vidět. Pomocný rozhodčí připraví na jeviště \odkaz{rekvizitu}{rekvizita}, kterou může být jakýkoli reálný předmět (deštník, židle, ramínko). Během kategorie se hráči snaží předmět použít co nejvíce způsoby a to tak, že si jej během improvizace předávají a hrají s ním jako by to byl ten předmět, který rekvizita právě představuje (dalekohled, pádlo, mobil, hřeben, atd.).  Předmět by přitom neměl být pojmenován. 
 
\subsection{ Varianty } \begin{itemize}
\item Předmět je vybrán přímo od diváků.
\end{itemize}
 
\subsection{ Běžné chyby } \begin{itemize}
\item Hráči vždy pojmenují, co má rekvizita představovat.
\item Rekvizita se během improvizace promění v malý počet předmětů.
\item Hráči nejsou schopní si rekvizitu technicky předat. Doporučuje se ji například upustit na zem a pokračovat následně v činnosti s imaginárním předmětem.
\end{itemize}
 
 
 
 
\needspace{5cm} \section{Schizofrenie} \label{schizofrenie} \katabox{téma + charakter pro každého hráče}{4}{2 min.} 
 
Ve \textif{Schizofrenie}{} jsou dvě postavy ztvárněné každá dvěma \odkaz{hráči}{hráč}, každá postava má dvě tváře/charaktery. 
 
\subsection{ Průběh } Často hrají jen dvě postavy. Každou postavu hrají dva hráči. Před začátkem improvizace vyzve \odkaz{MC}{mc} hráče, aby předstoupili a diváci jim zadají charakterové vlastnosti či silné \odkaz{emoce}{emoce} Často velice odlišné, ne však naprosto opačné: např. závistivý/zamilovaný, stydlivá/hašteřivá, flegmatik/cholerik, odvážná/bojácná. 
Při kategorii hrají 2 hráči a druzí dva hráči vyčkávají. V momentě kdy chtějí svého kolegu vystřídat tlesknou, hráči na jevišti \odkaz{štronznou}{štronzo}, hráč co tleskl přiběhne ke svému kolegovi, převezme jeho pozici a kategorie pokračuje. Časté proměny jsou funkční a je zábavné, jak se i během věty může charakter postavy zcela proměnit. Většinou platí pravidlo "Tleskejte hodně!". 
 
\subsection{ Časté chyby } \begin{itemize}
\item Charaktery se vybírají příliš dlouho.
\item Vyčkávající hráči málo střídají svého kolegu.
\item Hráč nejedná ve svém zadaném charakteru.
\end{itemize}
 
 
 
 
 
\needspace{5cm} \section{Sci-fi} \label{sci-fi} \katabox{téma  
}{2+libovolně 
}{cca 3 minuty} 
V prostředí sci-fi, vědeckofantastických filmů a knih se odehrává kategorie \textbf{Sci-fi}{} 
 
\subsection{Průběh} Probíhá stejně jako \odkaz{volná}{volná}, ale s využitím prostředí, reálií sci-fi.  
\subsection{Běžné chyby} \begin{itemize}
\item Stereotypní začínání ve stavu beztíže
\end{itemize}
 
 
 
 
\needspace{5cm} \section{Smrt v 1 minutě} \label{smrt v 1 minutě} \katabox{Prostředí 
}{2 
}{60 vteřin} 
 
 
V kategorii \textbf{Smrt v 1 minutě}{} jeden ze dvou hráčů zemře. 
 
\subsection{Průběh} Hráči v daném prostředí improvizují scénku, která vyustí v umrtí jedné postavy přesně po uplynutí jedné minuty. Kategorie se na \odkaz{zápasech}{zápas} často opakuje, protože hráči nesplní její podmínky. 
 
Průběh je podobný jako v kategorii \odkaz{Ryba (ve čtyřech)}{ryba (ve čtyřech)} 
 
\subsection{ Varianty } \begin{itemize}
\item Hraje se volně a k smrti se nějak dojde
\item Hráči generují co největší množství nechtěných zranění svému kolegovi či sobě
\end{itemize}
 
\subsection{ Běžné chyby } \begin{itemize}
\item Hráči nesledují čas
\item Nikdo nezemře v šedesáté vteřině
\item Smrt je příliš náhlá a očekávatelná
\item \odkaz{Blokování}{bránění ve hře} druhé postavy znemožněním pohybu (svázané končetiny, atd.)
\end{itemize}
 
 
 
 
\needspace{5cm} \section{Spoon river} \label{spoon river} \katabox{Dvojslovné téma, prostředí 
}{4-6 
}{Neomezený, bývá 5 a více minut, typicky alespoň na 3 promluvy za každého hráče. 
} 
V kategorii \textbf{Spoon river}{} uslyšíte příběh několika postav, které zdánlivě nemají nic společného. 
 
\subsection{Průběh} Hráči se postaví na kraj pódia (stejně jako v \odkaz{En face}{en face}), v náhodném pořadí se střídají v monolozích z pohledu jednotlivých postav. 
Při prvním monologu se zdá, že každá postava je na ostatních zcela nezávislá, s každou další promluvou se však čím dál tím víc zjišťuje, že tomu tak není, osudy postav se porůznu protínají. 
 
\subsection{Časté chyby} \begin{itemize}
\item Postavy se nepropojí.
\item Postavy se propojí příliš jednoznačně.
\end{itemize}
 
\subsection{Název} Odkazuje na knihu Spoon River Anthology, pojednávající historii jednoho městečka formou epitafu. 
 
 
 
 
\needspace{5cm} \section{Sportovní komentátor} \label{sportovní komentátor} \katabox{Vrcholový sport 
}{3 
}{3 minuty} 
 
V kategorii \textbf{Sportovní komentátor}{} uvidíte návštěvu vrcholového sportovce v televizním studiu s překladem do znakové řeči. 
 
\subsection{ Průběh } Hráči se rozdělí na moderátora, sportovce a překladatele. Moderátor uvítá a zběžně představí hosta, kterým je vrcholový sportovec sportu, který dali diváci. Celý jejich rozhovor je překládán do \odkaz{znakové řeči}{kvaziznaková řeč} překladatelem. Moderátor a host by měli být při rozhovoru co nejnehybnější a pouze mluvit, aby nestrhávali pozornost na sebe a diváci byli zaměřeni na překladatele. 
 
\subsection{ Varianty } \begin{itemize}
\item  \odkaz{Překlad do znakové řeči}{překlad do znakové řeči}
\item  Střídání překladatelů v polovině (nejčastěji na zápase)
\item  \odkaz{Televizní komentátor}{televizní komentátor}
\end{itemize}
 
\subsection{ Běžné chyby } \begin{itemize}
\item  Moderátor nebo host dělají výrazná gesta strhávají pozornost na sebe
\item  Moderátor nebo host mluví příliš rychle a překladatel nestíhá (pokud to není záměr)
\item  Překladatel se zraní během překladu
\end{itemize}
 
 
 
 
 
 
\needspace{5cm} \section{Stíhačka} \label{stíhačka} \katabox{Téma 
}{2, plus naskakující 
}{1:20, 40s na tým} 
 
V kategorii \textbf{Stíhačka}{} pro Vás hráči jednoho týmu vytvoří příběh a hráči druhého týmu určí, jak dopadne. 
 
 
\subsection{Průběh} \odkaz{MC}{mc} musí \odkaz{náhodně rozhodnout}{náhodné rozhodování}, který tým bude začínat. \odkaz{Hráči}{hráč} z prvního týmu začínají, rozehrají příběh na dané téma od diváků. Po 45 vteřinách se pískne a hráči druhého týmu \odkaz{přeberou pozici}{přebírání pozice} i \odkaz{postavu}{postava} a dohrávají příběh. Obecně funguje velice dobře, když se první dvojice drží vytvoření prostředí, definice postav a nastavení nějakého problému. Druhá dvojice pak už pouze rozvíjí. 
 
Na \odkaz{zápasech}{zápas} se obyčejně hraje \odkaz{dvojitá stíhačka}{dvojitá stíhačka}. 
 
\subsection{ Varianty } \begin{itemize}
\item \odkaz{Dvojitá stíhačka}{dvojitá stíhačka}
\item \odkaz{Trojitá stíhačka}{trojitá stíhačka}
\end{itemize}
 
\subsection{ Běžné chyby } \begin{itemize}
\item První dvojice nezvládne \odkaz{improvizační trojzubec}{improvizační trojzubec}. Druhá dvojice musí doplnit
\item Druhá dvojice nedoplní
\item Nepřejme pozici nebo postavu
\item Hráči špatně rozdělí úlohy ve výměnách
\end{itemize}
 
 
 
 
\needspace{5cm} \section{Tam a zpět} \label{tam a zpět} \katabox{Prostředí, slovo pro každého hráče 
}{libovolný (4 a více) 
}{neomezený,volný , 3 -5 minut (podle počtu hráčů)} 
 
 
V kategorii \textbf{Tam a zpět}{} hráči sehrají příběh, na scénu musí přijít či odejít pouze ve chvíli, kdy padne jejich klíčové slovo. 
 
 
\subsection{Průběh} \odkaz{MC}{mc} každému hráči určí jedno slovo, vetšinou z daného prostředí, a také který hráč má začínat na scéně. 
Hráč začíná v daném prostředí, a ve chvíli, kdy uvede klíčové slovo jiného hráče, ten musí bez meškání naběhnout na scénu a zapojit se do děje až do chvíle, kdy znovu padne jeho klíčové slovo, pak \odkaz{motivovaně odchází}{motivovaný odchod}. Padne li jeho klíčové slovo vícekrát za sebou, musí vícekrát přijít/ odejít. 
 
Na závěr poslední hráč odvolá sám sebe, či, lépe, předposlední při svém odchodu logicky odvolá posledního. 
 
\subsection{ Běžné chyby } \begin{itemize}
\item Hráči si nepamatují klíčová slova.
\item Hráči nereagují na svá klíčová slova.
\end{itemize}
 
 
 
 
 
\needspace{5cm} \section{Televizní komentátor} \label{televizní komentátor} \katabox{Téma, pro každé studio 
}{2x3 
}{3 minuty} 
 
V kategorii \textbf{Televizní komentátor}{} uvidíte náhled dvou televizních studií na různá témata. 
 
\subsection{ Průběh } Průběh je podobný jako v kategorii \odkaz{Sportovní komentátor}{sportovní komentátor}. Hráči se rozdělí na moderátora, hosta a překladatele do \odkaz{znakové řeči}{kvaziznaková řeč}, každý tým má vlastní studio. \odkaz{MC}{mc} přidělí témata studiím a vybere \odkaz{náhodným rozhodováním}{náhodné rozhodování}, kdo bude začínat. V průběhu kategorie \odkaz{MC}{mc} přepíná televizní studia. 
Studio, na které není fokus, zůstává ve \odkaz{štronzu}{štronzo}, ale čas v něm zatím běží, děj se přesouvá dál. 
 
\subsection{ Varianty } \begin{itemize}
\item  \odkaz{Sportovní komentátor}{sportovní komentátor}
\end{itemize}
 
 
 
 
\needspace{5cm} \section{Toaster} \label{toaster} \katabox{Prostředí 
}{4-7}{cca 6 minut  
} 
 
Hráči vyskakují jako topinky z topinkovače a pokračují v epickém příběhu ze zadaného prostředí. 
 
\subsection{ Průběh } Hráči se rozmístí po scéně, čekají na bobku, ve chvíli, kdy \odkaz{MC}{mc} pískne, náhodný počet hráčů se rychle postaví, ti všichni jsou v tu chvíli ve scéně a je nutné s nimi hrát.  Na každé další písknutí jdou stojící hráči k zemi, případně se můžou znovu ihned postavit. Po celou dobu si drží svoji postavu   a rozvíjí celkový příběh. 
 
 
\subsection{ Běžné chyby } \begin{itemize}
\item Uklízení, hledání něčeho.
\item Vyskakování vždycky na každou druhou dobu.
\end{itemize}
 
 
 
\needspace{5cm} \section{Vnitřní hlasy} \label{vnitřní hlasy} \katabox{Téma 
}{4 
}{2-3 minuty} 
 
 
V kategorii \textbf{Vnitřní hlasy}{} vám 2 hráči zahrají situaci na Vaše téma a další dva hráči vám řeknou, co si ti dva při tom myslí. 
 
 
\subsection{Průběh} Hrají 4 hráči, 2 na scéně a 2 na straně, tak že nejsou vidět, kteří dostanou mikrofony.  Dva hráči tvoří v podstatě jednu postavu. Hráči na scéně rozehrají volně situaci a druzí dva jsou jejich Vnitřní hlasy - doplňují akci a promluvu svého kolegy nějakou myšlenkou, kterou si postava myslí a bývá zpravidla protichůdného rázu. Např. situace První rande. Hráč na scéně: Jé, tobě to dneska ale sluší. Vnitřní hlas: Ale v těch šatech jsi docela tlustá. 
 
\subsection{ Varianty } \begin{itemize}
\item Muže se objevit třetí postava na scéně v nějaké malé roli, která nemá vnitřní hlas.
\end{itemize}
 
\subsection{ Běžné chyby } \begin{itemize}
\item Hráči na scéně neustále mluví a nedávají prostor svým Vnitřním hlasům.
\item Hráči dostanou vypnutý mikrofon.
\end{itemize}
 
 
 
 
\needspace{5cm} \section{Volná} \label{volná} \katabox{téma 
}{2+libovolně 
}{cca 10 minut} 
V kategorii \textbf{volná}{}  hráči odehrají krátký příběh. 
 
\subsection{Průběh} V kategorii nejsou žádná omezení, začínají dva hráči a odehraje se nějaký příběh. 
 
\subsection{Varianty} \begin{itemize}
\item \odkaz{Sci-fi}{sci-fi}, \odkaz{Western}{western}, \odkaz{Červená knihovna}{červená knihovna} - je přidáno omezení využívání daného žánru.
\end{itemize}
 
 
 
\needspace{5cm} \section{Vyprávěná} \label{vyprávěná} \katabox{téma 
}{neomezen 
}{2 minuty} 
 
V kategorii \textbf{Vyprávěná}{} se jeden z hráčů stává vypravěčem a ostatní hráči hrají volně na podněty od vypravěče. 
 
\subsection{ Průběh } První hráč, který do kategorie naběhne, je vyprávěč. Vypravěč nás uvede do příběhu, prostředí a představí nám postavy, které vstupují na scénu. Ostatní hráči reagují na to, co řekl vypravěč, ale současně mají možnost sami děj vytvářet. Vypravěč má také možnost posouvat děj v čase a diváci mohou vidět, co bylo příčinou současného jednání té které postavy, když se jí před dvaceti lety něco stalo a poté se vrátíme zpět do současnosti. Děj je také možné posouvat v čase dopředu, zvlášť v případech, kdy se nic neděje, případně když hráč navrhne nějakou zdlouhavou aktivitu. 
 
\subsection{ Běžné chyby } \begin{itemize}
\item Vypravěč stále mluví a nedává tak prostor ostatním hráčům.
\item Vypravěč popisuje to, co hráči dělají = neposouvá děj.
\item Hráč hraje jen to, co řekl vypravěč = neposouvá děj.
\end{itemize}
 
 
 
 
 
\needspace{5cm} \section{Vysílání FM} \label{vysílání fm} \katabox{Téma 
}{5 a více, obvykle se zapojí všichni 
}{volný, obvykle 5 - 10 minut} 
 
 
Diváci  si připomenou, jak mohl vypadat poslech rádia. 
 
 
\subsection{ Průběh } \odkaz{Hráči}{hráč} se usadí do půlkruhu a v sále se \textbf{zhasne}{}. Začíná se zvuky éteru, než se na sebe všichni naladí. Ve studiu sedí moderátor stanice a někdy i host. Každý hráč se zapojuje volně, například vytvořením znělky rádia, zavoláním do studia nebo jen děláním zvuků. 
 
Po ukončení kategorie se hráči akusticky představí nějakou vlastní replikou z vysílání, aby diváci mohli \odkaz{hlasovat}{hlasování}. 
 
\subsection{ Varianty } \begin{itemize}
\item Do studia mohou volat i diváci
\end{itemize}
 
\subsection{ Běžné chyby } \begin{itemize}
\item Po vyhlášení kategorie se zjistí, že v sále nejde udělat tma
\item Nic se neděje a moderátor plká
\item Hudebník se nezapojí
\end{itemize}
 
 
 
 
 
 
\needspace{5cm} \section{Věty z kapsy} \label{věty z kapsy} \katabox{Téma, prostředí 
}{2 
}{neomezený,volný} 
 
 
V kategorii \textbf{Věty z kapsy}{} vám hráči zahrají na Vaše téma a přitom použijí věty, které jste jim napsali na papírky. 
 
 
\subsection{Průběh} Před zápasem mají diváci možnost napsat krom dvouslovného téma také libovolnou větu na jiný papírek. Témata a věty se vybírají zvlášť. Před kategorií si dva hráči přibližně rovnoměrně rozdělí lístečky s větami a ty si připraví do kapes. Rozehrají improvizaci na dané téma a během ní čas od času vytáhnou lísteček s větou, kterou musí okamžitě a co nejobratněji začlenit do improvizace.  
 
 
\subsection{ Varianty } \begin{itemize}
\item Hráči tahají věty z kapsy podle svého uvážení.
\item Hráči tahají věty z kapsy na povel \odkaz{MC}{mc} či na povel od diváků "Věta!"
\end{itemize}
 
 
\subsection{ Běžné chyby } \begin{itemize}
\item Věty jsou vybrány do košíčku na témata.
\item Divák napíše větu a téma na stejný lístek.
\item Hráči vytažení každého lístečku uvádí bezpečnostní větou - např. Maminka mi vždy říkávala ...
\item Hráč otálí před použitím věty po jejím vytažení.
\end{itemize}
 
 
 
 
 
\needspace{5cm} \section{Western} \label{western} \katabox{téma  
}{2+libovolně 
}{cca 3 minuty} 
 
V prostředí divokého západu, Ameriky konce předminulého století, se odehrává kategorie \textbf{Western}{} 
 
\subsection{Průběh} Probíhá stejně jako \odkaz{volná}{volná}, ale s využitím prostředí, reálií divokého západu.  
 
\subsection{Běžné chyby} \begin{itemize}
\item Stereotypní začínání v saloonu
\end{itemize}
 
 
 
 
\needspace{5cm} \section{Změna} \label{změna} \katabox{Téma 
}{2 
}{2 min} 
 
 
V kategorii \textbf{Změna}{}, občas nazývané \textbf{Jinak}{} hrají hráči volnou improvizaci na téma, ale často se v ní mění \odkaz{repliky}{replika} a pohyby na povel \odkaz{MC}{mc}. 
 
 
\subsection{ Průběh } Kdykoli v průběhu improvizace může od \odkaz{MC}{mc} zaznít povel "změna"{} nebo "jinak". \odkaz{Hráč}{hráč}, který mluvil poslední musí změnit svoji repliku včetně pohybu, kterým ji doprovodil. Pokud není \odkaz{MC}{mc} ani s druhou variantou spokojen, může povel opakovat, ale nemělo by to být na úkor improvizace. 
 
\subsection{ Varianty } \begin{itemize}
\item  \odkaz{MC}{mc} má k dispozici dva povely: \textbf{Jinak}{} = změna poslední \odkaz{repliky}{replika}, \textbf{Jiný pohyb}{} = změna posledního pohybu
\item  Může se zapojit více hráčů
\end{itemize}
 
\subsection{ Běžné chyby } \begin{itemize}
\item Hráči se zablokují po změně a přestanou se hýbat
\item Hráči mění příliš málo.  (Rád bych koupil máslo. Změna. Rád bych koupil margarín. Změna. Rád bych koupil chleba.)
\end{itemize}
 
 
 
 
 
\needspace{5cm} \section{Zpívaná} \label{zpívaná} \katabox{Téma, prostředí 
}{libovolný (3 a více) 
}{neomezený,volný , 3 -5 minut (podle počtu hráčů)} 
 
 
V kategorii \textbf{Zpívaná}{} vám hráči dohromady zazpívají písničku na Vaše téma. 
 
 
\subsection{Průběh} Hráči stojí v těsném půlkruhu a zpívají. Začínají refrénem (2x opakovaný), následně jednotlivci zazpívají sloku, vždy následovanou refrénem, na konci se refrén zpívá vícekrát a vygradovává. Když jednotlivec zpívá sloku (sloky), vystoupí z půlkruhu a zpívá na mikrofon.  
 
S ohledem na plynulost je nutné, aby již na závěr refrénu byl u mikrofonu připraven hráč, který bude zpívat jako další. 
 
 
\subsection{ Varianty } \begin{itemize}
\item  Každý přistoupí k tématu z trošku jiného úhlu pohledu.
\item  Příběh navazuje.
\item  "Stará Zpívaná"{} aká muzikál. Hráči jsou implicitně v půlkruhu jako chór a před nimi se normálně odehrává příběh, však zpíván. Hráči mohou libovolně vstupovat a vystupovat jak se jim zamane. Refrén může být i nemusí. Podobné muzikálu při žánrech.
\end{itemize}
 
 
\subsection{ Běžné chyby } \begin{itemize}
\item  Nikomu se nechce dopředu, vzniká mrtvá pauza vyplněná čekáním.
\end{itemize}
 
 
 
 
 
\needspace{5cm} \section{Červená knihovna} \label{červená knihovna} \katabox{Téma 
}{libovolný 
}{2-3 minuty} 
 
 
V kategorii \textbf{Červená knihovna}{} vám hráči zahrají příběh plný klišé z romantické literatury. 
 
 
\subsection{Průběh} Začínají dva hráči, kteří rozehrávají příběh plný klišé, která známe z televizních novel a volně se k nim přidávají další postavy. Co by se mělo objevit: tajná láska, utajené rodinné vztahy, přetrvávající rodové rozepře, souboj, exotická jména, atd. 
 
\subsection{ Varianty } \begin{itemize}
\item V prvním poločasu je rozehrán začátek příběhu a v druhém poločasu jeho dokončení.
\end{itemize}
 
 
\subsection{ Běžné chyby } \begin{itemize}
\item Zapomenutí jména postavy
\end{itemize}
 
 
 
 
 
\needspace{5cm} \section{Čtverec} \label{čtverec} \katabox{4 * prostředí, a 1 téma navíc 
}{4 
}{neomezen} 
 
Ve \textbf{Čtverci}{} dva hráči na každé straně čtverce dostanou jedno prostředí, na které pak hrají, když je jejich strana natočená k divákům. 
 
\subsection{Průběh} Připraví se čtyři hráči a na scéně zaujmou postavení tvaru čtverce - dva dopředu, dva dozadu. \odkaz{MC}{mc} určí každé dvojici prostředí, které vybere od diváků a hráči si je ještě pro jistotu unisono zopakují. Každý hráč tak hraje ve dvou prostředích. Během improvizace \odkaz{MC}{mc} může čtvercem libovolně otáčet po směru či proti směru hodinových ručiček o jedno až tři prostředí. Běžně se při prvním hraní v daném prostředí (strany čtverce) odehraje úvod, v druhém zápletka a v třetím závěr (takže pak má kategorie 4*3 =12 scén), pokud MC neukončí dřív. 
Při "otočení"{} se přesně přejímají pozice hráčů. 
 
\subsection{ Varianty} \begin{itemize}
\item Technicky podobná je kategorie \odkaz{Poslední věta}{poslední věta} a dá se (zejména na \odkaz{improshow}{improshow}) s touto kategorií celkem bezproblémově kombinovat.
\item Hráči přejímají pozice.
\item Hráči mají shodnou postavu v obou prostředích, a tvoří tak větší příběh.
\end{itemize}
 
\subsection{Běžné chyby} \begin{itemize}
\item Shodné charaktery
\item Hráči se ztratí v otáčení čtvercem
\item Hráč zapomene v jakém prostředí se právě nachází
\end{itemize}
 
 
 
 
 
\needspace{5cm} \section{Šumlovaná} \label{šumlovaná} \katabox{Téma}{2 a více 
}{2 min 
} 
 
V kategorii \textbf{Šumlovaná}{} hráči improvizují v jazyce \odkaz{gibberish}{gibberish}. 
 
\subsection{ Průběh } \odkaz{MC}{mc} se zeptá diváků na nějaký jazyk, může být vymyšlený nebo reálný, ale jiný než kterým hráči mluví. Hráči poté rozehrají volnou improvizaci v \odkaz{gibberish}{gibberish} tak, aby připomínala zvolený jazyk. Význam se do vymyšleného jazyka dodává pohybem. 
 
\subsection{ Běžné chyby } \begin{itemize}
\item  Hráči hrají více hlasem než pohybem a z improvizace se vytrácí smysl
\end{itemize}
 
\subsection{ Podobné kategorie } \begin{itemize}
\item  \odkaz{Šumlovaná se simultánním překladem}{šumlovaná se simultánním překladem}
\item  \odkaz{Duál}{duál}
\end{itemize}
 
 
 
 
\needspace{5cm} \section{Šumlovaná se simultánním překladem} \label{šumlovaná se simultánním překladem} \katabox{Téma}{2 a více 
}{2 min 
} 
 
V této kategorii hráči improvizují v jazyce \odkaz{gibberish}{gibberish} a mají k dispozici překladatele. 
 
\subsection{ Průběh } Na začátku je určen překladatel. Průběh je podobný jako v kategorii \odkaz{Šumlovaná}{šumlovaná}. Hráči rozehrají volnou improvizaci v \odkaz{gibberish}{gibberish} a při hraní dávají prostor překladateli (podobně jako vypravěči v kategorii \odkaz{Vyprávěná}{vyprávěná}), který neutrálním hlasem překládá jejich řeč.  
 
\subsubsection{ Varianty } V zápasové formě se v polovině času překladatelé vystřídají. 
 
\subsection{ Podobné kategorie } \begin{itemize}
\item  \odkaz{Šumlovaná}{šumlovaná}
\item  \odkaz{Duál}{duál}
\end{itemize}
 
 
 
 
\needspace{5cm} \section{Ženy vs muži (kategorie)} \label{ženy vs muži (kategorie)} \todo{Video} 
 
\katabox{dvouslovné téma/situace 
}{skupina mužů, skupina žen 
}{2*2minuty} 
 
 
V kategorii \textbf{Ženy vs Muži}{} uvidíte dva podobné příběhy, jednou z mužského a podruhé z ženského pohledu. 
 
 
\subsection{Průběh} Hraje se jako \odkaz{Volná}{volná}, nejprve hrají hráči jen jednoho pohlaví, a následně na stejné zadání hráči druhého pohlaví... 
 
\subsection{ Varianty } \begin{itemize}
\item  Muži hrají ženy a ženy hrají muže
\item  Pohlaví, které hraje jako druhé, jde na první část za dveře.
\item  Druhá skupina hraje takřka doslova stejný příběh jako první skupina, jen výrazně pozmění některé gendrové detaily.
\end{itemize}
 
 
 
 
\needspace{5cm} \section{Živé rekvizity} \label{živé rekvizity} \katabox{Prostředí, dva dobrovolníci 
}{2 
}{2 minuty} 
 
V kategorii \textbf{Živé rekvizity}{} vám 2 hráči předvedou, kolika různými způsoby se dá využít nevinný divák jako rekvizita. 
 
\subsection{ Průběh } \odkaz{MC}{mc} požádá o dva dobrovolníky z publika. Mezitím se připraví dva hráči, kteří budou hrát. \odkaz{MC}{mc} přidělí každému hráči jednoho dobrovolníka, který pak musí být svému hráči plně k dispozici a nechat  sebou jakkoli manipulovat neboť se stává všemi rekvizitami, které hráč v improvizaci použije. Hráči pak rozehrávají improvizaci na zadané téma a přitom hojně využívají rekvizit, které si modelují z dobrovolníků. Divák se tak může stát dalekohledem, batehem, dveřmi, stolem, šálkem, činkou či jinými předměty. 
 
\subsection{Viz také } \begin{itemize}
\item \odkaz{S rekvizitou}{s rekvizitou}
\item \odkaz{Zvuky}{zvuky}
\end{itemize}
 
 
\subsection{ Běžné chyby } \begin{itemize}
\item Dlouhé hledání dobrovolníků
\end{itemize}
 
 
 
 
 
 
\needspace{5cm} \section{Žánrová vyprávěná} \label{žánrová vyprávěná} \katabox{Téma, 4 žánry 
}{4 
}{3 minuty} 
 
V kategorii \textbf{Žánrová vyprávěná}{}, známé i jako \textbf{Vyprávěná napřeskáčku}{}, či \textbf{Společné vyprávění}{}  vám 4 hráči vám poví společně jeden příběh v různých žánrech. 
 
\subsection{ Průběh } Připraví se 4 hráči, stoupnou si do středu jeviště vedle sebe s pohledem upřeným do diváků. \odkaz{MC}{mc} vybere pro každého z \odkaz{hráčů}{hráč} žánr, ve kterém bude vyprávět. Téma je pro všechny stejné. \odkaz{MC}{mc} podle vlastního uvážení přepíná (ukazuje), vypravěče a tím i žánr příběhu. 
 
\subsection{ Varianty } \begin{itemize}
\item  Neurčují se žánry, ale emoce. Hráči, kteří zrovna nevypráví dělají křoví v emoci, ve které je příběh.
\item  Nepřepíná MC ale divák, který je usazen před hráče zády k ostatním divákům.
\item  Hráčům nejsou dány žánry ani emoce, ale jsou jim zadávány instrukce které musí plnit - nepoužívat spojku a, mluvit ve slovech neobsahující určité písmeno. Pokud chybují vypadávají až do posledního hráče.
\end{itemize}
 
\subsection{ Běžné chyby } \begin{itemize}
\item  Nepozornost a ztracení se v příběhu
\item  Prodleva při přepínání/dlouhé pasáže v jednom žánru
\end{itemize}
 
 
 
 
 
 


\chapter{Kategorie na improshow}\label{další kategorie}
\label{:kategorie:kategorie na improshow}
Tyto kategorie bývají hrány na improshow.
%\documentclass[main.tex]{subfiles}\begin{document}
\needspace{5cm} \section{Abeceda} \label{abeceda} \katabox{Téma 
}{2 
}{2 min} 
 
V kategorii \textbf{Abeceda}{} vám 2 hráči zahrají na Vaše téma, ale tak, že jejich repliky musí začínat písmeny tak, jak jsou seřazeny v abecedě. 
 
\subsection{Průběh} Připraví se dva hráči a je jim zadáno téma či prostředí. Hráči poté rozehrají improvizaci, ale s omezením, že jejich repliky musí začínat písmeny podle pořadí v abecedě. V případě, že se hráč splete může být nahrazen jiným hráčem ze svého týmu. Pokud hráči dojdou až na konec abecedy, začíná se opět od písmene A. 
 
\subsection{ Varianty } \begin{itemize}
\item  Hraje se i na písmena s háčky a čárkami.
\item  Mohou se zapojit další hráči.
\item  Kategorie končí písmenem Z.
\end{itemize}
 
\subsection{ Běžné chyby } \begin{itemize}
\item Hráč se ztratí v abecedě.
\end{itemize}
 
 
 
 
 
\needspace{5cm} \section{Báseň (cvičení)} \label{báseň (cvičení)} Na zadané téma diváků, které může být buď inspirací nebo přímo názvem \odkaz{básně}{báseň}, hráči postupně chodí dopředu a každý přednáší  
jeden verš. Verše se rýmují \odkaz{ABAB}{abab}. 
 
V případě více slok, hráči, kteří složili 1. sloku opouští hrací prostor a vrací se k ostatním hráčům a kdokoliv z nich může začít tvořit 
další sloku. 
 
\subsection{Varianty} \begin{itemize}
\item  báseň končí po po první sloce.
\end{itemize}
  
\subsection{Viz také} \begin{itemize}
\item \odkaz{Poetická}{poetická}
\item \odkaz{Divadlo poezie}{divadlo poezie}
\item \odkaz{Báseň od publika}{báseň od publika}
\end{itemize}
 
\subsection{Běžné chyby} \begin{itemize}
\item Důležitým prvkem u básně je zvýraznit poslední řádek, dát mu důraz, aby posluchač cítil, že touto větou báseň končí.
\item Tím, že verše vznikají na jevišti a nejistotou hráčů, mohou básně působit občas plitce. Je tedy třeba dbát v tomto cvičení na přednes a sebevědomí hráčů.
\end{itemize}
 
 
 
 
\needspace{5cm} \section{Break-up song} \label{break-up song} \katabox{Dobrovolník z publika, židle 
}{2 
}{neomezený, 3 - 5 minut (podle počtu hráčů)} 
 
Hráči zazpívají o rozchodu na motivy informací od jednoho dobrovolníka. 
 
\subsection{ Průběh } Průběh je téměř stejný jako u kategorie \odkaz{Love song}{love song}. \odkaz{MC}{mc} ale z publika požádá diváka, který v nedávné době prošel rozchodem. Píseň je mířená na jeho protějšek. 
 
\subsection{ Varianty } \begin{itemize}
\item \odkaz{Love song}{love song}
\end{itemize}
 
 
 
 
 
\needspace{5cm} \section{Divadlo poezie} \label{divadlo poezie} \katabox{Téma 
}{4 a více 
}{4 minuty} 
 
 
\textbf{Divadlo poezie}{} poetickým způsobem ztvární vaše téma. 
 
 
\subsection{Průběh} Hráči společně tvoří pomocí zvuků, rýmu a pohybu zvukomalebný příběh. 
Využívají přitom opakování slov, zvuků, pohybů, celá skupina se organicky přelévá a využívá synonyma, asociace a přirovnání. 
 
{{todo|rozšířit text, přidat videa.}} 
 
 
 
 
 
\needspace{5cm} \section{Dopředu, dozadu} \label{dopředu, dozadu} \katabox{Téma 
}{2 a naskakující 
}{neomezený} 
 
V této kategorii se podíváme na příběh několikrát a zajímavé části si přetočíme dopředu nebo zpátky. 
 
\subsection{ Průběh } Hraje se volně příběh. \odkaz{MC}{mc} může kdykoliv přepnout směr děje příběhu dopředu nebo dozadu. Pokud \odkaz{hráči}{hráč} hrají příběh směrem dopředu a zazní povel "Dozadu!", opakují pozpátku všechny pohyby a repliky v opačném pořadí než zazněly. Řeč se při směru dozadu neobrací, mluví se normálně. Příběh se odehrává až na začátek nebo dokud \odkaz{MC}{mc} opět nepřepne směr. Poté, co \odkaz{MC}{mc} přepne opět dopředu, všechno se zopakuje, dokud se \odkaz{hráči}{hráč} nedostanou k bodu, kde ještě příběh nebyl zahrán a pokračují v ději. 
 
\subsection{ Běžné chyby } \begin{itemize}
\item  Příliš mnoho povídání, statická scénka
\item  Hráči zapomenou, co se v příběhu stalo
\item  \odkaz{MC}{mc} pouští příliš dlouhé úseky dozadu
\end{itemize}
 
 
 
 
 
\needspace{5cm} \section{Dva začátky} \label{dva začátky} \katabox{Téma 
}{libovolný 
}{5 minut} 
 
 
Diváci si ze \textbf{dvou začátků}{} příběhu vyberou, který chtějí slyšet. 
 
 
\subsection{Průběh} Dva hráči krátce shrnou příběh, přednesou anotaci či ukázku - dva rozličné příběhy na stejné téma. 
Následně se \odkaz{hlasováním}{hlasování} určí, který příběh bude pokračovat, ten následně hráči sehrají. 
 
\subsection{ Varianty } \begin{itemize}
\item  \odkaz{Režiséři}{režiséři}
\end{itemize}
 
 
 
\needspace{5cm} \section{Dvě repliky} \label{dvě repliky} \katabox{Dvě věty pro kažého (až na jednoho) hráče 
}{libovolný (3 a více) 
}{3 minuty} 
 
Hráči mají jen \textbf{dvě repliky}{} na to, aby vyjádřili vše, co chtějí říci. 
 
\subsection{Průběh} Jeden hráč není ve svém vyjadřování nijak omezován, další hráči mají přidělen každý dvě jiné obecné věty (např. Nesnáším tě.  Tohle by se za Marie Terezie nestalo.). Jenom \odkaz{replikami}{replika} mohou reagovat na jakékoli promluvy centrálního hráče, který své nabídky volí tak, aby ostatním tyto reakce umožnil. 
 
\subsection{Běžné chyby} \begin{itemize}
\item  Neomezovaný hráč nenahrává těm ostatním.
\item  Příliš obskurní věty pro hráče.
\end{itemize}
 
 
 
\needspace{5cm} \section{Facebookový status} \label{facebookový status} \katabox{Poslední facebookový status některého diváka 
}{libovolný  
}{3  minuty} 
 
Hráči sehrají příběh ze života inspirovaný  \textbf{Facebookovým statusem}{} některého diváka. 
 
\subsection{Průběh} Hraje se příběh statusem inspirovaný, situace která končí tímto statusem.  
 
 
 
\needspace{5cm} \section{Handicapovaná pohádka} \label{handicapovaná pohádka} \katabox{Známá pohádka a handicap pro každého hráče 
}{5 
}{podle vybrané pohádky} 
 
 
Modifikace známé pohádky od diváků pomocí handicapů. 
 
\subsection{Průběh}  
Hraje obvykle 5 hráčů, každý z nich má jeden handicap (verbální, fyzický, mentální,...) a společně hrají známou pohádku zadanou od diváků. Rozhodčí předem určuje handicapy a jednotlivé postavy, které hráči hrají. 
Pohádku musí všichni hráči znát, v opačném případě je vhodné např. nechat diváka odříkat celou pohádku. 
 
\subsection{Varianty} \begin{itemize}
\item Pohádka s charakterovou vadou - hráči mají zadanou pro svou postavu výraznou emoci/charakterovou vlastnost (podobně jako do \odkaz{Schizofrenie}{schizofrenie})
\end{itemize}
 
 
 
 
\needspace{5cm} \section{Jestli víš, co tím myslím} \label{jestli víš, co tím myslím} \katabox{Prostředí 
}{2+  
}{2  minuty} 
 
Každé slovo může získat zcela nový význam, \textbf{jestli víte, co tím myslím}{}. 
 
\subsection{Průběh} Dva hráči rozehrávají scénu v daném prostředí, ale každou svou \odkaz{repliku}{replika} zakončí  
větou jestli víš, co tím myslím. Všechny věty tak získávají jakousi "dvousmyslnost", která se ještě podpoří volbou slov před tímto zakončovačem repliky.   
 
 
 
\needspace{5cm} \section{Kaskadéři} \label{kaskadéři} \katabox{Téma 
}{4 (dva páry) 
}{orientačně 3 min} 
 
Přepínání mezi akčním a všedním životem 
 
\subsection{ Průběh }  
Hrají 4 hráči, najednou ale vždy jen dva. Jedna dvojice hraje civilní příběh na téma diváků, ale jakmile rozhodčí zakřičí “kaskadéři”, tak se vymění za druhou dvojici, která pokračuje ve stejném příběhu a stejných postavách, ale vše je akční, přeháněné a zběsilé, jakoby kaskadéři zaskočili za ty náročné části příběhu. Po chvíli rozhodčí zakřičí “civil” a vrací se první dvojice. Takto se civilní dvojice a kaskadérská několikrát během improvizace vystřídají. 
 
\subsection{ Běžné chyby } \begin{itemize}
\item  Ztráta napojení s druhou dvojicí
\end{itemize}
 
 
 
\needspace{5cm} \section{Kramářská píseň} \label{kramářská píseň} \katabox{Téma 
}{libovolný (4 a více) 
}{4 minuty} 
 
\textbf{Kramářská píseň}{} sloužila předkům jako bulvární zdroj informací. 
 
\subsection{Průběh} Jeden hráč se stává zpěvákem, na jeho tlesknutí se velmi rychle mění \odkaz{živý obraz}{živý obraz} tvořený ostatními hráči. 
Hráč zpěvem popisuje děj, a obrazy, mezi slokami obvykle přechází do vyvolavačské komunikace s diváky ( Ale to ještě není vše, přijďte blíž a uvidíte, přijďte blíž a uslyšíte.) 
 
Pro zjednodušení této obtížné kategorie se občas používá melodie z "Pes jitrničku sežral" 
 
\subsection{Varianty} \begin{itemize}
\item  Na zápase mohou být dva zpěváci, kteří se střídají po sloce a promluvě.
\end{itemize}
 
\subsection{Viz také} Podobná je kategorie \odkaz{Diapozitivy}{diapozitivy}. 
 
 
 
 
 
\needspace{5cm} \section{Muzikál} \label{muzikál} \katabox{Slovní téma 
}{2 - 5 
}{5-15 minut} 
 
Muzikál je improvizační kategorie založena na příběhu (pro lepší přehlednost zpravidla jednoho hlavního hrdiny) do kterého jsou vkládány písně, které posunují děj vpřed. Důležitá a hojně opomíjena je složka taneční, choreografická. Není ovšem nutné stále jen zpívat (narozdíl od kategorie \odkaz{opera}{opera}) v muzikálu jsou stejně nutné i dialogy. 
 
\subsection{Jak na to}  
Důležité je mít příběh (nejlépe hlavní hrdina někam musí jít a něco udělat aby něco se stalo). Často se hledí hlavně na zpěv, ten však je pouze prostředek ke sdělení nějaké informace (proto když vám prostě nejde rým tak tam nerýmujte, ale veďte příběh dál). Tip na tvoření choreografií: postavě se tak, aby jeden byl kousek před ostatními ten pak udává pohyby které ti za ním zrcadlí (vytváří tak efekt nacvičených choreografií) 
 
{{Todo| 
\begin{itemize}
\item  Doplnit průběh
\item  Doplnit video
\end{itemize}
}} 
 
 
 
 
 
\needspace{5cm} \section{Občanky} \label{občanky} \katabox{Občanské průkazy od diváků podle počtu hráčů 
}{3-4 
}{volný, dokud nedojdou nápady} 
 
Kategorie, ve které herci ztvární postavy vybraných diváků. 
 
\subsection{Průběh}  
Hraje 4-5 hráčů, každý si z diváků vybere jednoho člověka, který má stejné pohlaví, jako hráč. Divák mu půjčí občanku, řidičák, nebo pas a hráč má pak krátký čas na “nastudování” informací, jelikož poté hraje v improvizaci právě toho diváka, kterého si vybral. Je zajímavépřejmout i charakteristiky/vystupování postavy, kterou hrajeme. 
 
 
 
\needspace{5cm} \section{Oprsklá} \label{oprsklá} \katabox{téma/prostředí/vztah/situace 
}{2 
}{neomezen 
} 
 
V kategorie \textbf{Oprsklá}{} jsou hráči opravdu překvapeni každou informací, kterou se dozví. 
 
\subsection{Příprava} Každý hráč obdrží naplněnou sklenici vody.  
 
\subsection{Průběh} Kategorie spočívá v tom, že si navzájem sdělují informace, které druhého překvapí. Výrazem překvapení je následné vyprsknutí vody na spoluhráče. Hraje se, dokud je v půllitru nějaká voda.  
 
\subsection{Úklid} Po kategorii je scéna i hráči mokrá. Doporučujeme uklidit. 
 
\subsection{Běžné chyby} \begin{itemize}
\item Hráč ve své větě, větách nepřinese žádnou informaci.
\item Hráč nereaguje udiveně.
\item Hráč neumí vodu vyprsknout (doporučujeme natrénovat)
\end{itemize}
 
\subsection{Viz také} Kategorie je částečně podobná \odkaz{Co tím chceš jako říct}{co tím chceš jako říct}, hodí se zejména na horké letní festivaly. 
 
 
 
\needspace{5cm} \section{Orákulum} \label{orákulum} \katabox{otázka 
}{všichni 
}{volný} 
 
V této kategorii se z hráčů stává mýtické všeznalé stvoření, kterého se diváci ptají na záludné otázky. 
 
\subsection{ Průběh } Hráči se k sobě přitisknou a dohromady odpovídají na otázky z publika. Diváci se ptají na cokoliv, často i řečnické otázky. Orákulum by mělo být schopné na ně originálně a jednohlasně odpovědět. 
 
\subsection{ Běžné chyby } \begin{itemize}
\item Dlouhé odmlky
\item Všichni mluví přes sebe
\item Orákulum má viditelně jednoho hlavního mluvčího, po kterém všichni opakují
\end{itemize}
 
 
 
 
\needspace{5cm} \section{Příběh po slovech} \label{příběh po slovech} \katabox{Téma 
}{libovolný,minimálně 2 
}{volný} 
 
 
Napojení improvizátoři vypráví příběh po slovech. 
 
 
\subsection{ Průběh }  
Herci stojí vedle sebe v řadě a po jednom slově vyprávějí příběh. 
 
Jedno ze základních cvičení, které se používá jako součást kategorií \odkaz{Dopis}{dopis} nebo \odkaz{Divadlo poezie}{divadlo poezie}, ale také samostatně. 
 
 
\subsection{ Varianty }  
\begin{itemize}
\item  Diváci zadají ponaučení místo tématu. Příběh skončí přesnou citací zadaného ponaučení.
\end{itemize}
 
\subsection{ Běžné chyby }  
\begin{itemize}
\item  Přemýšlení nad dalším slovem a sebehodnocení
\item  Přehnaná kreativita
\end{itemize}
 
 
 
 
 
 
 
\needspace{5cm} \section{Režiséři} \label{režiséři} \katabox{Téma pro každého z režisérů}{neomezen}{neomezen} 
 
V kategorii \textbf{Režiséři}{} též známé jako \textbf{Super scene}{} několik hráčů v roli režiséra představí divákům své filmy. Diváci mají poté možnost rozhodovat o tom, který film chtějí vidět. vyberte si, který chcete vidět. 
 
\subsection{Průběh} Vyčlení se hráči, kteří se stávají režiséry. Každému z režisérů je zadán název jeho filmu (téma od diváků) a jeho žánr. Každý režisér má poté max. 30s na krátkou slovní upoutávku tohoto filmu, kterým se snaží diváky navnadit. Po představení všech filmů diváci hlasují o tom, který film vidět vůbec nechtějí. Z ostatních filmů se poté odehraje jejich první část, během níž může režisér komentovat děj filmu zákulisními informacemi nebo děj přímo ovlivňovat. Každý režisér první část svého filmu ukončí. Ke slovu se opět dostává \odkaz{MC}{mc} a diváci. Opět je zde hlasování po kterém opět jeden z filmů vypadne. Takto kategorie pokračuje až do konce jediného celého filmu. 
 
 
\subsection{ Varianty } \begin{itemize}
\item  \odkaz{Dva začátky}{dva začátky}
\end{itemize}
 
 
 
\needspace{5cm} \section{Scénář / Souboj s knihou} \label{scénář / souboj s knihou} \katabox{Připravená kniha či scénář 
}{2 
}{3 minuty} 
 
 
V této  kategorii  vám hráči sehrají příběh, který bude částečně shodný s  některým známým dílem. 
 
\subsection{Varianty začátku} \begin{itemize}
\item \textbf{Scénář}{} - Jeden hráč dostane od \odkaz{MC}{mc} přidělený kus scénáře - dialogu, ve kterém je vše co říká jedna postava vyškrtáno. Hráč musí používat pouze věty z dialogu té druhé postavy a to v původním pořadí.
\item \textbf{Souboj s knihou}{} - Jeden hráč dostane přidělenou knihu.  Všechny věty, které pronese, musí být \odkaz{repliky}{replika} z této knihy.
\end{itemize}
 
\subsection{Průběh} Sehraje se scénka, kdy druhý hráč musí citlivě reagovat a pečlivě poslouchat prvního hráče, v případě Scénáře navíc odhadovat, co se asi v scénáři může stát, aby reakce druhého hráče vypadaly logicky. 
 
\subsection{Viz také} \begin{itemize}
\item  \odkaz{Dvě repliky}{dvě repliky}
\end{itemize}
 
 
 
\needspace{5cm} \section{Shakespeare pozpátku} \label{shakespeare pozpátku} \katabox{Téma 
}{neomezený 
}{cca 10 minut} 
 
Uvidíme pravé a hutné drama jak od Shakespeara, ale celé \textbf{pozpátku}{}. 
\subsection{Průběh} Hraje se příběh rozdělený do 5 dějství pojmenovatelných dle modifikovaných částí \odkaz{příběhu}{příběh}.  
Tyto dějství se ale hrají v opačném pořadí - začínáme tedy od konce.  
\begin{itemize}
\item Tragedie
\item Konflikt/Intriky
\item Zápletka
\item Láska/Vztahy
\item Úvod
\end{itemize}
 
V rámci dějství jsou scény řazeny standardně, chronologicky.  
 
 
 
\needspace{5cm} \section{Statusy} \label{statusy} \katabox{Téma a lístky s čísly 
}{4 
}{volný} 
 
 
Hádání statusu. 
 
 
\subsection{ Průběh }  
Hrají 4 hráči a každý má na čele vylosovaný lístek s číslem (1-10), který ale neviděl a vylosoval si ho poslepu. Poté se zadá prostředí, kde se často vyskytují lidé s různými společenskými postaveními a přichází první hráč. Ten by měl začít neutrální činností, jelikož neví, jaké číslo (jaký status) má. Obecně platí, že čím vyšší číslo, tím vyšší společenské postavení. Hráči po skončení improvizace tipují, jaké číslo mají na čele. 
 
 
 
\needspace{5cm} \section{Stojí, leží, sedí} \label{stojí, leží, sedí} \katabox{Téma}{3}{2 min} 
 
V Kategorii \textbf{stojí, leží, sedí}{} hrají 3 hráči a vždy musí jeden z hráčů stát, jeden sedět a poslední ležet. 
 
\subsection{ Průběh } Jsou vybrání tři hráči, je jim zadáno téma a dána k dispozici jedna židle. Hráči zaujmou dané pozice - jeden zůstane stát, druhý si sedne a třetí lehne. Poté rozehrají improvizaci na dané téma a přitom musí zachovávat, aby byly stále zastoupeny všechny tři pozice. Hráči mohou své pozice během improvizace měnit. Není určeno pořadí v jakém musí pozice jít po sobě. Ležící hráč si může rovnou stoupnout a stejně tak i obráceně.  
 
\subsection{ Časté chyby } \begin{itemize}
\item Hráči nedávají pozor a opomenou jednu pozici
\end{itemize}
 
 
 
 
\needspace{5cm} \section{Tříhlavý song} \label{tříhlavý song} \katabox{první verš básně}{3}{cca 3 min} 
V této takřka nehrané kategorii hráči zazpívají píseň po jednom slově. 
\subsection{Průběh} Hráči se k sobě přitisknou a dohromady zazpívají píseň s tím, že se střídají po jednom slově. První verš, vyžádaný od publika, na začátku zopakují, čímž se i dostanou do nějakého tempa.  
 
\subsection{Běžné chyby} \begin{itemize}
\item Dlouhé odmlky
\item Běžné veršovací chyby
\end{itemize}
 
\subsection{Viz také} \begin{itemize}
\item Kategorie \odkaz{Opilecká píseň}{opilecká píseň} - zpěv se strřídáním po jednom verši.
\item Cvičení \odkaz{Příběh po jednom slově}{příběh po jednom slově}
\item Rozcvička \odkaz{Vtip}{vtip}
\end{itemize}
 
 
 
 
\needspace{5cm} \section{Tři židle} \label{tři židle} \katabox{není nutné 
}{3 
}{5-15 minut} 
 
Každá skupinka stojících hráčů má svůj příběh. 
 
\subsection{Průběh} Připraví se 3 židle, na které se neutrálně usadí stejný počet herců. 
Hráči dle svého rozhodnutí vstávají, vždy však stojí alespoň jeden hráč. Každá možná podskupinka vypráví nějaký příběh, jeho vyprávění je však dočasně přerušeno, když se postaví nebo si sedne některý z hráčů.   
Kategorie končí v momentě, kdy jsou všechny příběhy dovyprávěny, hráči se proto snaží je dovyprávět cv podobný okamžik. 
 
\subsection{Princip} Kategorie je divácky zábavná díky změnám dynamice vyprávění. Herci nemají žádné \odkaz{askfor}{askfor}, mají proto úplnou svobodu ve výběru  žánru, formy a délce příběhů. 
 
\subsection{Varianty} \begin{itemize}
\item  Dvě židle (3 příběhy)
\item  Čtyři židle (15 příběhů)
\item   Pět židlí (31 příběhů)
\end{itemize}
 
\subsection{Podobné} \begin{itemize}
\item  \odkaz{Spoon river}{spoon river}
\item  \odkaz{Toaster}{toaster}
\end{itemize}
 
\subsection{ Může se plést s } \begin{itemize}
\item  \odkaz{Židle}{židle}
\item  \odkaz{Dvě židle}{dvě židle}
\end{itemize}
 
 
 
 
 
 
 
\needspace{5cm} \section{Věty, které by neměly zaznít} \label{věty, které by neměly zaznít} \katabox{Situace/činnosti (3 - 5) 
}{libovolný 
}{neomezený} 
 
 
Hráči říkají věty, které by v určitých situacích nebo na určitých místech neměly zaznít. 
 
 
\subsection{ Průběh } Podobně jako v kategorii \odkaz{1000 způsobů jak}{1000 způsobů jak} vybere \odkaz{MC}{mc} od publika několik situací např. první kroky dítěte, přijímací pohovor, na operačním sále, atd. \odkaz{Hráči}{hráč} se postaví do půlkruhu a poté vždy ten s nápadem na větu vyběhne na forbínu, kde tuto \odkaz{repliku}{replika} prezentuje a opět se zařadí do půlkruhu. Situace či prostředí se nechává do vyčerpání nápadů nebo na pokyn \odkaz{MC}{mc}. 
 
\subsection{ Varianty } \begin{itemize}
\item  \odkaz{1000 způsobů jak}{1000 způsobů jak}
\item  \odkaz{Metafory}{metafory}
\item  Některé věty se hodí do více situací a hráči toho hojně využívají.
\end{itemize}
 
\todo{doplnit video} 
 
 
 
\needspace{5cm} \section{Vtip} \label{vtip} \katabox{nemusí být 
}{alespoň dva}{cca 1 min 
} 
Hráči vám po jednom slově řeknou zbrusu nový vtip. 
\subsection{Průběh} Hráči se seřadí vedle sebe, jeden začne mluvit s tím, že každý říká jen jedno slovo, drží se pořadí a odříká se vtip.  Vtip končí pointou a úklonou. 
\subsection{Varianty} \begin{itemize}
\item Předložky a spojky se nepočítají jako slovo.
\end{itemize}
 
\subsection{Viz také} \begin{itemize}
\item \odkaz{Tříhlavý song}{tříhlavý song}
\item \odkaz{Příběh po jednom slově}{příběh po jednom slově}
\end{itemize}
 
 
 
\needspace{5cm} \section{Zemři} \label{zemři} \katabox{Téma 
}{3 - 20 
}{do vyčerpání} 
 
Viděli jste někdy chobotnici vyprávět příběh? 
 
\subsection{ Průběh }  
Hrají obvykle všichni hráči na place a na tuto improvizaci si stoupnou do řady za sebe čelem k divákům. Vždy ten, kdo je v řadě první (nejblíž divákům), tak vypráví poutavý příběh na zadané téma. Dělá u toho výrazná gesta a pohyby, využívá celou šířku jeviště a ostatní jeho pohyby přesně napodobují. Kdykoliv může kterýkoliv z hráčů (kromě toho prvního) zakřičet “zemři” a tím přesunou prvního hráče na konec řady a jeho místo a také vyprávění přebírá ten hned za ním. 
 
\subsection{ Tipy } \begin{itemize}
\item  Někdy je vtipné, když se “zemři” řve často a někteří hráči se tak třeba ani nedostanou k vyprávění.
\end{itemize}
 
 
 
\needspace{5cm} \section{Zvuky} \label{zvuky} \katabox{Téma 
}{2 
}{2-3 minuty} 
 
 
V kategorii \textbf{Zvuky}{} vám hráči zahrají situaci, kterou dotvoří vybraní diváci svými zvuky. 
 
 
\subsection{Průběh} \odkaz{MC}{mc} na začátku kategorie vybere z diváků dva dobrovolníky. Každý z nich dostane mikrofon a je mu přidělen hráč, kterému bude divák dotvářet zvuky podle jeho činností. Je dobré ještě před začátkem kategorie si jeden zvuk vyzkoušet nanečisto např. hráč otevírá dveře a divák doplní vrzavý zvuk pantů. V samotné kategorii hráči hrají volně na téma a snaží se svým "zvukařům"{} nabízet co nejvíce akcí, které se dají ozvučit. Divák má možnost být pouze pasivním a doplňovat zvuky akcí, které se dějí ale pokud má nápad může také zvuky nabídnout hráčům další podnět (mluvící papoušek, hlášení z rádia, atd). 
 
\subsection{ Varianty } \begin{itemize}
\item Zvuky vydávají všichni diváci najednou, bez mikrofonů
\end{itemize}
 
\subsection{ Běžné chyby } \begin{itemize}
\item Diváci zapomenou, koho mají zvučit.
\item Divákovi je dán vypnutý mikrofon.
\end{itemize}
 
 
 
 
\needspace{5cm} \section{1000 způsobů jak} \label{1000 způsobů jak} \katabox{Činnosti (3-5) 
}{libovolný, obvykle se zapojí všichni 
}{volný, dokud nedojdou nápady} 
 
 
Nebo taky \textbf{Tisíc způsobů jak}{}. Hráči hledají různé způsoby, jak vykonat činnost, například \odkaz{výměna žárovky}{výměna žárovky}. 
 
 
\subsection{Průběh}  
\odkaz{MC}{mc} vybere od publika činnosti. Poté se hráči postaví do půlkruhu a střídají se. Každý hráč, který vyběhne, předvede volně nějaký svůj alternativní způsob, jak vykonat zadanou činnost. Obvykle MC po chvíli činnost vymění, zadá jinou. 
 
\subsection{ Varianty } \begin{itemize}
\item  \odkaz{Věty, které by neměly zaznít}{věty, které by neměly zaznít}
\item  \odkaz{Metafory}{metafory}
\item  \odkaz{1000 hitů o...}{1000 hitů o...}
\end{itemize}
 
\subsection{ Běžné chyby } \begin{itemize}
\item Sklouzávání k záchodovému humoru
\item Přehnané použití dynamitu
\end{itemize}
 
 
 
 
 
\needspace{5cm} \section{4, 2, 1} \label{4, 2, 1} \katabox{Téma nebo prostředí 
}{4 
}{1 min} 
 
 
Nebo taky \textbf{čtyři, dva, jedna}{}. Improvizátoři ztvární příběh, který opakují se zmenšeným počtem herců, ale při zachování počtu postav. 
 
 
\subsection{ Průběh }  
Hrají 4 hráči na zadané téma víceméně volnou improvizaci, kde ale dávají důraz na pohyb a akci. Tato improvizace je krátká, max 2 minuty. Po skončení diváci vyberou 2 hráče, kteří ten stejný příběh 4 postav zahrají sami (střídají postavy, přeskakují z jedné do druhé). Posléze diváci vyberou jen jednoho, který celý příběh odehraje naposledy a úplně sám (opět hraje všechny 4 postavy). 
 
Může se plést s \odkaz{Počet slov}{počet slov}, přezdívanou často '1, 3, 5, 7' 
 
\subsection{ Běžné chyby } \begin{itemize}
\item  Postavy příliš mnoho mluví (děj se nedá zapamatovat)
\item  Úvodní scénka v počtu 4 herců je příliš dlouhá
\end{itemize}
 
\subsection{ Podobné kategorie }  
\begin{itemize}
\item  \odkaz{Polovina času}{polovina času}
\end{itemize}
 
 
 

\end{document}

\chapter{Rozcvičky}\label{rozcvičky}
\label{:kategorie:rozcvičky}

\documentclass[main.tex]{subfiles}
 
\begin{document}\needspace{5cm} \label{kategorie:rozcvičky} Některé kategorie se často používají v průběhu \odkaz{veřejné rozcvičky}{veřejná rozcvička}. 
Leckterá z \odkaz{kategorií na improshow}{:kategorie:kategorie na improshow} se dá také použít. 
 
 
 
 
\end{document}
%\needspace{5cm} \section{Běh přes překážky} \label{běh přes překážky} Hráči stojí na zadní hraně jeviště, čekají na start závodu. Na pokyn \odkaz{MC}{mc} Připravit se, ke startu, pozor, teď   zakleknou do bloků a odstartují běh na místě. V průběhu běhu překonají (ideálně všichni najednou)přeskokem tři překážky, po té třetí se okamžitě přepínají do zpomaleného pohybu, během kterého se snaží dosáhnout přední hrany jeviště a zároveň zabránit ostatním hráčům tohoto cíle dosáhnout pomocí hraných útoků, úderů apod. 
 
\subsection{Varianta - štafeta} Tato varianta se obzvlášť hodí při malé scéně/velkém počtu hráčů. 
Na zadní hraně jeviště jsou dvě řady hráčů, přední zakleknutá na zemi, druhá stojící. Po odstartování běží na místě pouze zadní řada, po třetí překážce předávají štafetový kolík přední řadě která pokračuje dál, již ve zpomaleném režimu. 
 
 
 
 
\needspace{5cm} \section{Emocionální autobus} \label{emocionální autobus} \katabox{nic 
}{1+3-6 
}{neomezen} 
 
V našem \textbf{Emo autobusu}{} přebírá řidič a ostatní hráči (cestující) emoci/charakter toho, kdo právě nastoupil jako poslední.  
 
\subsection{Průběh} První na scéně je řidič, přichází k autobusu v nějaké \odkaz{emoci}{emoce}/charakterové vlastnosti/charakterové vadě a vydává se na cestu. Ve chvíli kdy zastaví a nabere pasažéra, přebírá od pasažéra jeho emoci/charakter. Při nabírání dalšího pasažéra přebírá jeho emoci i ostatní cestující. Když jsou v autobusu všichni hráči, začnou postupně vystupovat v opačném pořadí než nastupovali s tím, že autobus přebírá původní charakter toho, kdo bude vystupovat. 
 
\subsection{Varianta} \begin{itemize}
\item Když je autobus plný, všichni hráči vystoupí na další zastávce. Postupně, zachovává se stejný princip, jen se zrychlí konec kategorie.
\end{itemize}
 
\subsection{Scéna} Bývají připravené židličky. 
 
 
\needspace{5cm} \section{Horská dráha} \label{horská dráha} Hráči na pokyn nastoupí do vozíku horské dráhy - např. 6 lidí ve třech řadách, 
již při nastupování dbají na prostorové rozmístění sedadel a vrátek. 
Na hráče se zaklapnou zábrany, a začíná jízda, kterou hráči ztvárňují nakláněním do stran či dopředu/dozadu. 
 
\needspace{5cm} \section{Hospodská rvačka} \label{hospodská rvačka} Hráči rozehrají jednoducho scénu z baru/hospody. V jednu chvíli se scéna zvrtne ve rvanici, v tu chvíli všichni pokračují ve \odkaz{zpomaleném pohybu}{zpomalený pohyb}, dokud nezůstane poslední přeživší (stojící) hráč. 
 
 
 
\needspace{5cm} \section{Jsem, beru} \label{jsem, beru} Klasické cvičení, používané často na tréninzích i ve veřejných rozcvičkách. 
\subsection{Průběh} První hráč nabíhá, zaujme nějakou pózu, \odkaz{štronzo}{štronzo} , a pojmenuje se. 
 
Nabíhá druhý hráč, a zapojí se do obrazu, jak pozicí tak ztvárněným předmětem.  
 
Dtto třetí hráč. 
První hráč slovy Beru ... odvolá sebe a jednoho z dalších hráčů, ten následně začíná další kolo. 
\subsection{Varianty} \begin{itemize}
\item Odvolávání tlesknutím jednoho z čekajících hráčů, který následně tvoří základ dalšího obrazu
\item Veršovaně
\item Předmět, problém, řešení
\item Zapojí se všichni
\item Ze zadaného prostředí
\item Všichni se musí dotýkat
\item Beze slov
\item S opakováním  pohybu dokreslujícího danou věc
\end{itemize}
 
 
 
 
 
 
\needspace{5cm} \section{Krátká báseň} \label{krátká báseň} \katabox{téma či název básně 
}{4 
}{volný, max minuta} 
 
 
Hráči vám poví krátkou báseň na zadané téma 
 
 
\subsection{ Průběh } Hráči postupně přicházejí a rýmují na sebe \odkaz{ABAB}{abab}. 
Každý hráč řekne jeden verš.  
 
Rozcvička se obvykle opakuje dvakrát nebo třikrát. 
\subsection{ Varianty } \begin{itemize}
\item  \odkaz{MC}{mc} určí počet slok, po zaznění každé slok obvykle \odkaz{hudebník}{hudebník} proloží báseň krátkou vyhrávkou. Pořadí hráčů se v dalších slokách nemusí držet.
\end{itemize}
 
\subsection{ Běžné chyby } \begin{itemize}
\item  Běžné veršovací chyby
\end{itemize}
 
 
 
 
\needspace{5cm} \section{Orchestrion} \label{orchestrion} Na scénu přichází první hráč a začíná opakovat rytmický či melodický motiv. Připojí se druhý a následně i třetí hráč. Ve chvíli, kdy přijde čtvrtý hráč, první hráč odchází (ideálně po postupném utlumení vlastního motivu) a je připrave na scénu naběhnout znova. 
Po chvíli je možné jako vlastní motiv zpívat frázi z jednoho či dvou slov, držící se vytvořeného témata. 
 
Je vhodné, když alespoň jeden hráč drží nějaký základní rytmus, který pak všichni ostatní využívají. 
 
 
 
 
\needspace{5cm} \section{Roztočený brankář} \label{roztočený brankář} Základní improvizační cvičení na procvičení asociací, reakce a komunikace. 
 
Hráči stojí ve větším kruhu či v půlkruhu. Jeden hráč vstoupí do středu či před stojící hráče a začne se pohybovat. Pohyb není nijak omezen, může zaujímat rozličné pozice (sedat si, lehat si, postavit se na hlavu, ...), ale měl by být stále v pohybu. Kdokoli z ostatních hráčů může kdykoli během jeho pohybu tlesknout, což je pro pohybujícího pokyn \odkaz{štronzo}{štronzo}. Hráč, který tleskl zaujme vůči svému kolegovi též nějakou pozici, kterou mu pozice \odkaz{asociuje}{asociace} a rozehraje situaci. Hráč opouští štronzo a pokračuje také v rozehrané situaci. Hráči by měli co nejdříve objasnit kdo jsou, kde jsou a co tam dělají. Jakmile je toto objasněno, může opět tlesknout kdokoli z přihlížejících hráčů, což je opět povel tentokrát pro oba hrající hráče \odkaz{štronzo}{štronzo}. Hráč, který tentokrát tleskl přichází ke dvojici, dotykem určí hráče, který odejde, zaujme opět pozici k hráči, který zůstává a opět rozehrají jinou situaci. Toto cvičení se dá libovolně dlouho opakovat. 
 
 
 
 
\needspace{5cm} \section{Samuraj} \label{samuraj} Ze skupiny hráčů je vybrán samuraj, může se sám přihlásit nebo je určen cvičícím, který hráči předává imaginární samurajský meč. Ostatní hráči jsou samurajem ohroženi a snaží se před ním uniknout, ale současně se drží jako semknutá skupina. Zpravidla je skupina v jednom rohu prostoru a samuraj v protilehlém. Samuraj se snaží teatrálně a relativně pomalými pohyby meče ostatní hráče zasáhnout. Může mířit na hlavy a hráči se skrčí, může mířit na nohy a hráči nadskočí, případně sekne svisle do skupiny a skupina se rozestoupí, aby zásahu unikla. Pozice samuraje a skupiny se mohou během cvičení také prohodit, když po nezdařeném zásahu skupina proběhne na místo samuraje a samuraj zaujme původní místo skupiny. Závěr cvičení může být rozsekání skupiny na kousíčky zdařenými zásahy samuraje, nebo tím, že samuraj pod tíhou svého neúspěchu (žádný zasažený hráč) spáchá harakiri. 
 
 
 
 
 
\needspace{5cm} \section{Souboj otázek} \label{souboj otázek} V tomto cvičení se hráči rozdělí na dvě skupiny, které stojí proti sobě v zástupu. Cvičící zadá prvním dvěma hráčům prostředí, ve kterém jsou a ti spolu pak musí komunikovat pouze otázkami. První, kdo nepoloží otázku vypadává, jde na konec řady a na jeho místo nastupuje další hráč. Prostředí může cvičící střídat podle svého uvážení. 
 
\subsection{ Běžné chyby } \begin{itemize}
\item Hráči vrství stejné otázky na sebe - "Ty si myslíš, že já si myslím, že ty si myslíš, že jsi přišel pozdě?"
\item Hráči ve svých otázkách nepřináší nic nového.
\end{itemize}
 
  
 
 
 
\needspace{5cm} \section{Stroj} \label{stroj} Skupině je zadáno co bude daný stroj vyrábět nebo k čemu slouží. První hráč poté ztvární pohybující se součástku a doplní ji nějakým ozvučením. Ostatní hráči se postupně přidávají jako další součástky až vznikne celý stroj. Součástky by se měly dotýkat nebo navazovat na akci součástky, která je vedle ní. Když je stroj hotov cvičící zastaví jeho pohyb, ale hráči zůstávají v pozicích a pojmenují své součástky. Poté cvičící může stroj opět spustit a zkoušet různé rychlosti chodu stroje 50\%{}, 100\%{}, 150\%{}, 200\%{} až se stroj porouchá či exploduje. 
 
\subsection{ Varianty } Při velkém prostoru se dá vyzkoušet nekonečný stroj. Hráči opět představují jednotlivé součástky stroje, ale při zapojení všech hráčů ten první přebíhá za posledního a vytváří další originální součástku, atd. V tomto případě odpadá pojmenovávání součástek. 
 
 
 
 
 
\needspace{5cm} \section{Thai Chi} \label{thai chi} Hráči se rozestaví do čtverce (hrany, vrcholy i vnitřek) otočeni všichni stejným směrem. 
Všichni se pohybují stejně, pohyby,cviky i natočení určuje ten, kdo je na vrcholu směrem, kterým se všichni dívají. 
Cílem je vypadat stejně jako lidé cvičící Thai Chi. 
Po chvíli se začnou hráči náhodně střídat ve vyprávění příběhu po krátkých částech vět. 
\subsection{Varianty} \begin{itemize}
\item Čtení dopisu
\item Báseň
\item Vyprávění příběhu
\end{itemize}
\subsection{Běžné chyby} \begin{itemize}
\item Vůdčí hráč se pohybuje příliš rychle/ zaujímá pózy, které ostatní nezvládnou.
\item Hráči nepozorně sledují a kopírují pohyby
\item Vyprávění není plynulé
\end{itemize}
  
\needspace{5cm} \section{Židle} \label{židle} Uprostřed scény stojí židle, ke které se na písknutí ze stran rozeběhnou dva hráči. Ve chvíli, kdy se židle dotknou, na okamžik \odkaz{strnou}{štronzo} a pak sehrají krátkou scénu vycházející z pozic, ve kterých skončili. 
 
\subsection{ Viz také } \begin{itemize}
\item  \odkaz{Roztočený brankář}{roztočený brankář}
\end{itemize}
 
\subsection{ Může se plést s } \begin{itemize}
\item  \odkaz{Dvě židle}{dvě židle}
\item  \odkaz{Tři židle}{tři židle}
\end{itemize}
 
 



\chapter{Příprava zápasu}
\documentclass[main.tex]{subfiles}
 
\begin{document}\needspace{5cm} \label{příprava zápasu} Je li z vás \textbf{Garant akce}{} - organizátor, zde se dozvíte pár věcí, na které je dobré myslet. 
\section{Rozdělení rolí účinkujících} \section{Garant akce} Garant akce je člověk, který zašťiťuje celkovou organizaci. Má veškeré informace. Rozděluje role ostatním a vykoná všechno, co nemůže rozdělit. Typicky na veřejných akcích zastává roli konferenciéra, a nebo uvaděče. Pokud je člověk přizván pořadatelem většího celku, jako jsou festivaly, zápas druhého týmu, je výhodou pokud přímo garant akce sestavuje tým hráčů. Dostává tak vedle velké zodpovědnosti i dostatečné kompetence. 
 
\section{Trenér} Člověk v roli \odkaz{trenéra}{trenér} má za úkol připravit trénink před akcí. Tato role se často slučuje s rolí rozhodčího. 
\section{Hudebník} Pokud není formátem akce stanoveno jinak, tak tato role hudebně doprovází improvizaci po celou dobu představení. 
Je důležité, aby byl \odkaz{hudebník}{hudebník} dostupný už na tréninku před akcí. 
Pozor na dopravu nástroje na místo. Může se stát, že hudebník nebude schopen svůj nástroj sám dopravit na místo. 
 
\section{Konferenciér nebo rozhodčí} Tyto dvě role se řídí manuálem \odkaz{konferenciéra}{konferenciér} a manuálem \odkaz{rozhodčího}{rozhodčí}, které zatím neexistují, pročtěte si heslo \odkaz{Zápas}{zápas}. 
Je dobré, aby věděli 
\begin{itemize}
\item Kdo hraje (seznam \odkaz{hráčů}{hráč})
\item Kdo je \odkaz{hudebník}{hudebník}
\item Od kdy do kdy se hraje
\item Kategorie, které se mohou nebo budou hrát
\item Téma představení, pokud nějaké je
\end{itemize}
\section{Zvukař} Zvukař má na starosti ovládání hlasitosti hudebních nástrojů a mikrofonů po celou dobu představení. Jeho role se obvykle spojuje s osvětlovačem. Je důležité, aby věděl všechny technické parametry prostoru. 
V případě nedostupnosti některých pomůcek, jako například mikrofonů, má zvukař na starosti zařídit jejich dostupnost. 
Pokud bude potřeba přehrávat nějaká média (například empétrojky a nebo pouštění videí na plátno), řeší technické problémy s tím spojené. 
Může požadovat dopravu techniky. 
\section{Osvětlovač} Osvětlovač má na starosti ovládání světel a vedlejšího osvětlení prostoru po celou dobu představení. Jeho role se obvykle spojuje se zvukařem. Je důležité, aby věděl všechny technické parametry prostoru. 
Může požadovat dopravu techniky. 
\section{Pomocňák nebo pomocný rozhodčí} \odkaz{pomocní rozhodčí}{pomocný rozhodčí}  mají na starosti přípravné práce před začátkem představení a obvykle během tréninku před představením. Garant by se s nimi měl dohodnout, aby přijeli před začátkem. 
\section{Hráč} Hráč by měl mít na starost pouze dopravit sám sebe včas spolu s oděvem vhodným pro konkrétní akci. 
\section{Prostory k hraní} Prostor k hraní musí být zajištěn předem. K prostoru je potřeba vědět všechny technické věci 
\begin{itemize}
\item Velikost pódia včetně nepřemistitelných překážek
\item Dostupnost dvou mikrofonů
\item Jaké mají osvětlení
\item Jestli se dá úplně zhasnout v sále (\odkaz{vysílání FM}{vysílání fm}, \odkaz{Rozhlasová hra}{rozhlasová hra})
\item Kapacita hlediště
\item Časové možnosti – od kdy do kdy se může hrát
\item Dostupnost odposlechů
\item Sezení pro diváky
\item Plátno a projektor (\odkaz{Dabing filmu}{dabing filmu})
\end{itemize}
 
\section{Propagační nástroje} \begin{itemize}
\item Plakát
\item Akce na Facebooku a Google Plus, Twitter
\item Vstupenky
\item Newsletter
\end{itemize}
 
\section{Časová omezení} Organizace akce by včetně tisku a rozvěšení plakátů měla být dokončena nejméně týden před představením. 
\subsection{Sumarizace pro účinkující} Uvést v bodech 
\begin{itemize}
\item Kde se hraje
\item Jak se tam jede
\item Kdo hraje
\item Co si má kdo vzít sebou
\item Individuální úkoly
\item After party
\end{itemize}
\section{Improzápas} \subsection{Pomůcky} \begin{itemize}
\item Tabule skóre pro každý tým (čísla 0-9) a kolíčky na trestné body
\item Píšťalka a kazoo
\item Hlasovací kartičky v barvách týmových dresů (Je nutné domluvit dostatečně předem)
\item Papuče nebo míčky na házení
\item Papírky na témata
\item Půjčovací tužky na témata
\item Košík na témata
\item Kostým pro rozhodčího a pomocňáky
\end{itemize}
\end{document}

\chapter{Předzápasový trénink}
\documentclass[main.tex]{subfiles}\begin{document}
\needspace{5cm} \label{předzápasový trénink} Týmy se scházejí zpravidla 2 hodinky před zápasem, aby se na sebe mohli hráči, hudebník, rozhodčí a konferenciér naladit na předzápasovém tréninku. Je to také čas, kdy si týmy projdou kategorie, protože některé kategorie mají více variant. 
 
Předzápasový trénink by měl vést jiný než hrající člověk, protože pak on sám není rozehřátý jako ostatní a může se cítit mimo tým. 
 
Na tréninku by se měl \odkaz{trenér}{trenér} zaměřit na témata jako je společné naladění, důvěra mezi hráči, energii týmu, rozdýchání a rozmluvení, rozjetí fantazie a nápadů. Nemá smysl se zaobírat dovednostmi, které hráčům chybí, protože na to není před zápasem čas a hráči se mohou zablokovat, protože vidí, že jim nějaká činnost nejde. 
Některé týmy ani nezkouší kategorie kromě zpívané, protože, když se kategorie nepovede, může to vést ke špatně náladě týmu. Nevadí však, když si poslední cvičení vezme na konci tréninku samotný rozhodčí. 
 
Je vhodné, když cvičení na tréninku jsou hráčům známá, nemusíme tolik času trávit vysvětlování pravidel. 
Hráči i zpravidla vědí na co je cvičení zaměřené a mohou se věnovat jejich naplňování. 
 
 
\section{ Ukázka cvičení předzápasového tréninku: }  
 
\begin{itemize}
\item Honička se speciálními pravidly (na rozehřátí těla, hráči se dostávají do přítomnosti)
\item Dýchání ve dvojících a uvědomování si těžiště, kdy hráči držením těla mají společné těžiště (uvědomování si svého těla, těla partnera a těžiště pomáhá odstraňovat nervozitu, dýchání je důležité pro hlas, ve dvojicích také hráči získávají důvěru mezi sebou)
\item \odkaz{Gordický uzel}{gordický uzel} (cvičení zaměřený na kontakt, důvěru ve skupinu, že rozpletení společně vyřeší)
\item Společná masáž, vlastní masáž obličeje (společná masáž pomáhá zvyknout si na dotyk ostatních, uvolňuje svalové napětí, dělá dobře)
\item Hlasové cvičení (zajistí sílu hlasu a správné používání hlasivek)
\item \odkaz{Asociace}{asociace (cvičení)} (rozehřívá myšlení, cvičení může být obohaceno o pravidlo např., že všechna slova začínají na jedno písmeno apod.)
\item \odkaz{Jsem a beru}{jsem, beru} (do slov zapojujeme pohyb, pomáhá zapojovat obě hemisféry)
\item Počítání do dvaceti (pomáhá skupinu naladit se na společnou vlnu, zvládnout společně úkol a soustředění se)
\item Společná báseň (hráči si zkouší spolupráci, brát nápady od ostatních a posouvat je dál)
\item \odkaz{Ready Go}{ready go} (hráči si zkouší přejímat emoce a posouvat příběh dál)
\item Jamování (hráč se zaměřují na společný rytmus, doplňování hudby pomocí svého těla)
\item Společná píseň (naladění se na hudebníka, který si také může ověřit, jak jsou na tom hráči se zpěvem)
\item \odkaz{Kdo ukradl sušenky a ze stolu je vzal}{kdo ukradl sušenky a ze stolu je vzal} (rytmické a energetické cvičení na závěr)
\end{itemize}
 
Na závěr předzápasového tréninku hráči s \odkaz{rozhodčím}{rozhodčí} vybírají ze seznamu kategorií ty, které umí oba týmy, případně se vybírají varianty těchto kategoríí. 
 
 

\end{document}

\chapter{Cvičení}
%\documentclass[main.tex]{subfiles}\begin{document}
\needspace{5cm} \section{Asociace (cvičení)} \label{asociace (cvičení)} Základní a často používané improvizační cvičení na rozvíjení rychlé reakce a představivosti. 
 
\subsection{Průběh} Hráči stojí v kruhu. Ten který je na řadě, řekne slovo, které mu nějakým, libovolným způsobem navazuje (asociuje) na slovo předchozí. Na toto nové slovo pak pokračuje následující hráč. 
 
\subsection{Varianty} \begin{itemize}
\item říká se druhá asociace ještě přes jedno slovo (sériové)
\item říká se druhá asociace na příchozí slovo (paralelní)
\item jdou dvě (tři) slova najednou, vymýšlí se nová slova pro každou z těchto řad samostatně.
\end{itemize}
 
 
 
 
 
 
 
\needspace{5cm} \section{Bunny bunny} \label{bunny bunny} Jedná se o velice bláznivé cvičení, které má za úkol rozehřát hráč a nabít energií. Hráči stojí v kruhu. Všichni společně drží rytmus poklepáváním dlaněmi o stehna (a občas skandováním "umpa"{} či "umča"{} ). Jeden z hráčů vysílá signál předvedením hlavičky králíka pomocí prstů. Palcem, ukazováčkem a prostředníčkem klepne o sebe 2x směrem k sobě, 2x směrem hráči, kterému signál posílá. Pohyb doprovází slovem "bunny". Sousedící hráči podporují králíka pohybem rozevřených dlaní doleva a doprava na úrovni hrudníků a pohyb doprovází slovem "tyky", "taky". Oba sousedé předvádí stříhání králičích uších, dlaně mají nad hlavou a opět doprovází pohyb dohodnutým zvukem ("flappy"). 
 
Pokud chceme hru zamotat ještě více, může se ten, kdo byl králíčkem jako poslední, otočit jednou kolem osy s hláškou "I am not bunny any more!"{} či "Já už nejsem králíček", tento pohyb se však vykonává až v okamžiku, kdy hráč nemá žádnou jinou povinnost ("bunny", "tyky,taky", "flappy") . 
 
Běžné je také (podobně jako u dalších podobných her) zrychlování až do závratného tempa. 
 
 
 
 
\needspace{5cm} \section{Dřepy} \label{dřepy} {{Kategorie 
|hraci=libovolný,minimálně 5 
|cas=volný}} 
 
Vyřazovací hra pro větší počet účastníků, která probíhá na kola. Na začátku hráči stojí v prostoru, tak aby každý měl svoje místo. V každém kole mají hráči za cíl si dřepnout na bobek. Hru je potřeba řídit zvenku. 
 
\begin{itemize}
\item  Hráči, kteří se dřepnou ve stejný okamžik vypadávají
\item  Hráči, kteří se pohnou do dřepu ve stejný okamžik vypadávají.
\item  Poslední hráč, který si dřepne a nebo to nestihne, vypadává
\end{itemize}
 
Lektor verbálně určuje kdo vypadl. Vypadlí hráči okamžitě opouští hrací plochu. 
 
Po vyřešení kola si hráči ponechávají pozici a pouze si opět stoupnou. Pokračuje se dalším kolem, dokud nezůstane buď jeden vítěz a nebo všichni poražení. 
 
 
 
 
\needspace{5cm} \section{Dvě židle} \label{dvě židle} Dva herci vyjdou na jeviště, každý si nese svoji židli a náhodně zaplňují prostor. Trenér zaplňování ukončí tlesknutím nebo písknutím a oba hráči okamžitě položí židli a zaujmou k ní náhodnou pozici. Poté se rozehrává scéna, která vychází ze vzájemné pozice židlí a herců. 
 
\subsection{Může se plést s} \begin{itemize}
\item  \odkaz{Židle}{židle}
\item  \odkaz{Tři židle}{tři židle}
\end{itemize}
 
Cvičení je vhodné na \odkaz{vnímání prostoru}{vnímání prostoru} a na \odkaz{dávání si na čas}{dávání si na čas}. 
 
 
 
\needspace{5cm} \section{Gordický uzel} \label{gordický uzel} Kontaktní rozcvička, která začíná v kruhu. Postup je následující: 
 
\begin{enumerate}
\item  Skupina vytvoří kruh
\item  Všichni zavřou oči
\item  Každý natáhne do kruhu pravou ruku nahoru, levou ruku dolů a snaží se chytit ruku někoho jiného. Čeká se, dokud se všichni někoho nedrží
\item  Všichni otevřou oči a snaží se rozmotat zpátky do kruhu \textbf{aniž by se rozpojily ruce}{}
\end{enumerate}
 
Ne všechny možnosti propojení se dají rozmotat. Pokud se stane, že skupina utvoří oddělené n-tice, je cílem rozmotat je odděleně. 
 
 
 
 
 
\needspace{5cm} \section{Heja hají} \label{heja hají} Skupina stojí v kruhu a posílá se signál. 
 
Možné signály zhruba seřazené dle četnosti výskytu a znalostí mezi improvizátory: 
\begin{itemize}
\item \textbf{Heja}{} - pokračuje ve směru
\item \textbf{Hold down}{} - obrací směr
\item \textbf{Hají}{} - jednoho přeskočí
\item \textbf{Hold hají}{} - otočí a přeskočí
\item \textbf{Beng}{} - na kohokoli ve kruhu
\item \textbf{Pinguin party}{} - hráči si vyměňují místa
\item \textbf{Blechy}{} -  dtto s chytáním blech
\item \textbf{Klarinet}{} -  přepíná jen na gesta, bez zvuků
\item \textbf{Hoboj}{} - přepíná zpět
\end{itemize}
 
Dbejme na čistotu gest, rytmičnost a pozornost. 
 
\subsection{Varianty} \begin{itemize}
\item Několik signálů v kruhu
\item Místo zvuků říkáme \odkaz{asociace}{asociace}
\end{itemize}
 
\subsection{Viz také} \begin{itemize}
\item  \odkaz{Posílání signálu po kruhu}{posílání signálu po kruhu}
\end{itemize}
 
 
 
\needspace{5cm} \section{Hot spot} \label{hot spot} Toto cvičení je zaměřeno na nalazení hráčů a vzájemnou podporu. Hráči stojí v kruhu. Jeden z nich jde do prostřed a začne zpívat jakoukoliv známou písničku, ostatní se k němu přidávají. A i když neznají slova, snaží se ho jakkoliv podpořit. Ve chvíli kdy někdo z dalších hráčů má asociaci na novou písničku, vstupuje do kruhu, jemně se dotkne hráče uprostřed, a vymění si s ním místo. Střídání probíhá velice rychle a často zaznívají jen první slova z písničky. 
 
 
 
\needspace{5cm} \section{Kdo ukradl sušenky a ze stolu je vzal} \label{kdo ukradl sušenky a ze stolu je vzal} Skupina stojí v kruhu. Přípravou na cvičení je opakování věty "Kdo ukradl sušenky a ze stolu je vzal"{} (dále jen věta), dokud všichni nedrží rytmus a neznají to nazpaměť. Pro toto cvičení je ideální čtyřčtvrťový rytmus. 
 
Určí se, kdo začíná. Poté celá skupina jednou zopakuje \textif{větu}{}. Začínají jedinec si náhodně někoho vybere a použije v rytmu jeho jméno. Vypadá to následovně: 
 
\begin{enumerate}
\item  Skupina: Kdo ukradl sušenky a ze stolu je vzal
\item  A: Kdo ukradl sušenky a ze stolu je vzal
\item  A: Martin
\item  B: Kdo, já?
\item  A: Jó ty!
\item  B: Já ne!
\item  A: Tak kdo?
\end{enumerate}
 
Poté skupina opět zopakuje větu a pokračuje hráč, kterého byl naposledy označen za "zloděje sušenek". 
 
Cvičení se opakuje, dokud cvičící uzná za vhodné. 
 
\subsection{ Pomocné techniky } \begin{itemize}
\item  Dupání, klepání, luskání do rytmu
\end{itemize}
 
 
 
 
 
\needspace{5cm} \section{Kung-fu} \label{kung-fu} Hráči stojí ve dvojici naproti sobě, na větší vzdálenost než je jejich natažená paže. Hráči se střídají v útocích, švihnutím paže. Ve chvíli, kdy útočí hráč A, hráč B uhýbá. V obou případech hráči doprovází svůj pohyb slabikou "ha". 
 
\subsection{ Možnosti útoku a úhybu } \begin{itemize}
\item  Švihnutím paže nad hlavu protihráče, druhý hráč uhýbá skrčením se.
\item  Švihnutím paže přímo před sebe, druhý hráč uhýbá uhnutím v boku.
\item  Švihnutím paže k nohám protihráče, druhý hráč uhýbá vyskočením do vzduchu.
\end{itemize}
 
 
 
\needspace{5cm} \section{Most} \label{most} Kontaktní rozcvička pro větší skupinu, ve které jeden předem určený člověk přechází po ostatních jako by byli most z bodu A do bodu B. Most se může jakýmkoliv způsobem transformovat a ulehčovat nebo i zesložiťovat cestu. 
 
Během cvičení se nemluví. Je vhodné účastníky cvičení upozornit na mírnou fyzickou zátěž. 
 
 
 
 
 
\needspace{5cm} \section{Ninjové} \label{ninjové} Hráči stojí v kruhu, v určeném pořadí (po směru hodinových ručiček) hráč, který je na tahu, útočí jedním plynulým pohybem na jiného hráče, cílem je zasáhnout dlaní či prsty předloktí druhého hráče a tím mu danou končetinu vyřadit. Hráč, na kterého je útočeno, smí jedním plynulým pohybem zkusit uniknout útoku. 
Hráč, který nemá ani jednu ruku, vypadává, stejně tak jako hráč, který útočí, když není na řadě. 
 
\subsection{Varianta (pokračování) } Ve chvíli kdy hráči ovládají hru v kruhu, můžeme přejít na vyhazování v reálném čase (bez určených tahů) s tím, že se člověk do hry vrací, když je vyhozen ten, kdo mu sebral poslední končetinu.  
 
 
\needspace{5cm} \section{Otázky (cvičení)} \label{otázky (cvičení)} Cílem této [[:Kategorie:Warm-up|warm-up]] hry je, aby účastníci na sebe prozradili něco nového a zároveň se nebáli akce a vystupování z řady. Moderátor předloží nahlas uzavřenou otázku, tedy takovou, aby na ní bylo možné odpovědět ano-ne. Ve větším prostoru mohou účastníci stát na jedné straně a v případě jejich kladné odpovědi, přejdou na druhou stranu místnosti. Moderátor následně menší skupinu vyzve také k přechodu, aby získal znovu 1 skupinu a pokládá další otázku. V menším prostoru se mohou účastníci přesunovat z židle na židli. Důležité je, aby bylo dostatečně viditelné, kdo odpovídá ano a akce byla výrazná. 
 
\subsubsection{Příprava hry:} Důležitá je příprava otázek, které mohou být od úplně jednoduchých po výrazně osobní. Správně vymyšlené otázky budou v účastnících vyvolávat také smích a tím dojde k uvolnění atmosféry. Tedy přesně toho, co od hry potřebujeme. 
 
\subsubsection{Délka:} 5 - 10 minut 
 
\subsubsection{Tipy otázek:} Máte radši kočky než psy? 
Věříte na horoskop? 
Zamilovali jste se někdy do cizince? 
 
 
 
 
\needspace{5cm} \section{Počítání do 20} \label{počítání do 20} Účastníci cvičení stojí v kruhu zády k sobě. Mají za úkol jako skupina napočítat do dvaceti. Každý kdo řekne číslo, musí v ten samý okamžik udělat krok vpřed. Pokud 
promluví víc než jeden naráz, všichni se vrací do výchozí pozice a začínají znovu. 
 
 
 
 
\needspace{5cm} \section{Pomluva} \label{pomluva} \textbf{Pomluva}{} též známá jako fáma. Hráči v tomto cvičení imitují šíření pomluvy, tak jak to známe z běžného šíření "zaručených zpráv". Cvičení se podobá tiché poště, ale hráči si předávají informaci nahlas. Hráči stojí vedle sebe a prvnímu z nich je zadána věta s jednoduchou informací. Hráč předává tuto informaci druhému hráči s drobnou obměnou. Druhý hráč předává informaci třetímu zase trošku obměněnou, atd. až k poslednímu hráči. Je nutné, aby hráči reagovali jen na informaci, kterou obdrželi a ne na informaci, kterou obdrželi hráči před ním. Také je dobré, aby už v polovině skupiny nehrozila globální katastrofa či zničení vesmíru, života a vůbec. 
 
 
\subsection{Příklad}  
Základní informace: V samoobsluze zdražili rohlíky. 
 
1. hráč: V potravinách zase podražilo pečivo. 
 
2. hráč: Letos se urodilo málo obilí, a tak je vše dražší. 
 
3. hráč: Zemědělci si zase stěžují na malou úrodu. 
 
4. hráč: Zemědělci chtějí dotace od EU. 
 
atd. 
 
 
 
\needspace{5cm} \section{Popcorn} \label{popcorn} Cvičení pro alespoň pět hráčů. Všichni si stoupnou do kruhu a na místě vyskakují. Během výskoku hráč musí tlesknout - tedy dřív než dopadne. Pokud se dva hráči stlesknou, tedy tlesknou oba současně, musí si na krátký interval, třeba 5 vteřin, dřepnout. Zvuk tohoto cvičení pro vnějšího pozorovatele zní, jako když v mikrovlnné troubě praská popcorn, proto ten název. 
 
 
 
 
 
\needspace{5cm} \section{Posílání signálu po kruhu} \label{posílání signálu po kruhu} Cílem cvičení je nechat hráče se na sebe naladit, uvědomovat si ostatní a přijmout vlastní zodpovědnost v rámci skupiny. Zadání pro skupinu zní, aby posílání signálu bylo držené v jednom rytmu, a vždy, když signál předáváme dál, měli s příjemcem oční kontakt a věděli, že příjemce je připravený signál přijmout. 
 
\subsection{  Verze } \begin{itemize}
\item  V1: V lehké verzi hráči posílají signál po kruhu
\item  V2: Hráči posílají signál libovolně na přeskáčku
\item  V3: Hráči posílají signál na přeskáčku a mění varianty signálů
\item  V4: Více signálů najednou
\end{itemize}
 
\subsection{ Varianty signálů } \begin{itemize}
\item  formou tleskání
\item  formou přejetí dlaní o dlaní
\item  házení míčku
\item  mrknutí
\item  říkání jmen apod
\item  \odkaz{Heja hají}{heja hají}
\end{itemize}
 
\subsection{ Těžkosti } \begin{itemize}
\item  V případě, když skupina neudrží rytmus posílání signálu, je vhodné nastavit ve skupině pomalé tempo a postupně zrychlovat.
\item  Pokud nejsou hráči při předávání signálu v očním kontaktu, je dobré zkusit nejdříve variantu s mrknutím, které nutí hráče být pozorní.
\end{itemize}
 
\subsection{Varianty} \begin{itemize}
\item \odkaz{Heja hají}{heja hají}
\item \odkaz{Koncepty}{koncepty}
\end{itemize}
 
 
 
\needspace{5cm} \section{Příběh po fázích} \label{příběh po fázích} {{Kategorie 
|hraci=libovolný 
|cas=volný}} 
 
Napojení improvizátoři vypráví příběh po rozdělených fázích. 
 
\subsection{ Průběh }  
Herci stojí nebo sedí vedle sebe v kruhu. Každý má za úkol odvyprávět jednu izolovanou fázi příběhu.  
 
\subsection{ Běžné chyby }  
\begin{itemize}
\item  Zapomínání na věci z expozice
\end{itemize}
 
 
 
 
 
\needspace{5cm} \section{Ready Go} \label{ready go} Cvičení \textbf{Ready Go}{} je dobré pro procvičování poslouchání svých spoluhráčů a současně na pohyb. 
 
\subsection{ Průběh cvičení } Hráči stojí v řadě vedle sebe. První hráč začne vyprávět nějaký příběh na téma, nebo může třeba vyprávět svůj denní zážitek a snaží se u toho výrazně gestikulovat a pohybovat. Ostatní hráči poslouchají a sledují co dělá. Hlavně druhý hráč, který musí být připraven. Třetí hráč totiž může kdykoli během vyprávění 1. hráče říci \textbf{Ready!}{}, což je povel pro 2. hráče a ten musí co nejvěrněji kopírovat pohyby prvního hráče. Takto se chvíli pokračuje, aby se druhý hráč naladil na prvního, současně musí dávat pozor na to, co první hráč vypráví, protože třetí hráč může opět kdykoli říci \textbf{Go!}{} V tuto chvíli přebírá vyprávění druhý hráč a současně své vyprávění doprovází svými pohyby. První hráč odchází a může zbytek cvičení sledovat. Roli třetího hráče, který dával povely se posouvá na dalšího v řadě. Takto se vystřídají všichni hráči. Poslední hráč příběh uzavírá a povely mu udílí první hráč nebo trenér. 
 
Příběh by měl být kontinuální a zachovává se pohlaví prvního hráče. Pokud tedy začala vyprávět žena i mužští hráči pokračují ve vyprávění v ženském rodě. 
 
 
\needspace{5cm} \section{Sloupce} \label{sloupce} Cvičení ve kterém si herci pomocí minimalizace výrazových prostředků procvičí vnímání ostatních herců na jevišti, vnímání sebe, vnímání diváků a skupinovou responsivitu. Vhodný počet herců je 4 - 7. 
 
\subsection{ Průběh }  
Herci se postaví do řady tak aby mezi sebou měli místo alespoň na délku ruky. Prostor pro cvičení není omezen. Trenér zadá výrazové prostředky, které herci smějí používat. Pokud herec použije výrazový prostředek jiný, než bylo zadáno je dobré ho upozornit, aby se držel zadání. 
 
Cvičení nemá specificky daný čas a skupina si ho ukončuje sama. 
 
\subsubsection{ Výrazové prostředky }  
\begin{itemize}
\item  Krok dopředu
\item  Otočení o 180°
\item  Dřep dolů
\item  Dřep nahoru
\item  Výskok
\end{itemize}
 
\subsection{ Varianty } V zadání může být, že herci mají vyprávět příběh. Cvičení potom skončí s příběhem. 
 
 
 
 
 
\needspace{5cm} \section{Stleskávání} \label{stleskávání} {{Kategorie 
|hraci=libovolný,minimálně 5 
|cas=volný}} 
 
Účastníci cvičení stojí v kruhu a posílají si tlesknutí. Cílem je tlesknout ve stejný čas a naladit se na další práci 
 
1. Navázat oční kontakt 
2. Tlesknout společně 
 
\subsection{ Tipy }  
\begin{itemize}
\item  Pomáhá udržovat rytmus tleskání
\end{itemize}
 
\subsection{ Varianty }  
\begin{itemize}
\item  Zrychlování
\item  Vyslání víc signálů na začátku
\end{itemize}
 
 
 
 
\needspace{5cm} \section{Tvary} \label{tvary} Lehce kontaktní rozcvička se zavřenýma očima a bez mluvení. Cílem je vytvořit konkrétní tvary ze svých těl, například kruh, čtverec, trojúhelník nebo klidně něco složitějšího, jako třeba jeřáb. Může se hrát ve stoje i na zemi. 
 
Pomáhá skupině na uklidnění se, soustředění a práci s chybou. Musí být vedena zvenčí. 
 
 
 
 
 
 
\needspace{5cm} \section{Zabírání kopce} \label{zabírání kopce} Cvičení na rozhýbání skupiny a rozproudění myšlenek. V cvičení půjde hodně o přebíjení, ale také o spolupráci. 
 
\subsection{Průběh} Stojí se v kruhu. První člověk začíná se větou typu "Vylezu na kopec a ten je můj". Následuje slovní souboj o kopec, kdy každý mírně upraví co se stane například "Vylezu na kopec, shodím tě dolů a kopec je můj". Skupina by měla dělat co nejmenší změny, aby bylo kam gradovat. Může se jet po kruhu a nebo náhodně podle toho kdo má nápad.  
 
\subsection{Varianta beze slov} Hra je čistě pantomimická. Důraz na co nejmenší změny. 
 
 
 
 
 
\needspace{5cm} \section{Zaplňování prostoru} \label{zaplňování prostoru} Pramáti všech improvizačních cvičení, využívaná v nespočetném množství variant. \odkaz{Hráči}{hráč} se pohybují po prostoru a snaží se jej rovnoměrně zaplňovat. 
 
\subsection{Varianty} \begin{itemize}
\item Rychlosti chůze  1-5 , úkoly - výskok, dřep
\item Zdravení lidí pohledem
\item Chůze v emoci/charakteru, potkávání ostatních a sehrávání krátkých scén
\item Čtení dopisu/ básně ve více lidech
\item Počítání do 20
\item \odkaz{Štronzo}{štronzo}, zavřít oči a pak ukázat na nějakého hráče (jména, popis)
\end{itemize}
 
 
 
\needspace{5cm} \section{Zip, zep, zop} \label{zip, zep, zop} Hráči stojí v kruhu a posilají si mezi sebou jeden signal ukázáním na někoho prstem s hláškou zip, zep nebo zop v tomto pořadí. 
 
Kolikrát padne slovo “zip”, tolikrát se musejí opakovat i další dvě hlášky. 
 
 
\textbf{Příklad:}{}

 
''Hráč A vysílá “zip”, pokud hráč B řekne “zep”, hráč C musí říct “zop”. 
Pokud hráč B zopakuje “zip”, musí i další hlášky padnout 2x a tak dále.''  
 
 
Další “zip” už restartuje počet a opět záleží kolikrát “zip” od hráčů padne. 
 
Ve chvíli, kdy hráč udělá chybu nebo zaváha, říká místo “zip, zep a zop”, “bip, bep a bop” a s každou další chybou dostává nové písmeno na začátek, třeba až do “g”. Pak může následovat ještě obměna a může říkat například místo “gop” “gong” apod. 
 
Hra zpravidla pokračuje do té doby dokud i poslední hráč neudělá chybu. 
 
\subsubsection{Návrh otázek k reflexi:} \begin{itemize}
\item  Co vám na hře přišlo složité?
\item  Jakou jste měli strategii, abyste hru vyhráli?
\item  Co vám pomáhalo neudělat chybu?
\end{itemize}
 
 
 
 
 
 
 
 
 
\needspace{5cm} \section{Živé obrazy} \label{živé obrazy} \katabox{Prostředí 
}{neomezený 
}{neomezený, obvykle 2-3 minuty} 
 
Je ve své prapodstatě \odkaz{štronzo}{štronzo} několika improvizátorů které vytváří situaci, či prostředí. V živém obraze není vhodné se hýbat, ani vydávat zvuky. 
 
{{Todo| 
\begin{itemize}
\item  Doplnit průběh
\item  Doplnit video
\end{itemize}
}} 
 
 
 
 

\end{document}

\chapter{Fauly}\label{fauly}

\label{:kategorie:fauly}
\documentclass[main.tex]{subfiles}\begin{document}
\needspace{5cm} \label{faul} Podobně jako v mnoha sportech, i na improvizačních \odkaz{zápasech}{zápas} může \odkaz{rozhodčí}{rozhodčí} písknout \textbf{faul}{}. Tím označuje chybu \odkaz{hráče}{hráč} nebo hráčů a může se rozhodnout udělit týmům \odkaz{trestné body}{trestný bod}. 
 
Ohlášení faulu provádí \odkaz{rozhodčí}{rozhodčí} kombinací písknutí a pohybu specifického pro každý faul. 
 
Konkrétní fauly najdete v \odkaz{kategorii Fauly}{:kategorie:fauly}. 
 
 

\end{document} 
%\documentclass[main.tex]{subfiles}\begin{document}
\needspace{5cm} \section{Bránění ve hře} \label{bránění ve hře} \faulbox{\obrazekmaly{Braneni.png}}{Rozhodčí zkříží ruce ve výši prsou}{1 } 
 
Bránění ve hře též blokování hry je faul pískán v případě, pokud hráč neguje nabídku svého spoluhráče případně pokud jej výrazně manipuluje. Dalším případem je, pokud spoluhráči zamezí možnost se projevovat nebo pohybovat. 
 
 
\needspace{5cm} \section{Hrubá chyba} \label{hrubá chyba} \faulbox{\obrazekmaly{Hruba-chyba.png}}{pohyb zalomenou paží nahoru a dolů}{2} 
Pokud hráč opakuje již faulované chování, může rozhodčí pískat hrubou chybu. V případě dále stupňovaného opakovaného chybování může rozhodčí hráče vyloučit z následující kategorie/kategorií. 
 
\needspace{5cm} \section{Klišé} \label{klišé} \faulbox{\obrazekmaly{Klise.png}}{Rozhodčí klepne dlaní o kotník}{1 } 
 
Pokud je v improvizaci použito jméno širokou veřejností známé značky, osobnosti, knižní nebo filmové postavy, citace z filmu apod., měl by být písknut \odkaz{faul}{faul} \textbf{Klišé}{}. Obecně se do tohoto výčtu počítají jakékoliv reference na známé věci. V zápasových kategorií je používání klišé zakázáno a pískáno jako faul, protože ve hře odkazuje na věci mimo hru a stává se vtipem, který hru rozbíjí, místo, aby ji rozvíjel. 
  
 
Výjimkou z tohoto pravidla jsou názvy \odkaz{improligových}{čili} týmů a názvy zemí. 
 
 
\needspace{5cm} \section{Mimo rámec} \label{mimo rámec} \faulbox{\obrazekmaly{Mimo-ramec.png}}{Rozhodčí do vzduchu kreslí rámeček}{1 } 
 
\odkaz{Hráč}{hráč} porušil stanovená pravidla \odkaz{kategorie}{kategorie} a nebo došlo k úplnému opuštění tématu improvizace. Tento \odkaz{faul}{faul} se nejčastěji objevuje v kategorii \odkaz{Smrt v 1 minutě}{smrt v 1 minutě}, když jedna postava nezemře po 60 vteřinách, případně ve \odkaz{Stíhačce}{stíhačka}, začnou li hráči hrát ještě před písknutím. 
 
 
\needspace{5cm} \section{Nátlak (faul)} \label{nátlak (faul)} \faulbox{\obrazekmaly{Natlak.png}}{Rozhodčí bouchá pěstí do dlaně}{1 } 
\odkaz{Hráč}{hráč} blokuje ostatní spoluhráče tím, že se snaží protlačit si vlastní způsob rozvoje příběhu a nebo je omezuje akcemi svojí postavy např. rozkazováním a nebo odmítáním. 
 
 
\needspace{5cm} \section{Nepovolená rekvizita} \label{nepovolená rekvizita} \faulbox{\obrazekmaly{Nepovolena-rekvizita.png}}{Rozhodčí se tahá za cíp dresu}{1 } 
 
Pokud není dohodnuto jinak a nebo nemá \odkaz{kategorie}{kategorie}, ve které se scénka hraje, vlastní pravidla o \odkaz{rekvizitách}{rekvizita}, jsou na \odkaz{zápasech}{zápas} povoleny pouze následující věci pro použití v improvizaci: 
 
\begin{itemize}
\item  vlastní tělo
\item  tričko nebo týmový dres
\item  ostatní \odkaz{hráči}{hráč}, kteří nejsou v \odkaz{postavě}{postava}
\end{itemize}
 
V momentě, kdy \odkaz{hráč}{hráč} použije jakýkoliv jiný objekt jako \odkaz{rekvizitu}{rekvizita}, je na místě tento \odkaz{faul}{faul} 
 
 
\needspace{5cm} \section{Nepovolený postup} \label{nepovolený postup} \faulbox{\obrazekmaly{Nepovoleny-postup.png}}{poklepávání na paži}{1 } 
 
 
Hráč či tým nedodržují pravidla kategorie či ducha improvizace, a nelze to postihnout jinak. 
\needspace{5cm} \section{Nepozornost} \label{nepozornost} \faulbox{\obrazekmaly{Nepozornost.png}}{ruka držící druhou v zápěstí}{1} 
 
O nepozornost se jedná, pokud hráč ignoruje vstup nových postav na pódium, 
nehraje s informacemi, které padly (včetně jmen postav), nebo pokud na scéně  
vládne obecné neporozumění. 
 
 
\needspace{5cm} \section{Nepřehlednost hry} \label{nepřehlednost hry} \faulbox{\obrazekmaly{Neprehlednost-hry.png}}{Rozhodčí dělá rukama mlýnek}{1} 
 
Pokud situace v rozehrané improvizaci není jasná, je na místě \odkaz{faul}{faul} \textbf{Nepřehlednost hry}{}. Zpravidla takové situace vznikají tím, že \odkaz{hráči}{hráč} vytvoří zbytečně komplikovaný příběh nebo prostředí a nebo dojde k vzájemnému nepochopení. 
 
 
\needspace{5cm} \section{Nerespektování postavy} \label{nerespektování postavy} \faulbox{\obrazekmaly{Nerespektovani-postavy.png}}{Rozevřenou rukou (dlaní) si zakrejeme obličej}{1} 
 
Faul nerespektování postavy píská \odkaz{rozhodčí}{rozhodčí}, pokud některý hráč  
jedná s jinou postavou v rozporu s tím, co o jeho postavě víme. 
Nejčastěji se jedná o zapomenutí jména či rodinného stavu. 
Dále v případě, že muž hraje ženu a naopak. 
\needspace{5cm} \section{Nevhodné chování} \label{nevhodné chování} \faulbox{\obrazekmaly{Nevhodne-chovani.png}}{Ruce v bok}{2 } 
 
 
Pod nevhodné chování se dají zahrnout rozličné chyby proti dobrému duchu 
improvizace, zejména pak používání sprostých a vulgárních termínů, 
nadbytečné urážení týmu protivníka, neustálé odmlouvání hlavnímu rozhodčímu,  
apod. Dále sem patří i opakované stupňované \odkaz{šaškování}{šaškování}. 
\needspace{5cm} \section{Plynulost hry} \label{plynulost hry} \faulbox{\obrazekmaly{Plynulost-hry.png}}{Rozhodčí krouží prstem nad hlavou}{1 } 
 
Pokud scénka vázne, je na místě písknout faul \textbf{plynulost hry}{}. Taková situace může nastat, pokud se příběh nikam dlouho neposune, což může být zapříčiněno více faktory. 
 
\begin{itemize}
\item  \odkaz{hráči}{hráč} nenastoupí včas do \odkaz{kategorie}{kategorie}
\item  postavy se pouze dohadují
\item  není definován \odkaz{improvizační trojzubec}{improvizační trojzubec}
\end{itemize}
 
 
\needspace{5cm} \section{Přečíslení} \label{přečíslení} \faulbox{\obrazekmaly{Precisleni.png}}{Rozhodčí se plácá dlaní do hlavy}{1 } 
 
Pokud vznikne situace, kdy je na scéně o dva a více hráčů z jednoho týmu víc než z druhého jedná se o \textbf{přečíslení}{}. Jeden tým totiž ubírá prostor tomu druhému a vytváří tak nerovné podmínky. Tento faul se nepíská v \odkaz{kategoriích porovnávacích}{:kategorie:porovnávací kategorie} 
 
 
\needspace{5cm} \section{Šaškování} \label{šaškování} \faulbox{\obrazekmaly{Saskovani.png}}{Rozhodčí ukazuje dlouhý nos}{1 } 
 
\odkaz{Hráč}{hráč} se snaží být vtipnější na úkor ostatních a sklouzává k trapným klišé nebo gagům, které neposouvají děj, případně nabourává děj jednorázovými fórky. 
 
 
\needspace{5cm} \section{Vypadnutí z role} \label{vypadnutí z role} \faulbox{\obrazekmaly{Vypadnuti-z-role.png}}{Ruka zaťatá s pěstí, zapumpujeme dvakrát nahoru a dolu   }{1} 
 
Faul vypadnutí z role píská \odkaz{rozhodčí}{rozhodčí} zejména ve chvíli, kdy se hráč odbourá - rozesměje,  
nebo pokud změní svoji postavu příliš razantním a pro děj narušujícím způsobem. 
 
\needspace{5cm} \section{Výstraha} \label{výstraha} \faulbox{\obrazekmaly{Vystraha.png}}{poklepávající ruce v bok }{2} 
 
Hráč může být v zápase varován před opakováním nějaké chyby.  
Při dalším opakování se pak může pískat výstraha následovaná např. vyloučením z další kategorie. 
 
 
 
 

\end{document} 

\chapter{Terminologie}\label{terminologie}
\label{:kategorie:terminologie}
%\needspace{5cm} \section{ABAB} \label{abab} Vzorec básně, ve kterém se rýmují sudé verše se sudými a liché verše s lichými. 
 
Příklad takové básně: 
	 
\begin{quote} 
Někdo se bojí vlastní ženy (A)

 
A někdo třeba upíra (B)

 
Někdo je v přítmí vyděšený (A)

 
Ve tmě už téměř umírá (B) 
 
Jan Burian, Zubař 
\end{quote} 
 
Tento vzorec se používá hlavně v kategoriích \odkaz{Krátká báseň}{krátká báseň} a \odkaz{Opilecká píseň}{opilecká píseň}. 
 
Dál se často používají střídání AABB a ABBA. 
 
 
 
\needspace{5cm} \section{Askfor} \label{askfor} \textbf{Askfor}{} - v překladu např. \textbf{zeptání se}{}. \odkaz{MC}{mc} si vyžádá od publika téma, na které se bude hrát. 
 
Může se jednat o: 
\begin{itemize}
\item \textbf{Košíček s tématy}{}, do kterého publikum před \odkaz{zápasem}{zápas} vhazuje napsaná dvojslovná témata
\item \textbf{Reálné předměty}{}, které diváci zapůjčí na \odkaz{Prskavku ala šanson}{prskavka ala šanson}
\item \textbf{Prostředí}{} udávané např na \odkaz{Toaster}{toaster}, \odkaz{Smrt v 1 minutě}{smrt v 1 minutě} či \odkaz{Poslední větu}{poslední věta}
\item \textbf{Emoce}{} pro \odkaz{Schizofrenii}{schizofrenie}
\item \textbf{Problém}{} pro \odkaz{Barman song}{barman song}
\item \textbf{Jazyk}{} pro \odkaz{šumlovanou}{šumlovaná}, \odkaz{šumlovanou se simultánním překladem}{šumlovaná se simultánním překladem}, \odkaz{duál}{duál} a \odkaz{šumlovaný sbor}{šumlovaný sbor}. Víz také \odkaz{Gibberish}{gibberish}.
\end{itemize}
 
Dobrovolník např. do \odkaz{Love songu}{love song} je svým způsobem také Askfor 
 
 
\needspace{5cm} \section{Asociace (terminologie)} \label{asociace (terminologie)} Asociace je termín, kterým označujeme okamžitý nápad, který se nám vybaví bez přemýšlení. Asociace tak může s prvním zadáním zdánlivě nesouviset např. strom -> ovoce -> počítač (Apple). Improvizátoři trénují svou mysl a paměť, aby byli schopní vždy okamžitě reagovat i když může být zadáním slovo, jehož význam neznají. 
 
 
\needspace{5cm} \section{Báseň} \label{báseň} Báseň v širším slova významu představuje literární dílo tvořené ve verších. 
V improvizaci se s tímto pojmem diváci mohou potkat jako s kategorií \odkaz{Poetická}{poetická} nebo rozcvičkovým cvičením \odkaz{Báseň}{báseň (cvičení)}. 
 
\subsection{Viz. také} \begin{itemize}
\item \odkaz{ABAB}{abab}
\item \odkaz{Divadlo poezie}{divadlo poezie}
\end{itemize}
 
 
\needspace{5cm} \section{Forma} \label{forma} \textbf{Forma}{} je zkrácené označení pro formát představení. Určuje postavy a způsob, kterým budou vystupovat po dobu trvání celého představení. Seznam forem naleznete v \odkaz{kategorii Formy}{:kategorie:formy}. 
 
 
\needspace{5cm} \section{Gibberish} \label{gibberish} \textbf{Giberish}{} je označení pro vymyšlený jazyk. Setkat se můžeme také s českým označením \textbf{šumlovaná}{} nebo také \textbf{svojština}{}. Improvizátor vytváří vlastní řeč, hledá pro něj pravidla. Inspirace pro znění jazyka by mělo vycházet z kontextu situace ve kterém se hráč nachází, jako postavu hraje, jaké jsou jeho emoce apod. 
 
Hráči předstírájí  znalost neznámého cizího jazyka, někdy \odkaz{zadaného publikem}{askfor}. 
 
Giberish může být vhodný nástroj pro rýmování, kdy hráč není zatěžován významem slov, a umožňuje mu soustředit se na rytmiku a výrazovost projevu. 
 
\subsection{Viz také} \begin{itemize}
\item \odkaz{kvaziznaková řeč}{kvaziznaková řeč}
\end{itemize}
 
 
\needspace{5cm} \section{Hlasovací kartička} \label{hlasovací kartička}  
 
Pomůcka, na kterou by se nemělo zapomenout při organizování klasického \odkaz{zápasu}{zápas}. Divák při vstupu do sálu dostane od pomocného rozhodčího či jiné pověřené osoby dvě hlasovací kartičky různých barev nebo jednu s kartičku s každou stranou jiné barvy. Hlasovací kartička je zhruba velikosti A6 většinou z papírového kartonu či zalaminovaná. Barva závisí na hrajících týmech, každý \odkaz{tým}{tým} má svou specifickou barvu či hraje v dresech dané barvy. Divák prostřednictvím hlasovací kartičky dává najevo, kterému \odkaz{týmu}{tým} by přidělil bod v průběhu \odkaz{hlasování}{hlasování} o proběhlé improvizaci. 
 
 
\needspace{5cm} \section{Hlasování} \label{hlasování} Na \odkaz{zápasech}{zápas} se využívá \textbf{hlasování}{}, aby se určilo, který \odkaz{tým}{tým} dostane bod za právě odehranou \odkaz{kategorii}{kategorie}. Organizátor a nebo \odkaz{MC}{mc} určují způsob, jakým mohou diváci hlasovat. 
 
Hlasování se používá jako varianta \odkaz{náhodného rozhodování}{náhodné rozhodování}. 
 
\subsection{ Obvyklé způsoby hlasování } \begin{itemize}
\item  Pomocí \odkaz{hlasovacích kartiček}{hlasovací kartička}
\item  \odkaz{Vteřinovým potleskem}{vteřinový potlesk}
\end{itemize}
 
 
\needspace{5cm} \section{Impro} \label{impro} Zkratka používaná pro divadelní improvizaci. Velice často se také spojuje s jinými slovy, jako \odkaz{improtým}{tým}, \odkaz{improkategorie}{kategorie}, \odkaz{Improtřesk}{improtřesk} 
 
 
\needspace{5cm} \section{Improshow} \label{improshow} Improvizované představení, které je podobné \odkaz{zápasu v improvizaci}{zápas}. Narozdíl od zápasu ale hraje jenom jeden tým a představení nemá \odkaz{rozhodčího}{rozhodčí}, ani \odkaz{pomocňáky}{pomocný rozhodčí}. 
 
\subsection{ Průběh } Improshow není standardizovaný formát představení, na kterém se hrají \odkaz{:Kategorie:Krátké formy}{:kategorie:krátké formy} a cokoliv se může změnit. Obvykle se hrají \odkaz{kategorie}{kategorie} v nějakém pořadí a celým večerem provází \odkaz{konferenciér}{konferenciér}. 
 
 
 
 
 
\needspace{5cm} \section{Improvizační trojzubec} \label{improvizační trojzubec} Pojem označující seřazenou sadu tří základních otázek, na které by si měla umět každá postava v improvizaci odpovědět. Pokud to nedokáže, měla by ve scénce vytvořit věci, které k tomu chybí. 
 
\subsection{ Kdo jste? } Byly všechny postavy dobře popsány? Ví všichni o všech dost na to, aby z toho mohli vycházet? 
 
\subsection{ Kde jste? } Bylo definováno prostředí? Dá se s ním pracovat? 
 
\subsection{ Co tam děláte? } Proč jsme v této situaci? Proč moje postava dělá to, co dělá? Zachovala by se takhle pokaždé? 
 
 
 
\needspace{5cm} \section{Kvaziznaková řeč} \label{kvaziznaková řeč} Způsob komunikace založený na pohybu celým tělem. Každý projev v této řeči je originální, protože pohyby obvykle vycházejí z první asociace na významovou část překládané věty. Tlumočník se obvykle snaží v reálném čase překládat řeč nebo rozhovor. Při překladu z verbální komunikace se nepřekládá každé slovo, ale pouze významová část. 
 
Při překladu se vyjadřuje emoce a její intenzita výrazem v obličeji a rychlostí pohybu. 
 
Tato technika se obvykle používá v kategorii \odkaz{Překlad do znakové řeči}{překlad do znakové řeči} a \odkaz{Sportovní komentátor}{sportovní komentátor}. 
 
 
\needspace{5cm} \section{Motivovaný odchod} \label{motivovaný odchod} Odchod zdůvodněný logikou scény. Např. při scéně s \odkaz{výměnou žárovky}{výměna žárovky} to může být pro vypnutí proudu. Je potřebný zejména pro \odkaz{pyramidu}{pyramida} a \odkaz{Tam a zpět}{tam a zpět}. 
 
 
\needspace{5cm} \section{Náhodné rozhodování} \label{náhodné rozhodování} Na \odkaz{zápasech}{zápas} je potřeba občas náhodně rozhodnout, který tým dostane výhodu. Převážně se jedná o rozhodování, který tým bude začínat s pokřiky a nebo v porovnávacích kategoriích. 
 
\subsection{ Puk } V ideálním případě má rozhodčí k dispozici puk. Kapitání týmů si vyberou, která strana podle nich padne. Výherce si vybírá, jestli chce začít. Občas se jako náhražka používá mince. Hod probíhá pod důsledným dozorem \odkaz{pomocných rozhodčích}{pomocný rozhodčí}. 
 
\subsection{ Náhodný generátor čísel } \odkaz{Rozhodčí}{rozhodčí} se zeptá týmů kdo chce začít a během toho v hlavě vygeneruje náhodné číslo. Tým, jehož hlasitost v decibelech při křičení "Ano"{} je blíže začíná. 
 
\subsection{ Diváci } Vždycky je možnost zeptat se diváků, pro kterou možnost budou \odkaz{hlasovat}{hlasování}. 
 
 
 
\needspace{5cm} \section{Pantomima} \label{pantomima} Pantomima je divadelní způsob vyjádření pomocí mimiky a gestikulace s prvky tance bez použití hlasu.  
 
 
 
\needspace{5cm} \section{Papuče} \label{papuče} Molitanová papuče či lehký míček slouží divákům k vyjádření nesouhlasu s verdiktem hlavního \odkaz{rozhodčího}{rozhodčí}. Papuče dostávají diváci při vstupu do sálu, vždy jeden kus na osobu. Během začátku zápasu instruuje \odkaz{konferenciér}{konferenciér} diváky, jak mohou papuči využít. Vyjádření nesouhlasu zpravidla bývá hlavně při udílení trestných bodů po proběhlé improvizaci, kdy hlavní rozhodčí upozorňuje na fauly, kterých se hráči v průběhu kategorie dopustili. Pokud se divákům zdá faul nespravedlivý nebo hráčů argument dobrý, ale rozhodčím neuznaný, mají možnost hodit papuči či míček na jeviště a tím vyjádřit svůj nesouhlas. Hlavní rozhodčí může vzít tento nesouhlas v potaz a trestný bod neudělit. 
 
\odkaz{Pomocní rozhodčí}{pomocný rozhodčí} obvykle papuče/míčky vrací zpátky divákům jako  přísun další munice.  
Zřídkakdy diváci použijí papuče k vyjádření nesouhlasu s chováním některého hráče. 
 
 
\needspace{5cm} \section{Porovnávací kategorie} \label{porovnávací kategorie} Při \odkaz{kategoriích}{kategorie} v tomto seznamu nehrají týmy společně, ale zvlášť. 
 
 
\label{:kategorie:porovnávací kategorie} 
\begin{multicols}{2}\begin{itemize} 
\item  \odkaz{Break-up song}{break-up song}  
\item  \odkaz{Báseň od publika}{báseň od publika}  
\item  \odkaz{Dvojitá stíhačka}{dvojitá stíhačka}  
\item  \odkaz{Love song}{love song}  
\item  \odkaz{Sportovní komentátor}{sportovní komentátor}  
\item  \odkaz{Stíhačka}{stíhačka}  
\item  \odkaz{Televizní komentátor}{televizní komentátor}  
\item  \odkaz{Trojitá stíhačka}{trojitá stíhačka}  
\item  \odkaz{Ženy vs muži (kategorie)}{ženy vs muži (kategorie)}  
\end{itemize} 
\end{multicols} 
\needspace{5cm} \section{Pískání} \label{pískání} \odkaz{Rozhodčí}{rozhodčí} dává písknutím do kazoo najevo, že došlo k chybě při improvizaci a pohybem těla signalizuje, o který \odkaz{faul}{faul} došlo. 
 
 
\needspace{5cm} \section{Přebírání pozice} \label{přebírání pozice} \textbf{Přebírání pozice}{} je termín, který popisuje situaci kdy je jeden hráč ve \odkaz{štronzu}{štronzo} a jeho pozici se snaží co nejpřesněji zaujmout druhý hráč. Ten by měl zohlednit co nejvíce detailů, případně i mimiku, kterou zaujímá první hráč. Přebírání pozic se používá např. v kategoriích \odkaz{dvojitá stíhačka}{dvojitá stíhačka}, \odkaz{poslední věta}{poslední věta} ale i dalších a různých cvičeních. 
 
 
\needspace{5cm} \section{Předmět} \label{předmět} Předmět může být od diváků vybírán ve formě inspirace pro danou kategorii či cvičení a hráči jej musí během cvičení či improvizace zohlednit. Případně se jeden či více improvizátorů může předmětem stát a my tak máme možnost krom sledování práce s předmětem sledovat i jeho emoční pochody. Předmět v zastoupení hráče ve většině případů nemluví, ale není to podmínkou. 
 
Druhá forma vybírání předmětů od diváků je jejich skutečné vybrání. Diváci mají možnost zapůjčit předmět, který mají u sebe a ten se stává rekvizitou či inspirací pro improvizátory. Po kategorii se předměty divákům vracejí... většinou v pořádku, jídlo ale většinou nepřežije. 
 
Rekvizity od diváků se používají v kategoriích \odkaz{prskavka ala šanson}{prskavka ala šanson}, \odkaz{s rekvizitou}{s rekvizitou}, a další. 
 
 
\needspace{5cm} \section{Reflexe} \label{reflexe} Část po \odkaz{zápase}{zápas} či tréninku, kdy účastníci mají možnost vyjádřit své připomínky k průběhu. 
 
Mezi známé způsoby vedení reflexe patří 
\begin{itemize}
\item každý má půl minuty na to, aby řekl, co ho tíží a těší
\item všichni mají dohromady 3 minuty aby se vyjádřili, osobní zpětné vazby se pak podávají v hospodě
\item kladné kolečko - co se povedlo, záporné kolečko - rady pro příště
\end{itemize}
 
Je vhodné v rámci reflexe nezabředávat do vysvětlování a obskurních diskuzí.    
    
 
\needspace{5cm} \section{Rekvizita} \label{rekvizita} V drtivé většině divadelních her se používají rekvizity. V \odkaz{impru}{impro} se nejčastěji setkáte s imaginární rekvizitou. 
 
 
\subsection{ Imaginární rekvizita } \odkaz{Hráč}{hráč} předvádí \odkaz{pantomimou}{pantomima}, že má u sebe rekvizitu, která reálně neexistuje. Pohybem těla simuluje pohyby, které by dělal, kdyby rekvizitu opravdu měl nebo držel. V závislosti na povaze imaginární rekvizity je dobré ji na začátku pojmenovat nebo s ní provést přirozenou akci\footnote{Akci přirozenou pro rekvizitu, například vyhrožování pistolí}, aby bylo jasné, o jaký předmět se jedná. 
 
\begin{quote}"Improvizátor má u sebe vždycky všechno"\end{quote} 
 
\subsection{ Personifikovaná rekvizita } \odkaz{Hráč}{hráč} předvádí samotnou rekvizitu. Může se v očích diváka stát konkrétním předmětem. Tento druh rekvizity přináší své výhody.  
 
\begin{itemize}
\item  S takovou rekvizitou se dá reálně pohybovat, nedá skrz ni projít
\item  \odkaz{Hráč}{hráč} hrající rekvizitu může předmětu přidávat emoce a výrazy, pokud je předmět uveden v činnost
\end{itemize}
 
Často k vidění v kategorii \odkaz{duety}{duet}. 
 
Podobně v kategorii \odkaz{Živé rekvizity}{živé rekvizity} je divák používán jako rekvizita. 
 
\subsection{ Zástupná rekvizita } Reálný předmět, který v divadelní hře představuje jiný předmět z děje. Například žlutá čepice zastupující královskou korunu. 
 
Často k vidění v kategorii \odkaz{s rekvizitou}{s rekvizitou}. 
 
\subsection{ Reálná rekvizita } Reálný předmět v příběhu. Například "auto na pódiu"{} v příběhu odehrávajícím se v autě. V improvizaci se s touto rekvizitou prakticky nesetkáte. 
 
  
 
 
\needspace{5cm} \section{Replika} \label{replika} Replika je jednotlivá věta/odpověď/reakce. 
 
\needspace{5cm} \section{Rozehřívání diváků} \label{rozehřívání diváků} Na začátku \odkaz{zápasu}{zápas} či \odkaz{improshow}{improshow} se \odkaz{MC}{mc}  snaží vybudit v divácích aktivitu a odezvu. 
Používají se například: 
\begin{itemize}
\item odstupňované učení se tleskat
\item zkušební \odkaz{vyžádání si témat}{askfor}
\item společný zpěv písně
\item vzájemné představování se, objímání či líbání diváků
\end{itemize}
 
 
\needspace{5cm} \section{Status} \label{status} Při kontaktu dvou osob neustále mimovolně probíhají drobné statusové hry, 
jedná se o to, kdo bude mít "navrch". Status ale můžete mít i vůči prostředí, jinak vstoupí do  
dveří ředitel firmy a jinak žebrák. 
 
Jakýkoli příběh oživí drobné statusové hry, či dokonce statusové výměny. 
 
 
Vysoký status bývá rozpoznatelný skrze: 
\begin{itemize}
\item vysoký tonus těla
\item pomalejší pohyby a mluva
\item pevný pohled
\item ruce volně podél těla
\end{itemize}
zatímco pro  nízký status bývají charakteristické 
\begin{itemize}
\item vyhýbavý těkavý pohled
\item klopení hlavy
\item nervózní chichotání
\item neurotická práce s rukama - točení mlýnku palci, upravování si oděvu
\end{itemize}
 
\needspace{5cm} \section{Střih} \label{střih} Cítí li některý z hráčů, že je nutné děj někam více posunout, obvykle tak činí nahlas pronesenou scénickou poznámkou  - např. O hodinu později na letišti.  nebo Ukažme si, jak se kdysi seznámili. 
 
V \odkaz{longformách}{longforma} je množství používaných uznávaných metod, kterými  lze ukončit scénu, podstatně širší. 
Jedná se např. o: 
 
\begin{itemize}
\item scénická poznámka
\item monolog postavy
\item mazání scény rychlým průchodem před ostatními hráči
\item opakování pohybu
\item opakování \odkaz{repliky}{replika}
\item rozehrátí nové scény bez kontaktování ostatních hráčů
\end{itemize}
 
Další změny scény může iniciovat osvětlovač či \odkaz{hudebník}{hudebník}. 
 
 
 
\needspace{5cm} \section{Trestný bod} \label{trestný bod} Trestný bod je udělován na \odkaz{zápasech}{zápas} jako forma trestu za \odkaz{fauly}{faul}. Typicky se vizualizuje připnutím kolíčku k tabulce s počtem bodů daného \odkaz{týmu}{tým}. Za každé tři trestné body je připočítán jeden bod do celkového skóre soupeřícího týmu. 
 
 
\needspace{5cm} \section{Trénink} \label{trénink} O improvizaci říkáme, že ji trénujeme, nikoli zkoušíme, protože se, narozdíl od klasického divadla 
nejedná o neustálé zkoušení/vylepšování konkrétního obsahu, ale 
o soustavnou práci na vylepšování vlastních schopností.  
 
\needspace{5cm} \section{Ukazování času} \label{ukazování času} Při \odkaz{zápasech}{zápas} je jednou z hlavních úloh \odkaz{pomocných rozhodčí}{pomocný rozhodčí} měření času a jeho ukazování \odkaz{hráčům}{hráč} a \odkaz{hlavnímu rozhodčímu}{rozhodčí}. 
Pro toto ukazování je několik ustálených gest.  
 
\subsection{Polovina času} Ruce zkřížené v zápěstí. 
 
\subsection{10 vteřin do konce} Kroucení dlaní s roztaženými prsty. 
 
\subsection{5 (4,-1) vteřin do konce} Ukazuje se rukou se zvednutým patřičným počtem prstů. 
 
\subsection{Konec} Kroucení rukou zatnutou v pěst 
 
 
 
 
\needspace{5cm} \section{Veřejná rozcvička} \label{veřejná rozcvička} Na začátku zápasu se hráči rozehřívají před diváky několika jednoduššími cvičeními. Cvičení pomáhá seznámit se s prostorem, světlem, jevištěm i diváky. 
Diváci si naopak také zvykají na hráče, kdy mohou pocítit již první prvky improvizace. 
Existuje několik známějších cvičení, se kterými se může divák i hráč na zápasech potkat. 
 
Moderátor by měl divákům vždy představit pravidla cvičení, aby diváci nezískali pocit, že je vše připravené a mohli se soustředit na obsah, ne na hledání porozumění pravidel. 
 
Seznam cvičení na veřejnou rozcvičku naleznete v kategorii \odkaz{Rozcvičky}{:kategorie:rozcvičky}. 
 
 
\needspace{5cm} \section{Vteřinový potlesk} \label{vteřinový potlesk} \textbf{Vteřinový potlesk}{} je zvláštní formou \odkaz{hlasování}{hlasování}. Diváci v tomto případně nemají hlasovací kartičky a k hlasování používají jen svůj hlas či jiné hlasité projevy. Průběh hlasování je v tomto případě následující: \odkaz{MC}{mc} určí, pro co či jaký tým se bude právě hlasovat, poté odstartuje vteřinový potlesk, během něhož mohou diváci křičet, tleskat, dupat a jinak se projevovat následně MC potlesk ukončí zpravidla gestem a hlasuje se pro druhou variantu či tým. \odkaz{MC}{mc} poté rozhodne, který potlesk byl hlasitější a na jeho základě je rozhodnuto. Toto hlasování je dobré při velmi krátkém čase na hlasování, velkém počtu diváků nebo v případě zapomenutí hlasovacích kartiček. 
 
\subsection{ Varianta } \begin{itemize}
\item Potlesk je měřen decibelometrem a rozhoduje vyšší hodnota.
\end{itemize}
 
 
 
\needspace{5cm} \section{Štronzo} \label{štronzo} Slovní pokyn, občas nahrazovaný písknutím na píšťalku. Hráči se ve chvíli, kdy uslyší slovo \textif{stronzo}{}, zastaví a nemohou (ano, jsou to hráči, takže vlastně nechtějí, i když po chvíli hodně chtějí) se hýbat. Cvičení používané jak na fyzickou zdatnost, tak na zklidnění hráčů a jejich koncentraci. Obvyklým výstupem je následný popis \textif{živých obrazů}{}. Vydržení ve \textif{stronzu}{} je základním znakem síly týmu. A sadismu rozhodčího či trenéra. 
 
Na zápase uvidíte ztuhnutí nejčastěji při rozcvičce, jako funkční prvek se používá  v \odkaz{živých obrazech}{živé obrazy}, \odkaz{pyramidě}{pyramida}, \odkaz{poslední větě}{poslední věta} a \odkaz{čtverci}{čtverec} . 
 
Opakem \textif{štronza}{} je \textif{portamento}{}, tedy rozpohybování se.  
 
 
 
\subsection{ Poznámky }  
Pozor při překladech a komunikaci se zahraničními skupinami. V některých jazycích může být výraz \textif{stronzo}{} \textbf{urážlivý}. 
 
 


\chapter{Příběh}
\documentclass[main.tex]{subfiles}
 
\begin{document}\needspace{5cm} \label{příběh} Příběhem rozumíme popis nějakého děje, události. 
 
Historicky známě několik osnov příběhu: 
 
\begin{enumerate}
\item Začátek
\item Prostředek
\item Konec
\end{enumerate}
 
Nebo také klasické drama 
\begin{enumerate}
\item \odkaz{Expozice}{expozice}
\item \odkaz{Kolize}{kolize}
\item \odkaz{Krize}{krize}
\item \odkaz{Peripetie}{peripetie}
\item \odkaz{Katarze}{katarze}/katastrofa/závěr
\end{enumerate}

\end{document}
%\documentclass[main.tex]{subfiles}
 
\begin{document}\needspace{5cm} \section{Expozice} \label{expozice} {{Fáze příběhu}} 
 
\textif{Expozice}{} je úvod \odkaz{příběhu}{příběh}, který většinou bývá na začátku\footnote{Některé příběhy mohou mít opožděnou expozici a nebo ji naleznete až na konci příběhu. viz kategorie \odkaz{Shakespeare pozpátku}{shakespeare pozpátku}}. Obecně se expozicí vytváří předpoklady k pochopení příběhu. Zpravidla to znamená nastavení \odkaz{prostředí}{prostředí} představení hlavních \odkaz{postav}{postava} a nastavení vztahů mezi postavami. 
 
\subsection{ Běžné chyby }  
\subsubsection{ Prolínání úvodních fází příběhu } Velice běžná chyba, se kterou se potýká každá třetí improvizace. Například pokud se stane, že při expozici definujeme nějaký příběhový problém. Expozice není dokončena, vyprávění poskočí do další fáze a vzniká zmatek na pódiu. 
 
\subsubsection{ Příliš mnoho detailů } U \odkaz{krátkých forem}{:kategorie:krátké formy} improvizace by expozice příběhu neměla být zdlouhavá. Vypravěč nebo vypravěči by se měli držet \odkaz{improvizačního trojzubce}{improvizační trojzubec}. Často se stává, že vypravěči nadefinují více \odkaz{rekvizit}{rekvizita} nebo \odkaz{postav}{postava}, než které je možné v příběhu zahrnout. Výsledkem je zmatený divák, který nechápe, kam některé rekvizity nebo postavy zmizely. 
 
  
 
 
 
\needspace{5cm} \section{Kolize} \label{kolize} {{Fáze příběhu}} 
 
\textif{Kolize}{} je fáze \odkaz{příběhu}{příběh}, při které dochází k zařazení problému, konfliktu nebo rozporu do hlavní příběhové linky. Cílem kolize je vyvolání napětí a určení, kterým směrem se celý \odkaz{příběh}{příběh} bude dále vyvíjet. Předpokladem pro funkční kolizi je definování \odkaz{postav}{postava}.\footnote{\odkaz{Postavy}{postava} se definují obvykle v \odkaz{expozici}{expozice}.} 
 
\subsection{ Běžné chyby }  
\subsubsection{ Okamžité řešení } V \odkaz{krátkých formách}{:kategorie:krátké formy} by se nemělo stát, že se v \odkaz{příběhu}{příběh} objeví problém a okamžitě se vyřeší, protože takový problém nemá žádnou příběhovou hodnotu. Výsledkem okamžitého řešení je ztráta 20 vteřin, protože se v \textif{kolizi}{} musí pokračovat, dokud se nějaký konflikt nebo problém nenajde. V \odkaz{delších improvizačních formách}{:kategorie:delší formy} to smysl má, protože je zajímavá a příběhově hodnotná reakce postav na problém, který se řeší sám. 
 
\subsubsection{ Povrchní konflikt } Po \textif{kolizi}{} zpravidla přichází další fáze \odkaz{příběhu}{příběh}, kde se s problémem, což je výstup \textif{kolize}{}, nějakým způsobem pracuje. Pokud nemá problém nalezený v této fázi nějakou hloubku, může se stát, že další fáze \odkaz{příběhu}{příběh} budou působit povrchně. Způsobů, jak nastavit "hloubku"{} je více a mohou se kombinovat: 
 
\begin{itemize}
\item  Problém by neměl nepřesahovat \odkaz{žánr}{:kategorie:žánry} \odkaz{příběhu}{příběh}
\item  Postavy by měly mít citovou vazbu k problému, "stává se to osobním"
\item  Problém se dotýká přímo publika - například situace, kterou divák řeší nebo řešil
\end{itemize}
 
{{todo|Pomozte prosím doplnit další způsoby, jak postavám pořádně okořenit problém}} 
 
  
 
 
 
\needspace{5cm} \section{Krize} \label{krize} {{Fáze příběhu}} 
 
Zjednodušeně je \textif{krize}{} prostor v \odkaz{příběhu}{příběh} pro "pořádné pohrabání se v problému", který byl vytvořen v jiné fázi \odkaz{příběhu}{příběh}. Divácky je to nejdůležitější fáze, protože už jsou nastaveny všechny předpoklady pro vžití se do děje. Zároveň je to fáze, ve které dostávaji hlavní postavy hodně prostoru. 
 
 
\subsection{ Průběh } Příběh v této fázi popisuje a rozvádí, čím vším musí hlavní hrdinové projít, než vyřeší problém, tedy výstup \odkaz{kolize}{kolize}. Zpravidla zde dochází k zhoršování problému a oddalování řešení. V různých žánrech funguje prohlubování problému jinak, ale obecně to v divácích vyvolává nějakou empatii k postavě, která se buď stává bezmocnou a nebo jí někdo jen komicky zlomyslně ubližuje. Pro diváky '''jsou nejdůležitější reakce hlavní postavy''' na vývoj problému a každé rozhodnutí, které udělá.  
 
 
\subsection{ Běžné chyby }  
\subsubsection{ Úplné přeskočení fáze } \odkaz{Krátké formy}{:kategorie:krátké formy} vedou improvizátory ke spěchu a tak se občas stane, že improvizátor skočí od zadání problému ihned k jeho řešení a příběh ztratí hodnotu, která by byla vytvořena \textif{krizí}{}. Tato chyba se dá opravit \odkaz{střihem}{střih} nebo posunutím času \odkaz{příběhu}{příběh}. V \odkaz{delších formách}{:kategorie:delší formy} k této chybě zpravidla nedojde, protože \textif{krize}{} tvoří více jak polovinu \odkaz{příběhu}{příběh}. 
 
\subsubsection{ Neschopnost práce s výstupem kolize } Kolikrát už jste viděli, že problém který byl nastaven \odkaz{kolizí}{kolize}, se nakonec v \odkaz{příběhu}{příběh} vůbec neobjevil? A kolikrát jste na něj, improvizátoři, úplně zapomněli? Problém se v této fázi nesmí ani zapomenout, ztratit, vyřešit, nic. Pouze prohloubit a rozvíjet. 
 
\subsubsection{ Zbytečná expozice } Velice často potkáváme improvizované příběhy, ve kterých se při krizi začnou objevovat další \odkaz{postavy}{postava} nebo \odkaz{rekvizity}{rekvizita}, které nebyly v expozici vůbec zmíněny. Je to přirozený proces - improvizátor si řekne, že by se v \odkaz{příběhu}{příběh} hodilo letadlo, tak ho tam doplní. Chyba, která se velice často opakuje je \textbf{přehnání}{} této vnořené \odkaz{expozice}{expozice}, kdy improvizátoři stráví polovinu času \textif{krize}{} definováním nových \odkaz{postav}{postava} nebo \odkaz{rekvizit}{rekvizita}. 
 
\subsubsection{ Ztráta na váze } V této fázi je opravdu nejdůležitější prohlubování problému. Pokud \odkaz{hráč}{hráč} hrající hlavní postavu nedokáže adekvátně zareagovat na situaci, problém najednou ztratí na váze a celý \odkaz{příběh}{příběh} se stane povrchním. 
 
\begin{quote} 
"Jsem tvůj otec!" 
"Tak mi to podrž." 
\end{quote} 
 
 
 
 
\needspace{5cm} \section{Peripetie} \label{peripetie} {{Fáze příběhu}} 
 
\textif{Peripetie}{} je fáze příběhu, ve které hlavní \odkaz{postavy}{postava} po dlouhé \odkaz{krizi}{krize} nacházejí možnosti řešení. V této fázi může dojít u \odkaz{postav}{postava} k takzvanému "change of heart"\footnote{Anglická fráze přeložitelná jako fundamentální změna názoru}. 
 
\subsection{ Průběh } V této fázi příběh nabízí hlavním \odkaz{postavám}{postava} způsoby, jak se dostat z \odkaz{krize}{krize}, čímž je zároveň definována podoba \odkaz{katarze}{katarze} celého \odkaz{příběhu}{příběh}. Mělo by dojít k urovnání všech příběhových referencí a přípravě k rozuzlení celého děje. Hlavní postava se dostává informace, rekvizity nebo které chyběly k vyřešení problému. 
 
Pro diváka je v této fázi důležitá \textbf{změna}{} v myšlení nebo chování hlavních postav, které k nabídnutým řešením zaujmou nějaký postoj. 
 
\subsection{ Běžné chyby }  
Klasické chyby postav, jako nedůvěryhodnost role nebo nespojení informací dohromady, dokáží anulovat tuto fázi. 
 
  
 
 
\needspace{5cm} \section{Katarze} \label{katarze} {{Fáze příběhu}} 
 
\textif{Katarze}{} je zakončení příběhu. Divák se dozvídá, jak \odkaz{příběh}{příběh} končí. 
 
\subsection{ Průběh }  
V ideálním případě by se všechny \odkaz{příběhové}{příběh} reference (výstup \odkaz{krize}{krize} a \odkaz{peripetie}{peripetie}) v této fázi měly spojit dohromady a mělo by dojít k nějakému řešení výstupu \odkaz{kolize}{kolize}. Je více přístupů, jak nahlížet na závěr \odkaz{příběhu}{příběh}. 
 
\subsubsection{ Dobro a zlo } Dopadne to dobře nebo špatně ve vztahu k hlavním \odkaz{postavám}{postava}? Skončí to katastrofou a nebo rozuzlením? Dokáže se hrdina poprat s problémem a nebo spáchá sebevraždu? Obojí je řešení problému. 
 
\subsubsection{ Monolog } Hlavní \odkaz{postava}{postava} nebo \odkaz{vypravěč}{vypravěč} vypráví, co bylo dál. Typické pro \odkaz{vyprávěnou}{vyprávěná} a \odkaz{noir}{noir}. 
 
\subsubsection{ Otevřený konec } Pokud hrdinové získají všechny informace nebo \odkaz{rekvizity}{rekvizita}, které potřebovali a dá se logicky vydedukovat přímé pokračování příběhu, dá se skončit otevřeně. 
 
\subsubsection{ Láska za každou cenu } Hlavní postava zjistí, že konflikt nebo problém se dá vyřešit tím, že se změní. 
 
\subsubsection{ Level up } Hlavní postavy obrátí celou krizi v osobní prospěch. 
 
\begin{itemize}
\item  \begin{quote}"Sice je válka, ale já na tom vydělám."\end{quote}
\item  \begin{quote}"Moje žena stejně byla proti polygamii."\end{quote}
\end{itemize}
 
\subsubsection{ Krveprolití } Brutální odstranění všech problémových postav, typické pro \odkaz{western}{western} a \odkaz{červenou knihovnu}{červená knihovna}. 
 
 
\subsection{ Běžné chyby }  
\subsubsection{ Není čas! } V \odkaz{krátkých formách}{:kategorie:krátké formy} se může stát, že na závěr příběhu nezbyde čas. V takovém případě může buď \odkaz{MC}{mc} čas přidat a nebo se dá volně přejít do \odkaz{monologu}{katarze}. 
 
\subsubsection{ Neukončené příběhové linky } Pokud příběh obsahuje více linek, je mnohem zajímavější dohrát postupně všechny, než ukončit jenom jednu. Po konci jedné linky by se mělo přepnout do prostředí, kde se odehrává další a zahrát její \textif{katarzi}{}. 
 
 
 
\end{document}


\chapter{Longformy}\label{longformy}
\needspace{5cm} \label{longforma} Longforma/longformy se používá jako označení delších \odkaz{forem}{forma}, kdy se nejedná 
o sled krátkých nezávislých scének v \odkaz{kategoriích}{kategorie} vybíraných \odkaz{MC}{mc},  
ale o delší sled scén (nebo jedna velká dlouhá scéna). 
 
Možnosti  a varianty longforem jsou široké, od \odkaz{Harolda}{harold}, či \odkaz{Motelu}{motel} s pevně danou strukturou 9 scén ze tří příběhů, 
přes \odkaz{Bar}{bar (forma)} s pevným začátkem a volnějším koncem,  \odkaz{Armanda}{armando} s freeformami přerušovanými monology hlavního vypravěče,   
až po \odkaz{freeformy}{freeformy} obvykle omezované pouze způsobem \odkaz{střihu}{střih} mezi scénami v neomezeném množství příběhu. 
 
V \odkaz{seznamu longforem}{:kategorie:longformy} se můžete podívat na seznam těch  zde zatím popsaných. 
 
 

\documentclass[main.tex]{subfiles}\begin{document}
\needspace{5cm} \section{Armando} \label{armando} Tato \odkaz{longforma}{longforma}, zpravidla se hraje dvacet až čtyřicet minut. 
 
Jeden z hráčů získává titul  Armando. Hráč je na jevišti sám za sebe, nehraje žádnou roli. Přichází na  scénu a žádá od diváků nějaké téma, větu, otázku, cokoliv o co požádá. 
Na základě námětu od diváků vypráví pravdivý monolog ze svého života  - příběh, postřeh apod. 
 
Pro ostatní hráče je tento monolog inspirací pro rozehrání několika různých scének - příběhů. Scénku zpravidla rozehrávají 2 hráči, ale může se zapojit více lidí. 
 
Hráči mohou zastávat v různých scénkách různé role, mohou zahrát jednu postavu ve více scénkách, důležité je udržet logiku rozehraných příběhů. 
Cyklus Armandův monolog/scény jím inspirované se může několikrát opakovat. 
 
Edity mezi scénky zpravidla zajištuje jeden hráč, který i rozehrává scénku následující. Edit vznikne tím, že hráč viditelně přeběhne jeviště bez očního kontaktu s ostatním hráči. 
 
 
\needspace{5cm} \section{Bar (forma)} \label{bar (forma)} \textbf{Bar}{}  (též \textbf{Věci zapomenuté v baru}{}) je druh \odkaz{longformy}{longforma}. Hráči na začátku nejprve představují předměty zapomenuté v baru, které si berou jako inspiraci pro vytvoření postavy. Poté už jako postavy hrají scénky tvořící jeden příběh. Mezi scénkami si hráči přepínají sami pomocí \odkaz{střihů}{střih}. \textbf{Bar}{} probíhá v třech fázích, kterými celá forma plynule prochází. 
 
\subsection{ Potřeby } \begin{itemize}
\item  \odkaz{MC}{mc}
\item  Hudba nebo hudebník
\item  5-7 hráčů
\end{itemize}
 
\subsection{ Průběh } \odkaz{Hráči}{hráč} začínají \odkaz{zaplňováním prostoru}{zaplňování prostoru} a nalaďují se na sebe, dokud nepřejdou ve \odkaz{štronzo}{štronzo}. To jim může určit \odkaz{MC}{mc} nebo se hráči sami zastaví.  
Funkční je sekání příbhu na části a střídání scén s použitím hudby. 
 
\subsubsection{ Fáze předmětů } \odkaz{Hráči}{hráč} jsou stále ve \odkaz{štronzu}{štronzo} a pojmenují se, co jsou za předmět (zapomenutý v baru) podle polohy, ve které skončili. Podobně jako v \odkaz{Jsem a beru}{jsem, beru}. Po představení všech předmětů opět zaplňují prostor. 
 
\subsubsection{ Fáze monologů } Při zaplňování prostoru postupně vyjdou hráči do popředí a mají monolog na dané téma ve třech kolech. Během chození po prostoru si hráč v hlavě utváří postavu. Dohromady každá postava řekne tři monology. Každý hráč musí svojí postavu nějak popsat, ale obecně platí, že čím méně toho řekne, tím lépe. Příliš mnoho konkrétních informací, může zablokovat nápady spoluhráčů. V ideálním stavu si hráč vybere jednu věc, která je pro tu postavu nejdůležitější a o té něco poví. 
 
\subsubsection{ Vztah k předmětu } V prvním monologu postava mluví o vztahu k předmětu, který "nechala v baru".  
 
\subsubsection{ Něco o sobě } V druhém monologu mluví postava krátce o sobě. Velice dobrý nápad je říct jméno své postavy. 
 
\subsubsection{ Vztah k ostatním } Ve třetím monologu postava mluví o vztahu k ostatním postavám. Pokud předtím někdo zmínil vztah k nějaké postavě, je dobré se toho vztahu okamžitě ujmout, aby se na něj nezapomnělo, ale není potřeba aby všechny postavy měly už od začátku jasný vztah k jiné. Opět to může být svazující pro spoluhráče a může být i oříšek se pak ve vztazích vyznat. 
 
\subsubsection{ Fáze příběhu } Po skončení monologů se \odkaz{hráči}{hráč} rozejdou na strany scény. Příběh se pak sestává ze série scének, které vždy rozehraje několik hráčů, kteří vyjdou na scénu. Podoba scénky není nijak omezena, ale každý hráč si po celou dobu drží svojí postavu a také se pracuje s tím, co už zaznělo. Scénku může uzavřít \odkaz{MC}{mc}, hráč, který zrovna nehraje a čeká na straně, ale mohou ji uzavřít i hráči, kteří právě hrají. Příběh by měl být logicky ukončen, což hlídá \odkaz{MC}{mc}. 
 
 
 
\needspace{5cm} \section{Freeformy} \label{freeformy} Improvizace hraná bez omezení daných \odkaz{kategoriemi}{kategorie}, žánry, či speciálními pravidly. V tomto typu \odkaz{longform}{longforma}  jsou \odkaz{hráči}{hráč} obvykle dohodnuti pouze na způsobu ukončování scén, ani ten však nemusí dodržovat, pokud to v dané chvíli cítí jinak. Stává se, že jednotlivé scény mají souvislost, třeba ve vystupujících postavách. 
 
 
\needspace{5cm} \section{Lavička} \label{lavička} Hraje 4 - 5 hráčů. Uprostřed jeviště se postaví lavička, nebo židle vedle sebe, nebo pohovka, nebo tak něco. Přichází první hráč a promlouvá delší monolog své postavy. Na cinknutí konferenciéra se v sále zhasne, na opětovné cinknutí se rozsvítí, ale to už na lavičce sedí dva, jednoho už známe, druhý je neznámý. Ti dva spolu odehrají krátkou improvizaci, ve které zjistíme jejich vztah a vazby a pak se opět setmí. V další scéně první hráč chybí a je tam pouze druhý, který pronese monolog své postavy. Pak se opět setmí a na jevišti se vedle druhého objevuje třetí a odehrává s ním krátkou improvizaci. Pak má opět samostatný monolog a tak dále, až se představí postupně všichni hráči. Pak jdou všichni na stranu a přicházejí k lavičce vždy ve dvojici a rozehrávají společný příběh. Konferenciér jej vždy ukončuje cinknutím a zhasnutím. 
 
\subsection{ Tipy }  
V ideálním případě by divák vůbec neměl vidět výměny hráčů, jejich příchody a odchody by se měly dít za úplné tmy. 
 
 
 
\needspace{5cm} \section{Lodička} \label{lodička} Hraje obvykle 4 - 5 hráčů, kdy všichni stojí blízko sebe, jakoby stáli na malé, vratké lodičce. Konferenciér nechá diváky určit prostředí a poté se postupně všichni hráči představí krátkým monologem své postavy. Poté diváci určují, kdo je nejméně významná postava a ta řekne své jméno a odchází z lodičky. Takto se vyřadí všichni hráči, až zbyde ten nejvýznamnější. Poté se odehrává dlouhá improvizace, ve které se postupně potkávají jednotlivé postavy s tím, že nejvíce prostoru a největším hybatelem děje by měl být ten, kdo byl diváky nejvýznamnější, a poté další a další, podle toho, jak je diváci vyřazovali. Délku setkávání (jednotlivých dialogů) si můžou hráči určovat buď sami, nebo je může stopovat konferenciér. 
 
 
 
\needspace{5cm} \section{Motel} \label{motel} V malém motelu uprostřed Ameriky je pokoj, který má zvláštní genius loci. 
\subsubsection{Scéna} Na scéně stojí např. židle, noční stolek, věšák a postel, ideálně stylově odpovídající ošuntělému hotýlku kdesi v Americe. 
\subsubsection{AskFor} Od diváků se vyžádá jedno slovo.  
Všech 6 hráčů začnou "uklízet"{} scénu, urovnávat, oprašovat apod., s opakováním tohoto klíčového slova, nakonec na scéně zůstávají jen dva hráči, kteří začínají prvním výstupem prvního dějství. 
Každý hráč si drží postavu a dvojici po celé představení. 
\subsubsection{Dějství} Každé dějství se sestává ze tří výstupů, výstup končí hráči na stranách novým "uklízením"{} scény - jejím návratem do původního stavu. 
\subsubsection{Mezihra} Po skončení prvního a druhého dějství postupně předstoupí z každé dvojice jedna postava a sdělí divákům v \odkaz{monologu}{monolog} nějaké tajemství, které ten druhý nezná. 
\subsubsection{Konec} Po skončení 3. dějství a úklidu scény hráči scénu neopouští, předstupují dopředu, na klaněčku. 
 
 
---- 
''Tento formát vycházející z \odkaz{Harolda}{harold} vyvinul \textbf{ Jstar Atlanta }'' 
 
\needspace{5cm} \section{Sedm slov (Longforma)} \label{sedm slov (longforma)} \subsection{ Průběh } Hraje opět 4 - 5 hráčů a tato longforma má několik oddělených částí 
 
\subsection{ Popis prostředí } Začínáme popisem prostředí a to tak, že hráči stojí v půlkruhu a postupně jeden po druhém přicházejí před diváky a popisují místo, kde se bude příběh odehrávat (začínají od nejvzdálenějšího obrazu a neustále se přibližují a konkretizují to místo). Je důležité, aby se popisovalo, jaké to místo je, jak to tam vypadá, jakou má náladu, ale vyvarujte se popisování osob, které se tam nacházejí. Toto se děje několikrát, až jsme všichni v jednom konkrétním prostředí, které si ideálně všichni představují stejně. 
 
\subsection{ Představení postav } V druhé části chodí hráči po prostoru jeviště a představují své postavy v namalovaném prostředí a to tak, že se vždy zastaví (s nimi se zastaví i ostatní) a monologem směrem k divákům řeknou něco o tom, jaká je jejich postava. Poté pokračují v chození tak dlouho, než se nepředstaví všichni hráči. 
 
\subsection{ Drby } Pak nastává fáze drbů, hráči opět stojí v půlkruhu a po jednom předstupují před diváky a prozrazují pikantnosti ze života všech postav a tím nám lépe osvětlují jejich vzájemné vztahy. V tuto chvíli je možné začít hrát i \odkaz{Lavičku}{lavička}. 
 
\subsection{ Příběhová část } Ve fázi příběhu vždy přicházejí dva hráči společně na scénu a dohromady vždy musí použít pouze 7 slov. Konferenciér jim ideálně ukazuje na prstech, kolik jim ještě zbývá a po 7 slově je cinknutím posílá pryč. 
 
\subsection{ Tipy } Vzhledem k časté komplikovanosti postav a jejich vazeb je vhodné, když si konferenciér dělá poznámky a po fázi představování a drbů vždy znova pro všechny zopakuje zásadní fakta. 
 
 
 
\needspace{5cm} \section{Věřte, nevěřte} \label{věřte, nevěřte}  
 
\subsection{ Původ }  
Forma je inspirována americkým pořadem Věřte, nevěřte, který v každé epizodě představoval několik příběhů, které se zdály být velice nepravděpodobné, někdy obsahovaly reference na paranormální jevy. Na konci pořadu byl divák postaven před kvíz, kde se příběhy rekapitulovaly a mohl si sebereflektivně zopakovat zda-li byl příběh pravdivý nebo ne. 
 
\subsection{ Typický žánr }  
Mysteriózní/Komedie 
 
\subsection{ Lokace }  
Celá scéna je rozdělena ať už fyzicky nebo světelným přechodem na Lobby a Jeviště. Příběhů se obvykle hraje více. 
 
\subsection{ Lobby }  
V lobby se nachází moderátor pořadu a zároveň vypravěč následujícího příběhu. Lobby je možno vylepšit rekvizitami od diváků, naaranžovat pro téma daného večera a nebo jinak upravit, aby bylo pořadu dodána atmosféra. 
 
\subsection{ Jeviště}  
Zde herci hrají příběh. Jeviště by mělo být prázdné. Pokud je potřeba v příběhu nějaká rekvizita, dodá ji vypravěč před spuštěním příběhu, jinak se pracuje s imaginárními rekvizitami. 
 
\subsection{ Průběh }  
Pořad se nijak neuvádí, není potřeba konferenciér. Na se jevišti se objeví obraz postavený z herců, kteří budou hrát. Po krátké pauze, do které je vhodná znělka, se v Lobby objeví moderátor a začne představovat první záhadu podobně jako vypravěč otevírá příběh v kategorii \odkaz{vyprávěná}{vyprávěná}. Jako hlavní inspirace pro moderátora je určený obraz z herců a rekvizity které najde v Lobby. Je vhodné použít obraz jako "snímek z příběhu"{} který uvidíte - všechny postavy se tam tak objeví. Není to ale nutnost. Po krátkém úvodu se příběh spustí - focus se předá z Lobby na Jeviště.  
 
Průběh příběhu je volný - nejsou daná žádná pravidla. Je dobré používat scénické přechody, ale vše záleží na nastavení skupiny. 
 
Po skončení příběhu se vrací focus a případně všechny rekvizity do Lobby. Moderátor příběh dokončuje pár větami tak, aby se každý divák zamyslel nad tím co je a není možné v našem světě.  
 
Po skončení příběhu se dělá pauza, kde se mohou vystřídat moderátoři, opět je vhodná znělka a průběh jede od začátku. 
 
\subsection{ Zakončení }  
Všechny příběhy jsou improvizované, proto by kvíz o pravdivosti příběhů dopadl asi jednoznačně. Je dobré příběhy zrekapitulovat, pokud se nevymyslí nic lepšího, tak nechat diváka ať si o tom povídá v hospodě ... 
 
\subsection{ Představení záhady }  
Funkční variantou představení záhady pro tuto formu je model popsaný níže. Nenechte se jím však svazovat. 
 
\begin{enumerate}
\item  Zvol si obyčejný fyzický předmět který se dá najít doma
\item  Pojmenuj ho
\item  Řekni nahlas jednu věc, kterou ten předmět dělá
\item  Vymysli si co v tomhle příběhu ten předmět dělá jinak
\item  Představ hlavní postavu
\item  Uveď prostředí ve kterém se příběh odehrává
\item  Předej rekvizity
\end{enumerate}
 
\subsection{ Poznámky } \begin{itemize}
\item  Na moderátora jsou kladeny vysoké nároky na schopnost prezentovat před divákem, zachování postavy, uvěřitelnost a udržení focusu.
\item  Vnímavý osvětlovač představení hodně pomůže
\item  Příběhy by v rámci této formy měli končit těsně před uzavřením
\end{itemize}
 
 

\end{document}


\chapter{Knihy}\label{knihy}
%\needspace{5cm} \section{Improknihovnička} \label{improknihovnička} Ahoj, 
zdraví tě malá \textbf{Improknihovnička}{}. 
Mám v sobě několik knih v jazyce anglickém, které si můžeš půjčit, přečíst a vrátit. 
  
Aktuálně obsahuji tyto knihy (řazeno dle data získání): 
\begin{itemize}
\item Keith Johnstone: Impro for storytellers
\item John Cremer: Improv: Enjoy life and success with the power of yes
\item Patricia Ryan Madson: Improv wisdom : Don't prepare, Just show up
\item Viola Spolin: Improvisation for the Theater
\item Mick Napier: Improvise: Scene from the Inside Out
\item Charna Halpern, Del Close, Kim "Howard"{} Johnson: Truth in Comedy: Manual of Improvisation
\item Ben Hauck: Long-Form Improv: The Complete Guide to Creating Characters, Sustaining Scenes, and Performing Extraordinary Harolds
\item Tom Salinsky and Deborah Francis-White: The Improv Handbook: The Ultimate Guide to Improvising in Theatre, Comedy, and Beyond
\item Jimmy Carrane: Improvising Better: A Guide for the Working Improviser
\item Nancy Howland Walker: Instant Songwriting: Musical Improv from Dunce to Diva
\item Matt Walsh a kol.: Upright Citizens Brigade Comedy Improvisation Manual
\item Joe Samuel & Heather Urquhart: Sing it! The essential guide to musical improvised comedy
\item Jo Salas: Improvising real life
\item Jason Chin: Long-form improvisation & the Art of Zen
\item William Hall: Improv games for performers
\item Mick Napier: Behind the scenes (Improvising Longforms)
\item T.J.Jagodowski, David Pasquesi with Pam Victor: Improvisationat the speed of life
\end{itemize}
 
 
Jak postupovat? 
 
\begin{itemize}
\item  Oslovíš \odkaz{Vaška}{uživatel:vatoz}, který ti knihu půjčí.
\end{itemize}
 
Knihovnička přivítá i tvá doporučení na další rozšířování, či sponzorské dary. 
\needspace{5cm} \section{Literatura} \label{literatura} Pár tipů na knížky ze světa dramatické výchovy a improvizace. Kvalita není zaručena. Hodnoťte, pokud můžete nebo přidávejte odkazy, pokud je kniha dostupná online. 
 
Česká literatura 
\begin{itemize}
\item  Budínská, H.: Hry pro šest smyslů
\item  Propp, V. J.: Morfologie pohádky a jiné studie
\item  Ulrychová, I.: Drama a příběh
\item  Svozilová, L.: Tvořivá dramatika
\item  Štembergová-Kratochvílová, Š.: Loutkářova řeč
\item  Štembergová-Kratochvílová, Š.: Metodika mluvní výchovy
\end{itemize}
 
 
Zahraniční literatura  
 
\begin{itemize}
\item  Piers, L.: Theatresports
\item  Spolin, V.: Theatre Games
\item  Spolin, V.: Improvisation for the Theater
\item  Keith Johnstone: Impro for storytellers  (vyšlo v českém překladu)
\end{itemize}
 
Pár knih v jazyce anglickém je možno si \odkaz{půjčit}{improknihovnička} 
 
Cenné můžou být také knihy přímo o improvizaci od českých autorů 
\begin{itemize}
\item  Martin Vasquez: Staňte se mistrem improvizace
\item  Jana Machalíková: Improvizace ve škole
\item  Martin Vasquez & Mája Škvorová: Cvičení mistrů improvizace
\end{itemize}
 
či bakalářské a disertační práce 
\begin{itemize}
\item  Martin Vasquez: Principy  divadelní  improvizace  (v  prostředí  zápasů  v  divadelní improvizaci)
\item  Vít Piskala: Česká improvizační liga v roce 2013
\item  Tomáš Kováč: Koncepce a pravidla improvizační ligy jako klíčový prvek určující charakter divadelní komunikace mezi herci a diváky
\end{itemize}
 
 
 
\todo{isbn českých knih} 



\chapter{Autoři}\label{autori}
\needspace{5cm} \section{Uživatel:Just-paja} \label{uživatel:just-paja} Ajťák Improligy. 
 
Máte problém, nikdo vám nechce poradit a nebo jste se jenom ztratili v Praze, protože jste sháněli šalinu a někdo Vás poslal na tramvaj? Pište mně. 
\needspace{5cm} \section{Uživatel:VandaGabi} \label{uživatel:vandagabi} Šéfová pražského týmu KIŠe, který je pro ni rodinou. Začínala jako člen KOPRU a následně OMILu. 
\needspace{5cm} \section{Uživatel:Vatoz} \label{uživatel:vatoz} Jmenuji se \textbf{Vašek}{}. 
Jsem členem \textbf{KIŠ}{}, \textbf{VIP}{} a revizní komise \odkaz{Improligy}{čili}. 
 
Kontakt:   vatoz@seznam.cz, případně tel 728 303 155 
 
Mám pár knih, které si můžete přečíst - \odkaz{Improknihovnička}{improknihovnička} 
 
Jsem spoluautorem této \odkaz{ImproWiki}{improwiki} včetně její tištěné verze. 
 
Pořádám \textbf{Impromaraton}, \textbf{Letní improúterky}, Improsvětice, pomáhám s organizací \odkaz{Improtřesku}{improtřesk}. 
 
Vedl jsem i workshopy  - např  \textbf{Ozvěny Dublinu}{} pro VIP, \textbf{Improligový zápas}{} pro  Improstor, \textbf{Základy}{} pro Vyšíváme s Martou. 
\needspace{5cm} \section{Uživatel:VojtechKopta} \label{uživatel:vojtechkopta} \textbf{Vojta Kopta}{} - kameraman, fotograf, zvukař, osvětlovač, rekvizitář a divák v jedné osobě, který byl na trojciferném počtu zápasů.  
 
Za \odkaz{zápasy}{zápas}, \odkaz{improshow}{improshow}, \odkaz{longformami}{longforma}, ale i jinými formami improvizace chodím po pražských scénách, ale zajedu si i do Pardubic, Olomouce, Hradce Králové, Opatovic nad Labem či do jiného města, kde se právě konají. 
 
V letech 2001 - 2006 aktivním hráčem týmu '''\textbf{Kl.I.Ka}''' - Klamovští Improvizační Kaskadéři. 


\chapter{Co se jinam nevešlo}\label{co se jinam nevešlo}
%\documentclass[main.tex]{subfiles}\begin{document}
\needspace{5cm} \section{Automatické psaní} \label{automatické psaní} Automatické psaní je proces psaní bez vědomé kontroly autora. V google se vám sice po zadání tohoto hesla objeví spoustu spirituálních článku, tato metoda je však velice uvolňující a pomáhá rozvíjet vaše myšlenky a nápady. 
 
Pro automatické psaní je důležité si najít čas, příjemnou polohu sedu, tužku (nejlépe měkkou normální tužku) a čistý bílý papír.  
Tužku nechat lehce nad papírem, prostě vypnout a nechat ruku psát všechny slova, která vám začnou plynout myslí. 
V rámci cvičení na psaní je možné si jasně omezit časem, třeba na 5 minut. Někomu více vyhovuje trochu stresu. 
 
Opakovaným zkoušením se budete zlepšovat a texty půjdou snadněji. 
Nevadí, pokud budete zapisovat i citoslovce nebo nesmyslná slova. Důležité je, že všechno zaznamenáváte bez jakékoliv vaší cenzury a přemýšlení nad obsahem. 
 
S textem lze i dále pracovat. Například se jím můžete nechat inspirovat pro psaní povídek. 
Text po napsaní si přečtěte, podtrhněte si zajímavá spojení nebo vystřihněte a zkuste z nich složit nový text, který si už budete vědomě řídit. 
\needspace{5cm} \section{Emoce} \label{emoce} Emoce jsou procesy vykládané jako subjektivní zážitky libosti a nelibosti provázané fyziologickými změnami, motorickými projevy (gestikulace, mimika), stavy s různě velkou aktivností a vztahují se vůči něčemu konkrétnímu. Tyto prvky vidíme u ostatních na první pohled, známe je, protože sami pořád prožíváme různé stavy emocí. 
 
A stejně tak i každá postava na jevišti musí vyjadřovat nějakou emoci, aby byla uvěřitelná, pokud zrovna hlavním hrdinou není robot a i ten teoreticky může něco cítit. Často se stává, že hráč ukazuje na jevišti spíše svoje rozpoložení než psychický stav své postavy. Pokud se mu daří přijít v uvědomělém emočním stavu, je často omezený na základní emoce, jako je radost, smutek, vztek, strach, znechucení. Přitom škála našich emocí je daleko širší. 
 
Pokud pracujeme v improvizačních cvičení s emocemi, je dobré vědět: 
\begin{itemize}
\item  emoce má nějaký důvod
\item  u hraní extrémních emocí je důležité, aby hráči ovládali neutrální stav a uměli stavy přepínat
\item  k hraní emocí můžeme jít přes fyzické napětí, gesta, mimiku apod.
\item  sledovat pravidelné dýchání, aby nedocházelo u hráčů k hyperventilaci nebo hypoventilaci
\end{itemize}
 
Existují různě dlouhé seznamy, a je otázka co všechno můžeme považovat za emoci. 
Důležitým parametrem je však to, že emoce má nějakou intenzitu, je časově omezená a je vůči něčemu zaměřena. 
 
\subsubsection{Seznam emocí} Vybrané jsou převážně ty emoce, které se snadno interpretují na scéně: 
 
\begin{multicols}{3} 
\begin{itemize}
\item  Afektovaně
\item  Bázlivě
\item  Blouznivě
\item  Bojácně
\item  Bojovně
\item  Bolestivě
\item  Bouřlivě
\item  Citlivě
\item  Živě
\item  Děsivě
\item  Divoce
\item  Dotčeně
\item  Drsně
\item  Drze
\item  Důstojně
\item  Energicky
\item  Chladně
\item  Intimně
\item  Jásavě
\item  Láskyplně
\item  Lhostejně
\item  Mazlivě
\item  Milostně
\item  Nadšeně
\item  Napjatě
\item  Nedůvěřivě
\item  Nerozhodně
\item  Nedůvěřivě
\item  Pobouřeně
\item  Pokorně
\item  Procítěně
\item  Přísně
\item  Radostně
\item  Rezignovaně
\item  Rozčileně
\item  Roztomile
\item  Rozlobeně
\item  Sebejistě
\item  Sladce
\item  Snivě
\item  Starostlivě
\item  Stydlivě
\item  Svěže
\item  Štítivě
\item  Toužebně
\item  Trapně
\item  Unaveně
\item  Uraženě
\item  Vítězně
\item  Výbušně
\item  Vztekle
\item  Zahanbeně
\item  Zamilovaně
\item  Zatvrzele
\item  Zle
\item  Zlobně
\item  Znepokojeně
\end{itemize}
\end{multicols} 
 
\needspace{5cm} \section{Hudební nálady} \label{hudební nálady} Mnoho hudebníků hraje pořád stejně. Zde jsou některé nálady jako inspirace pro další variace: 
\begin{itemize}
\item něžně
\item "natvrdle"{} ([https://www.youtube.com/watch?v=BPFKMco8AL0&list=PLKPYI9ACDonB7jrbv_eAPdJpQVgbjWoDZ&index=2 Chicago - All That Jazz])
\item pohodově
\item vesele
\item smutně
\item těžkopádně (\textbf{Prokofiev - Dance of the Knights})
\item tajemně ([https://www.youtube.com/watch?v=EBGvtkVVas0&index=1&list=PLKPYI9ACDonB7jrbv_eAPdJpQVgbjWoDZ Marilin Manson - Resident Evil])
\item akčně
\end{itemize}
\needspace{5cm} \section{Icebreaker} \label{icebreaker} viz \odkaz{Warm-up}{warm-up} 
 
Cvičení na rozehřátí, odhození nervozity, seznámení se ve skupině. Cvičení mají zpravidla jednoduché zadání, na které se hráči musí plně soustředit. Pomáhá jim tak odhodit myšlenky z minulosti a naopak se soustředit na přítomnost. 
\needspace{5cm} \section{ImproAktiv} \label{improaktiv} Tato společenská hra je inspirována deskovou hrou Aktivity a jí podobných. Vznikla na improvizovaném setkání, s nadsázkou pojmenovém “ImproNarcis” v A Maze in Tchaiovna. 
 
{| class="wikitable" 
|\textbf{Pro:}{} 4  - 12 
|- 
|\textbf{Čas:}{} cca 1h 
|- 
|\textbf{Pomůcky:}{} měřič času, počítání bodů, karty s oblastmi 
|} 
 
Cílem hry je nasbírat co nejvíce bodů pro váš tým po předem dohodnutém počtu kol. Hráči se losem rozdělí do 2 týmů, v případě většího množství lidí, doporučujeme vytvořit více skupin. 
 
\begin{enumerate}
\item  \textbf{krok:}{} První člověk z týmu A (předvádějící) si vylosuje z karet oblast, podle které dostane od protivníka téma. Tip na oblasti najdete níže. Na výrobu dalších oblastí se meze nekladou. Vybranou oblast odhazujeme a dobráním karet končí hra. Vylosovaná oblast je veřejná.
\item  \textbf{krok:}{} Jeden člověk z týmu B vymýšlí téma, které šeptá losujícímu hráči.
\item  \textbf{krok:}{} 2. hráč z týmu B naopak vymýšlí limity improvizace, viz příklady níže, aniž by ještě v danou chvíli znal téma.
\item  \textbf{krok:}{} Hráči se musí shodnout (předvádějící a zadávající téma) , zda je věc hratelná a zda si rozumí v domluvených pravidlech.
\item  \textbf{krok:}{} Ostatní hráči z týmu A mají přesně 1 minutu, aby uhádli téma. V případě porušení pravidel (limitů) získává tým trestné body. Za 3 trestné body, dávají protihráči 1 bod.
\end{enumerate}
 
Hra končí buď a) dobráním všech karet b) po domluveném počtu kol. Hru lze protáhnout o 2. level, kdy na 1 oblast a 1 téma hrají hráči z týmu A a B zároveň. Téma předvádějí společně a je hezké, aby scéna byla založena na spolupráci. Hráč, který zadával téma pro dané kolo nehraje. 
 
\textbf{Tipy na oblasti:}{} 
předmět denní potřeby, výrobek, povolání, vymyšlená postava, reálná postava, historická událost, přísloví, název knihy, název filmu, vtip 
 
\textbf{Tipy na limity:}{} pantomima, báseň, repliky mohou začínat od 1 písmene, zpěv, scéna zahraná pozpátku, repliky je možné pronášet jen ve chvíli, kdy se hráč celý hýbe, jako moderátor zpráv, ve scéně se mohou vyskytnout pouze zvířátka, hráč se za celou dobu nesmí pohnout, improvizace v \odkaz{gibberish}{gibberish}, předvádění pouze rukama a hlasem, na židli bez dotknutí se země apod.. 
 
 
\needspace{5cm} \section{Improtřesk} \label{improtřesk} Několikadenní akce pro improvizátory, sestávající se z několika workshopů, večerních představení a dalších souvisejících aktivit, networkingu improvizátorů, konference atd. 
 
\begin{itemize}
\item 2011 - Improtřesk 2011 - oslava 10 let Improligy - DDM Klamovka
\item 2012 - Improtřesk 2012 - SVČ Lužánky
\item 2013 - Improtřesk 2013 - Divadlo do improligy, improliga do divadla - SVČ Lužánky
\item 2014 - Improtřesk 2014 - Get inspired - SVČ Lužánky
\item 2015 - Improtřesk 2015 - Sladíme se.   - KD Milevsko
\item 2016 - Improtřesk 2016 - Na stejné vlnové délce - KD Milevsko
\item 2017 - Improtřesk 2017 - Na vlastní kůži - KD Milevsko
\item 2018 - Improtřesk 2018 - Otevřeno nonstop - KD Milevsko
\end{itemize}
 
Improtřesk 2012 organizovalo SVČ Lužánky,  
Improtřesk 2013 a 2014 skupina Bafni, Improtřesk 2015 -2018  \odkaz{ČILI}{čili} a přátelé. 
 
\textbf{Web}{} - http://www.improtresk.cz/ 
 
{{todo|Kdo organizoval 2011 (Asi \odkaz{ČILI}{čili})?}} 
\needspace{5cm} \section{Jazykolam} \label{jazykolam} Jazykolamem se nazývá větné spojení, které je těžko vyslovitelné. Výhodou češtiny je, že jazykolam není zas tak těžké vytvořit. :) 
Pro improvizátory je i dobré rozvíjet cit k jazyku a vytvářet si vlastní věty a zároveň si je aspoň chvíli pamatovat tak, že je dokáží zopakovat vícekrát za sebou. 
 
 
\subsection{ České jazykolamy }  
Spoustu jazykolamů známe ze škol a jsou už trochu ohrané, tak pro inspiraci méně známé: 
 
\begin{itemize}
\item  Měla babka v kapse vrabce, vrabec babce v kapse píp. Zmáčkla babka vrabce v kapse, vrabec babce v kapse chcíp.
\item  Smrž pln skvrn zvlhl z mlh.
\item  Piksla sklapla, sklapla piksla. Kapsa splaskla, splaskla kapsa.
\item  Letěl jelen jetelem, jetelem letěl jelen.
\item  Patří rododendron do čeledi rododendronovitých či nerododendronovitých?
\item  Já rád játra, ty rád játra, ty rád játra, já rád játra, co nám brání dát si játra?
\end{itemize}
 
\subsection{ Anglické jazykolamy }  
\begin{itemize}
\item  Three witches watched three watches. Which witch watched which watch?
\item  I scream, you scream, we all scream for ice cream.
\item  Black back bat. (x3)
\item  Good blood, bad blood. (x3)
\item  Fuzzy Wuzzy was a bear. Fuzzy Wuzzy had no hair. Fuzzy Wuzzy wasn’t fuzzy, was he?
\item  Eleven elves licked eleven little liquorice lollipops.
\item  A box of biscuits, a box of mixed biscuits and a biscuit mixer!
\item  She sells sea-shells on the sea-shore.
\item  A proper copper coffee pot.
\end{itemize}
 
 
\needspace{5cm} \section{Kontaktní improvizace} \label{kontaktní improvizace} Jedná se o formu improvizace, kde se využívá tanec, pohyb a dotyk pro komunikaci mezi hrajícímí. Ve hře se pracuje s fyzickým vedením, využíváním tlaku, předáváním váhy těla, emocemi a důvěrou. V české republice nabízí  kurzy kontaktní improvizace např. sdružení Druna, kde lektoruje Zita Pavlištová a mnoho dalších center po celé české republice. 
 
Kontaktní improvizace rozvíjí citlivost vnímání sebe sama i druhého člověka. Pracuje s pozorností, koordinací pohybu, tlakem, uvolněním, gravitací apod. 
 
Základní rozdělení: 
 
\begin{itemize}
\item  jam (volná improvizace)
\item  akrobatická
\item  v plném kontaktu
\item  bez doteku s partnerem
\item  komunitní tvorba
\end{itemize}
 
\needspace{5cm} \section{Postava} \label{postava} Smyšlená osoba, identita nebo entita, kterou ztvárňují \odkaz{hráči}{hráč} tím, že ji hrají. Postava se projevuje pohybem, statusem, myšlením, důležitostí v příběhu dalšími projevy. 
 
{{todo|Doplnit projevy improvizované postavy}} 
\needspace{5cm} \section{Tým} \label{tým} \odkaz{Improligový}{čili} \textbf{tým}{} je skupina lidí věnující se společně improvizaci, většinou pod vedením \odkaz{trenéra}{trenér}, soustavně se připravují na \odkaz{zápasy}{zápas} i další formy. 
 
Aktuální verze seznamu týmů uznávaných  Českou Improvizační Ligou  \textbf{je dostupný online}. 
\needspace{5cm} \section{Výměna žárovky} \label{výměna žárovky} Primitivní činnost, často používaná jako příklad pro kategorii \odkaz{1000 způsobů jak}{1000 způsobů jak}. 
 
 
---- 
 
 
''Kolik je potřeba vývojářů Microsoftu na výměnu žárovky?'' 
 
''-- Ani jeden, Microsoft definuje tmu jako standard'' 
\needspace{5cm} \section{Výuka improvizace} \label{výuka improvizace} Jak si vyzkoušet improvizaci, potkat s novými lidmi, které mají podobný zájem, nabrat nové zkušenosti a inspiraci od lektorů? Nabízíme základní přehled, kam se obrátit. 
 
\subsection{ V České republice }  
Dále několik \odkaz{improligových}{čili} týmů kromě představení pořádá pravidelně víkendové či semestrální kurzy improvizace. 
Jsou to např: 
\begin{itemize}
\item  Brno, Praha - \textbf{Bafni}
\item  Pardubice, Hradec Králové - \textbf{Paleťáci}
\end{itemize}
 
Mimo Improligu existují také profesionální instituce, na které se můžete obrátit: 
 
\begin{itemize}
\item  Praha - \textbf{Škola improvizace} -  Ondřej Nečas, \odkaz{Vanda Gabrielová}{uživatel:vandagabi} , Pája Sedláčková, Táňa Jírová a další ; \textbf{Impro Institut}  - Martin Vasquez ; \textbf{Nandaru} - Pavel Wieser, Kateřina Nováková
\item  Pardubice - \textbf{Improvision} - Andrea Moličová a Tomáš Jireček
\item  Brno - \textbf{ImproŠpíl} - Džem Theatre (středoškolská improliga)
\item  Ostrava, Olomouc - \textbf{Improacademy} - Vladislav Kos a Alexandr Dvořák
\item  Praha - [http://www.blrtheatre.com/ Blood, Love & Rhetoric] - Dan Brown, Marc Cram, Jim High a John Poston -  vedeno v angličtině
\item  Brno - \textbf{Proraptor} - Bára Herucová, David Tchelidze, Deni Tchelidze, Jana Netíková a další
\end{itemize}
 
Dále je možné domluvit si workshop s konkrétními lektory. Ti využívají improvizační techniky zpravidla v kombinaci s jejich profesionální oblastí.  
 
 
Neposlední možností je návštěva festivalu \odkaz{Improtřesk}{improtřesk}, kde je v nabídce několik zajímavých workshopů. 
 
\subsection{ V zahraničí }  
Po celém světě probíhají improvizační festivaly. Vedle večerních představení se zpravidla setkáte s workshopy vedenými v angličtině. V Německu je to například \textbf{IMPRO} pořádané Die Gorillas, v Polsku můžete zavítat do Krakova na \textbf{IMPROFEST}. Aktuální seznam festivalů v Evropě najdete na stránkách \textbf{Improworld.com} 
 
 
Velkou výzvou může být i návštěva Improvizačních škol, které jsou nejvíce rozšířené v USA a Kanadě. Například v Chicagu, USA najdete hned dvě vedle sebe \textbf{Improv Olympic} a \textbf{Second City} pořádající mimo jiné i intenzivní letní kurzy. 
 
 
20 škol Improvizace  http://www.rantlifestyle.com/2014/01/04/20-schools-learn-improv-comedy/ 
\needspace{5cm} \section{Web Improliga.cz} \label{web improliga.cz} Web improliga.cz je realizován na publikačním systému Wordpress s na míru psaným rozšířením které 
\begin{itemize}
\item  Obsahuje správu skupin-členů (Včetně polí město, geokoordináty, web, facebook)
\item  Obsahuje správu skupin-nečlenů (pouze facebook)
\end{itemize}
Tyto dva seznamy jsou automaticky/naplánovanou úlohou (každý by měl přijít na řadu alespoň jednou denně) procházeny na seznam nových událostí skupiny a událostí skupinou sdílených. 
Nalezené události se v okamžiku prvního nasosnutí zkopírují do nové události,  na základě hledání v názvu a v popisu se jim přiřadí zařazení mezi workshopy / představení. 
Jednou za čas admin (nyní \odkaz{Uživatel:Vatoz}{uživatel:vatoz}) poopraví zařazení, případně promazává nesmysly, a doplňuje u zápasů druhý tým. 
Rozšíření dále umí vykreslit mapu (kde je u týmů psáno buď místo tréninků nebo domovská scéna) 
Nedílnou součástí je i generování polotovaru na pravidelné týdenní výpisy na facebook, které se pak dopisují, zpřesňují (primárně rozdělení zápas, improshow, longformy, site-specific, sólová improvizace, EN , vystoupení účastníků kurzu ) a publikují. 
 
Další součástí webu je child téma pro šablonu (čili obsahuje minimum kódu) které realizuje vykreslování týmu, seznamu týmů a seznamu představení. 
 
Kód je na https://github.com/vatoz/wpimprov a případně https://github.com/vatoz/wpimprov_vantage_child 
\needspace{5cm} \section{Zpomalený pohyb} \label{zpomalený pohyb} Scény jako je rvačka, souboj, milostné scény se mohou hrát ve zpomaleném pohybu. Tato forma umožňuje hráčům zvýraznit a gradovat veškeré emoce, pohyb, reakci mezi sebou. 
Umět svůj pohyb a gesta zpomalit však znamená výraznou fyzickou náročnost, propracovanou techniku v ovládání svalů a fyziologické znalosti. 
 
Na youtube existují tutoriály v angličtině, které si můžete vyhledat a nebo se můžete inspirovat moderními filmy využívají střídání rychlých a pomalých záběrů. 
 
\needspace{5cm} \section{ČILI} \label{čili} Jedná se o zkratku názvu Česká improvizační liga z.s..   
 
\textbf{ČILI}{} zaštiťuje \odkaz{zápas}{zápas} v improvizaci a je držitelem ochranné známky  \textbf{Improliga.cz}{} 
 
Členem valné hromady se může stát kdokoliv, kdo chodí do improvizačního týmu uvedeného na stránkách improliga.cz 
 
Valná hromada je svolávána jednou za rok výborem a o akci jsou informováni aktuální členové a skrze intranet Improligy i členové jednotlivých týmů. 
 
Výbor se schází častěji. 
 
Ve výkonném výboru  Čili jsou v tuto chvíli  Táňa Jírová, Martin Dočkal a \odkaz{Václav Černý}{uživatel:vatoz}. 
V revizní komisi jsou v tuto chvíli  Pavel Wieser, Jan Formánek a \odkaz{Vanda Gabrielová}{uživatel:vandagabi}. 
 
ČILI pro propagaci a komunikaci provozuje \textbf{stránky improliga.cz} (včetně subdomén) a \textbf{facebookový kanál} . 
 
Podívat se můžete na oficiální kanál \textbf{Youtube}, kde naleznete sestřihy kategorií od různých týmů. 
\needspace{5cm} \section{Zpívané kategorie} \label{zpívané kategorie} \label{:kategorie:zpívané kategorie}Seznam \odkaz{kategorií}{kategorie}, ve kterých se zpívá, obvykle za doprovodu \odkaz{hudebníka}{hudebník}. 
 
\begin{multicols}{2}\begin{itemize} 
\item  \odkaz{Barman song}{barman song}  
\item  \odkaz{Break-up song}{break-up song}  
\item  \odkaz{Kam zmizel ten starý song}{kam zmizel ten starý song}  
\item  \odkaz{Kramářská píseň}{kramářská píseň}  
\item  \odkaz{Love song}{love song}  
\item  \odkaz{Mikrofon}{mikrofon}  
\item  \odkaz{Muzikál}{muzikál}  
\item  \odkaz{Opilecká píseň}{opilecká píseň}  
\item  \odkaz{Prskavka ala šanson}{prskavka ala šanson}  
\item  \odkaz{Tříhlavý song}{tříhlavý song}  
\item  \odkaz{Vysílání FM}{vysílání fm}  
\item  \odkaz{Zpívaná}{zpívaná}  
\end{itemize} 
\end{multicols} 
\needspace{5cm} \section{Warm-upy} \label{warm-upy} Warm-up cvičení, nebo také \textbf{Icebreaker}{} či \textbf{Rozehřívačky}{}, mají za úkol hráče zaktivovat, rozehřát, dostat do dobré nálady, najít důvěru k ostatním hráčům a zkoncentrovat se. Zpravidla se tyto cvičení nepoužívají pro zlepšování speciálních dovedností. Jsou vkládány na začátek workshopů nebo před vystoupením. Dobře zabaví i během výletů, při čekání na vlak. 
 
Mají zpravidla jednoduché zadání, na které se hráči musí plně soustředit. Pomáhá jim tak odhodit myšlenky z minulosti a naopak se soustředit na přítomnost. 
\begin{multicols}{2}\begin{itemize} 
\item  \odkaz{Asociace (cvičení)}{asociace (cvičení)}  
\item  \odkaz{Bunny bunny}{bunny bunny}  
\item  \odkaz{Gordický uzel}{gordický uzel}  
\item  \odkaz{Heja hají}{heja hají}  
\item  \odkaz{Hot spot}{hot spot}  
\item  \odkaz{Kdo ukradl sušenky a ze stolu je vzal}{kdo ukradl sušenky a ze stolu je vzal}  
\item  \odkaz{Kung-fu}{kung-fu}  
\item  \odkaz{Most}{most}  
\item  \odkaz{Otázky (cvičení)}{otázky (cvičení)}  
\item  \odkaz{Popcorn}{popcorn}  
\item  \odkaz{Posílání signálu po kruhu}{posílání signálu po kruhu}  
\item  \odkaz{Samuraj}{samuraj}  
\item  \odkaz{Tvary}{tvary}  
\item  \odkaz{Zabírání kopce}{zabírání kopce}  
\item  \odkaz{Živé obrazy}{živé obrazy}  
\end{itemize} 
\end{multicols} 

\end{document}

\chapter{Bonusy}
\pagebreak
\newcommand{\btbinfo}[6]{

\ifx Z#6 \textbf{#1} \else #1 \small{  (i)} \fi   & \pageref{#2} &  \small{#3} & \small{#4} & \small{#5}\hline
}
\renewcommand{\btbinfo}[6]{
\ifx Z#6 \large{#1}   \fi

} 



\begin{multicols}{2}

 \begin{tabular}[t]{|p{4cm}|p{4cm}    |}
 \hline
 \multicolumn{2}{|c|}{ } \\
 \multicolumn{2}{|c|}{\large{Zápas}} \\
 \multicolumn{2}{|c|}{ } \\
 
  \hline
Rozhodčí:  & Konferenciér: \\
 \hline
\multicolumn{2}{|l|}{ Pomocňáci: }   \\ \hline
Hudba: & Světla: \\ \hline \hline
Organizace: & Grafika: \\ \hline \hline
Tým: &  Tým:\\ \hline
1. \small{(cpt)}  & 1.\small{(cpt)} \\ \hline
2.  & 2. \\ \hline
3. & 3. \\ \hline
4.  & 4. \\ \hline
Body:  & Body: \\ \hline

\end{tabular}
\vspace{3mm}
\documentclass[main.tex]{subfiles}\begin{document}
\btbinfo{Abeceda}{abeceda}{2 min}{2 
}{Téma 
}{I}
\btbinfo{Barman song}{barman song}{neomezeně}{neomezeně 
}{Problém pro každého hráče 
}{Z}
\btbinfo{Báseň od publika}{báseň od publika}{cca 6 min}{2 
}{Několik(4) slov pro každého hráče 
}{Z}
\btbinfo{Break-up song}{break-up song}{neomezený, 3 - 5 minut (podle počtu hráčů)}{2 
}{Dobrovolník z publika, židle 
}{I}
\btbinfo{Co tím chceš jako říct}{co tím chceš jako říct}{volně, více jak 3 minuty}{2+ 
}{Téma, prostředí 
}{Z}
\btbinfo{Červená knihovna}{červená knihovna}{2-3 minuty}{libovolný 
}{Téma 
}{Z}
\btbinfo{Čtverec}{čtverec}{neomezen}{4 
}{4 * prostředí, a 1 téma navíc 
}{Z}
\btbinfo{Dabing ABC}{dabing abc}{neomezený}{3 
}{Téma 
}{Z}
\btbinfo{Dabing filmu}{dabing filmu}{omezen ukázkou}{2-3 
}{Téma 
}{Z}
\btbinfo{Dabovaná}{dabovaná}{2 minuty}{4 
}{Téma 
}{Z}
\btbinfo{Diapozitivy}{diapozitivy}{3 minuty}{6+ 
}{Téma 
}{Z}
\btbinfo{Divadlo poezie}{divadlo poezie}{4 minuty}{4 a více 
}{Téma 
}{I}
\btbinfo{Dopis}{dopis}{30 vteřin, 3 minuty 
}{2, následně neomezeno 
}{Obvykle dvě postavy(kdo a komu píše) 
}{Z}
\btbinfo{Dopředu, dozadu}{dopředu, dozadu}{neomezený}{2 a naskakující 
}{Téma 
}{I}
\btbinfo{Duál}{duál}{2 min. 
}{2 a více 
}{Téma}{Z}
\btbinfo{Duet}{duet}{volný, do 30 s 
}{dvojice 
}{jednoslovné téma nebo předmět pro každou dvojici 
}{Z}
\btbinfo{Dům a zahrada}{dům a zahrada}{neomezený}{libovolný (ideálně 3 na každé straně) 
}{Téma, prostředí 
}{Z}
\btbinfo{Dva na pět}{dva na pět}{neomezený}{dva 
}{prostředí 
}{Z}
\btbinfo{Dva začátky}{dva začátky}{5 minut}{libovolný 
}{Téma 
}{I}
\btbinfo{Dvě repliky}{dvě repliky}{3 minuty}{libovolný (3 a více) 
}{Dvě věty pro kažého (až na jednoho) hráče 
}{I}
\btbinfo{Dvojitá stíhačka}{dvojitá stíhačka}{4x40s 
}{Dva týmy}{Téma 
}{Z}
\btbinfo{Emocionální autobus}{emocionální autobus}{neomezen}{1+3-6 
}{nic 
}{I}
\btbinfo{En face}{en face}{3 min 
}{4-6 
}{Téma 
}{Z}
\btbinfo{Facebookový status}{facebookový status}{3  minuty}{libovolný  
}{Poslední facebookový status některého diváka 
}{I}
\btbinfo{Filmové žánry}{filmové žánry}{neomezený, 5 - 15 minut}{2 + 2n 
}{Téma, několik existujících filmových a divadelních žánrů 
}{Z}
\btbinfo{Handicapovaná pohádka}{handicapovaná pohádka}{podle vybrané pohádky}{5 
}{Známá pohádka a handicap pro každého hráče 
}{I}
\btbinfo{Jestli víš, co tím myslím}{jestli víš, co tím myslím}{2  minuty}{2+  
}{Prostředí 
}{I}
\btbinfo{Kam zmizel ten starý song}{kam zmizel ten starý song}{cca 3 minuty na každou píseň}{libovolný (3 a více) 
}{Téma, povolání/specializace 
}{Z}
\btbinfo{Kaskadéři}{kaskadéři}{orientačně 3 min}{4 (dva páry) 
}{Téma 
}{I}
\btbinfo{Komiks}{komiks}{cca 1 minuta}{4 
}{Téma 
}{Z}
\btbinfo{Kramářská píseň}{kramářská píseň}{4 minuty}{libovolný (4 a více) 
}{Téma 
}{I}
\btbinfo{Krátká báseň}{krátká báseň}{volný, max minuta}{4 
}{téma či název básně 
}{I}
\btbinfo{Kvartet}{kvartet}{volný 
}{4 
}{jednoslovné téma nebo předmět 
}{Z}
\btbinfo{Loutky}{loutky}{3 minuty}{sudý počet, 4 a více 
}{Téma 
}{Z}
\btbinfo{Love song}{love song}{neomezený, 3 - 5 minut (podle počtu hráčů)}{2 
}{Dobrovolník z publika, židle 
}{Z}
\btbinfo{Metafory}{metafory}{volný, dokud nedojdou nápady}{libovolný, obvykle se zapojí všichni 
}{Dva názvy předmětů 
}{Z}
\btbinfo{Mikrofon}{mikrofon}{neomezený,volný}{libovolný 
}{Téma 
}{Z}
\btbinfo{Monolog}{monolog}{neomezený,volný}{libovolný 
}{Téma 
}{Z}
\btbinfo{Muzikál}{muzikál}{5-15 minut}{2 - 5 
}{Slovní téma 
}{I}
\btbinfo{Nátlak (kategorie)}{nátlak (kategorie)}{3 minuty, či do splnění úkolu}{2 
}{úkol pro jednoho hráče 
}{Z}
\btbinfo{Občanky}{občanky}{volný, dokud nedojdou nápady}{3-4 
}{Občanské průkazy od diváků podle počtu hráčů 
}{I}
\btbinfo{Opilecká píseň}{opilecká píseň}{volný}{4 
}{Téma 
}{Z}
\btbinfo{Oprsklá}{oprsklá}{neomezen 
}{2 
}{téma/prostředí/vztah/situace 
}{I}
\btbinfo{Orákulum}{orákulum}{volný}{všichni 
}{otázka 
}{I}
\btbinfo{Počet slov}{počet slov}{3 minuty}{libovolný (3 a více) 
}{Téma,počet pro každého hráče (ten může být přidělený) 
}{Z}
\btbinfo{Poetická}{poetická}{3 minuty}{2 
}{Téma 
}{Z}
\btbinfo{Polovina času}{polovina času}{Např 60, 30, 15, 7.5, 3.75, 1.8   
}{2 a více 
}{téma}{Z}
\btbinfo{Poslední věta}{poslední věta}{do nalezení pointy, 5-10 minut}{dvojice (2 a více) 
}{Téma a prostředí 
}{Z}
\btbinfo{Prskavka}{prskavka}{30 vteřin 
}{libovolný, každý sám za sebe 
}{jednoslovné téma pro každého hráče 
}{Z}
\btbinfo{Prskavka ala šanson}{prskavka ala šanson}{cca 40 - 90 vteřin na každého hráče 
}{libovolný, každý jednotlivě 
}{Reálné předměty z publika 
}{Z}
\btbinfo{Překlad do znakové řeči}{překlad do znakové řeči}{3 min}{libovolný (3 a více) 
}{Téma,odbornost experta 
}{Z}
\btbinfo{Příběh po slovech}{příběh po slovech}{volný}{libovolný,minimálně 2 
}{Téma 
}{I}
\btbinfo{Pyramida}{pyramida}{volný}{4-6 
}{Prostředí, Téma nebo činnost 
}{Z}
\btbinfo{Reklamace}{reklamace}{neomezený, volný}{4 
}{Co se bude reklamovat 
}{Z}
\btbinfo{Režiséři}{režiséři}{neomezen}{neomezen}{Téma pro každého z režisérů}{I}
\btbinfo{Rozhlasová hra}{rozhlasová hra}{neomezený}{všichni 
}{Téma 
}{Z}
\btbinfo{Ruce}{ruce}{cca 8 minut}{4+ (nebo 4 a MC) 
}{Bláznivý vynález 
}{Z}
\btbinfo{Ryba (ve čtyřech)}{ryba (ve čtyřech)}{3 minuty 30 vteřin}{4 
}{Prostředí 
}{Z}
\btbinfo{Scénář / Souboj s knihou}{scénář / souboj s knihou}{3 minuty}{2 
}{Připravená kniha či scénář 
}{I}
\btbinfo{Sci-fi}{sci-fi}{cca 3 minuty}{2+libovolně 
}{téma  
}{Z}
\btbinfo{Shakespeare pozpátku}{shakespeare pozpátku}{cca 10 minut}{neomezený 
}{Téma 
}{I}
\btbinfo{Schizofrenie}{schizofrenie}{2 min.}{4}{téma + charakter pro každého hráče}{Z}
\btbinfo{Smrt v 1 minutě}{smrt v 1 minutě}{60 vteřin}{2 
}{Prostředí 
}{Z}
\btbinfo{Spoon river}{spoon river}{Neomezený, bývá 5 a více minut, typicky alespoň na 3 promluvy za každého hráče. 
}{4-6 
}{Dvojslovné téma, prostředí 
}{Z}
\btbinfo{Sportovní komentátor}{sportovní komentátor}{3 minuty}{3 
}{Vrcholový sport 
}{Z}
\btbinfo{S rekvizitou}{s rekvizitou}{2 minuty}{2 
}{Reálný předmět , Prostředí 
}{Z}
\btbinfo{Statusy}{statusy}{volný}{4 
}{Téma a lístky s čísly 
}{I}
\btbinfo{Stíhačka}{stíhačka}{1:20, 40s na tým}{2, plus naskakující 
}{Téma 
}{Z}
\btbinfo{Stojí, leží, sedí}{stojí, leží, sedí}{2 min}{3}{Téma}{I}
\btbinfo{Šumlovaná}{šumlovaná}{2 min 
}{2 a více 
}{Téma}{Z}
\btbinfo{Šumlovaná se simultánním překladem}{šumlovaná se simultánním překladem}{2 min 
}{2 a více 
}{Téma}{Z}
\btbinfo{Tam a zpět}{tam a zpět}{neomezený,volný , 3 -5 minut (podle počtu hráčů)}{libovolný (4 a více) 
}{Prostředí, slovo pro každého hráče 
}{Z}
\btbinfo{Televizní komentátor}{televizní komentátor}{3 minuty}{2x3 
}{Téma, pro každé studio 
}{Z}
\btbinfo{Toaster}{toaster}{cca 6 minut  
}{4-7}{Prostředí 
}{Z}
\btbinfo{Trojitá stíhačka}{trojitá stíhačka}{6x40s 
}{Dva týmy nebo dvě dvojice}{Téma 
}{Z}
\btbinfo{Tříhlavý song}{tříhlavý song}{cca 3 min}{3}{první verš básně}{I}
\btbinfo{Tři židle}{tři židle}{5-15 minut}{3 
}{není nutné 
}{I}
\btbinfo{Věty, které by neměly zaznít}{věty, které by neměly zaznít}{neomezený}{libovolný 
}{Situace/činnosti (3 - 5) 
}{I}
\btbinfo{Věty z kapsy}{věty z kapsy}{neomezený,volný}{2 
}{Téma, prostředí 
}{Z}
\btbinfo{Vnitřní hlasy}{vnitřní hlasy}{2-3 minuty}{4 
}{Téma 
}{Z}
\btbinfo{Volná}{volná}{cca 10 minut}{2+libovolně 
}{téma 
}{Z}
\btbinfo{Vtip}{vtip}{cca 1 min 
}{alespoň dva}{nemusí být 
}{I}
\btbinfo{Vyprávěná}{vyprávěná}{2 minuty}{neomezen 
}{téma 
}{Z}
\btbinfo{Vysílání FM}{vysílání fm}{volný, obvykle 5 - 10 minut}{5 a více, obvykle se zapojí všichni 
}{Téma 
}{Z}
\btbinfo{Western}{western}{cca 3 minuty}{2+libovolně 
}{téma  
}{Z}
\btbinfo{Zemři}{zemři}{do vyčerpání}{3 - 20 
}{Téma 
}{I}
\btbinfo{Změna}{změna}{2 min}{2 
}{Téma 
}{Z}
\btbinfo{Zpívaná}{zpívaná}{neomezený,volný , 3 -5 minut (podle počtu hráčů)}{libovolný (3 a více) 
}{Téma, prostředí 
}{Z}
\btbinfo{Zvuky}{zvuky}{2-3 minuty}{2 
}{Téma 
}{I}
\btbinfo{Žánrová vyprávěná}{žánrová vyprávěná}{3 minuty}{4 
}{Téma, 4 žánry 
}{Z}
\btbinfo{Ženy vs muži (kategorie)}{ženy vs muži (kategorie)}{2*2minuty}{skupina mužů, skupina žen 
}{dvouslovné téma/situace 
}{Z}
\btbinfo{Živé obrazy}{živé obrazy}{neomezený, obvykle 2-3 minuty}{neomezený 
}{Prostředí 
}{I}
\btbinfo{Živé rekvizity}{živé rekvizity}{2 minuty}{2 
}{Prostředí, dva dobrovolníci 
}{Z}
\btbinfo{1000 způsobů jak}{1000 způsobů jak}{volný, dokud nedojdou nápady}{libovolný, obvykle se zapojí všichni 
}{Činnosti (3-5) 
}{I}
\btbinfo{4, 2, 1}{4, 2, 1}{1 min}{4 
}{Téma nebo prostředí 
}{I}

\end{document}
\end {multicols}
\renewcommand{\btbinfo}[6]{
\ifx Z#6  \else 
\large{#1 }

  \fi
}


\begin{multicols}{3}
\large{\textbf{Veřejné rozcvičky}}

\documentclass[main.tex]{subfiles}
 
\begin{document}\large{Báseň (cvičení)}

\large{Běh přes překážky}

\large{Horská dráha}

\large{Hospodská rvačka}

\large{Jsem, beru}

\large{Orchestrion}

\large{Roztočený brankář}

\large{Samuraj}

\large{Souboj otázek}

\large{Stroj}

\large{Thai Chi}

\large{Židle}
\end{document}
\end {multicols}


\begin{multicols}{3}
\large{\textbf{Kategorie na improshow}}

\documentclass[main.tex]{subfiles}\begin{document}
\btbinfo{Abeceda}{abeceda}{2 min}{2 
}{Téma 
}{I}
\btbinfo{Barman song}{barman song}{neomezeně}{neomezeně 
}{Problém pro každého hráče 
}{Z}
\btbinfo{Báseň od publika}{báseň od publika}{cca 6 min}{2 
}{Několik(4) slov pro každého hráče 
}{Z}
\btbinfo{Break-up song}{break-up song}{neomezený, 3 - 5 minut (podle počtu hráčů)}{2 
}{Dobrovolník z publika, židle 
}{I}
\btbinfo{Co tím chceš jako říct}{co tím chceš jako říct}{volně, více jak 3 minuty}{2+ 
}{Téma, prostředí 
}{Z}
\btbinfo{Červená knihovna}{červená knihovna}{2-3 minuty}{libovolný 
}{Téma 
}{Z}
\btbinfo{Čtverec}{čtverec}{neomezen}{4 
}{4 * prostředí, a 1 téma navíc 
}{Z}
\btbinfo{Dabing ABC}{dabing abc}{neomezený}{3 
}{Téma 
}{Z}
\btbinfo{Dabing filmu}{dabing filmu}{omezen ukázkou}{2-3 
}{Téma 
}{Z}
\btbinfo{Dabovaná}{dabovaná}{2 minuty}{4 
}{Téma 
}{Z}
\btbinfo{Diapozitivy}{diapozitivy}{3 minuty}{6+ 
}{Téma 
}{Z}
\btbinfo{Divadlo poezie}{divadlo poezie}{4 minuty}{4 a více 
}{Téma 
}{I}
\btbinfo{Dopis}{dopis}{30 vteřin, 3 minuty 
}{2, následně neomezeno 
}{Obvykle dvě postavy(kdo a komu píše) 
}{Z}
\btbinfo{Dopředu, dozadu}{dopředu, dozadu}{neomezený}{2 a naskakující 
}{Téma 
}{I}
\btbinfo{Duál}{duál}{2 min. 
}{2 a více 
}{Téma}{Z}
\btbinfo{Duet}{duet}{volný, do 30 s 
}{dvojice 
}{jednoslovné téma nebo předmět pro každou dvojici 
}{Z}
\btbinfo{Dům a zahrada}{dům a zahrada}{neomezený}{libovolný (ideálně 3 na každé straně) 
}{Téma, prostředí 
}{Z}
\btbinfo{Dva na pět}{dva na pět}{neomezený}{dva 
}{prostředí 
}{Z}
\btbinfo{Dva začátky}{dva začátky}{5 minut}{libovolný 
}{Téma 
}{I}
\btbinfo{Dvě repliky}{dvě repliky}{3 minuty}{libovolný (3 a více) 
}{Dvě věty pro kažého (až na jednoho) hráče 
}{I}
\btbinfo{Dvojitá stíhačka}{dvojitá stíhačka}{4x40s 
}{Dva týmy}{Téma 
}{Z}
\btbinfo{Emocionální autobus}{emocionální autobus}{neomezen}{1+3-6 
}{nic 
}{I}
\btbinfo{En face}{en face}{3 min 
}{4-6 
}{Téma 
}{Z}
\btbinfo{Facebookový status}{facebookový status}{3  minuty}{libovolný  
}{Poslední facebookový status některého diváka 
}{I}
\btbinfo{Filmové žánry}{filmové žánry}{neomezený, 5 - 15 minut}{2 + 2n 
}{Téma, několik existujících filmových a divadelních žánrů 
}{Z}
\btbinfo{Handicapovaná pohádka}{handicapovaná pohádka}{podle vybrané pohádky}{5 
}{Známá pohádka a handicap pro každého hráče 
}{I}
\btbinfo{Jestli víš, co tím myslím}{jestli víš, co tím myslím}{2  minuty}{2+  
}{Prostředí 
}{I}
\btbinfo{Kam zmizel ten starý song}{kam zmizel ten starý song}{cca 3 minuty na každou píseň}{libovolný (3 a více) 
}{Téma, povolání/specializace 
}{Z}
\btbinfo{Kaskadéři}{kaskadéři}{orientačně 3 min}{4 (dva páry) 
}{Téma 
}{I}
\btbinfo{Komiks}{komiks}{cca 1 minuta}{4 
}{Téma 
}{Z}
\btbinfo{Kramářská píseň}{kramářská píseň}{4 minuty}{libovolný (4 a více) 
}{Téma 
}{I}
\btbinfo{Krátká báseň}{krátká báseň}{volný, max minuta}{4 
}{téma či název básně 
}{I}
\btbinfo{Kvartet}{kvartet}{volný 
}{4 
}{jednoslovné téma nebo předmět 
}{Z}
\btbinfo{Loutky}{loutky}{3 minuty}{sudý počet, 4 a více 
}{Téma 
}{Z}
\btbinfo{Love song}{love song}{neomezený, 3 - 5 minut (podle počtu hráčů)}{2 
}{Dobrovolník z publika, židle 
}{Z}
\btbinfo{Metafory}{metafory}{volný, dokud nedojdou nápady}{libovolný, obvykle se zapojí všichni 
}{Dva názvy předmětů 
}{Z}
\btbinfo{Mikrofon}{mikrofon}{neomezený,volný}{libovolný 
}{Téma 
}{Z}
\btbinfo{Monolog}{monolog}{neomezený,volný}{libovolný 
}{Téma 
}{Z}
\btbinfo{Muzikál}{muzikál}{5-15 minut}{2 - 5 
}{Slovní téma 
}{I}
\btbinfo{Nátlak (kategorie)}{nátlak (kategorie)}{3 minuty, či do splnění úkolu}{2 
}{úkol pro jednoho hráče 
}{Z}
\btbinfo{Občanky}{občanky}{volný, dokud nedojdou nápady}{3-4 
}{Občanské průkazy od diváků podle počtu hráčů 
}{I}
\btbinfo{Opilecká píseň}{opilecká píseň}{volný}{4 
}{Téma 
}{Z}
\btbinfo{Oprsklá}{oprsklá}{neomezen 
}{2 
}{téma/prostředí/vztah/situace 
}{I}
\btbinfo{Orákulum}{orákulum}{volný}{všichni 
}{otázka 
}{I}
\btbinfo{Počet slov}{počet slov}{3 minuty}{libovolný (3 a více) 
}{Téma,počet pro každého hráče (ten může být přidělený) 
}{Z}
\btbinfo{Poetická}{poetická}{3 minuty}{2 
}{Téma 
}{Z}
\btbinfo{Polovina času}{polovina času}{Např 60, 30, 15, 7.5, 3.75, 1.8   
}{2 a více 
}{téma}{Z}
\btbinfo{Poslední věta}{poslední věta}{do nalezení pointy, 5-10 minut}{dvojice (2 a více) 
}{Téma a prostředí 
}{Z}
\btbinfo{Prskavka}{prskavka}{30 vteřin 
}{libovolný, každý sám za sebe 
}{jednoslovné téma pro každého hráče 
}{Z}
\btbinfo{Prskavka ala šanson}{prskavka ala šanson}{cca 40 - 90 vteřin na každého hráče 
}{libovolný, každý jednotlivě 
}{Reálné předměty z publika 
}{Z}
\btbinfo{Překlad do znakové řeči}{překlad do znakové řeči}{3 min}{libovolný (3 a více) 
}{Téma,odbornost experta 
}{Z}
\btbinfo{Příběh po slovech}{příběh po slovech}{volný}{libovolný,minimálně 2 
}{Téma 
}{I}
\btbinfo{Pyramida}{pyramida}{volný}{4-6 
}{Prostředí, Téma nebo činnost 
}{Z}
\btbinfo{Reklamace}{reklamace}{neomezený, volný}{4 
}{Co se bude reklamovat 
}{Z}
\btbinfo{Režiséři}{režiséři}{neomezen}{neomezen}{Téma pro každého z režisérů}{I}
\btbinfo{Rozhlasová hra}{rozhlasová hra}{neomezený}{všichni 
}{Téma 
}{Z}
\btbinfo{Ruce}{ruce}{cca 8 minut}{4+ (nebo 4 a MC) 
}{Bláznivý vynález 
}{Z}
\btbinfo{Ryba (ve čtyřech)}{ryba (ve čtyřech)}{3 minuty 30 vteřin}{4 
}{Prostředí 
}{Z}
\btbinfo{Scénář / Souboj s knihou}{scénář / souboj s knihou}{3 minuty}{2 
}{Připravená kniha či scénář 
}{I}
\btbinfo{Sci-fi}{sci-fi}{cca 3 minuty}{2+libovolně 
}{téma  
}{Z}
\btbinfo{Shakespeare pozpátku}{shakespeare pozpátku}{cca 10 minut}{neomezený 
}{Téma 
}{I}
\btbinfo{Schizofrenie}{schizofrenie}{2 min.}{4}{téma + charakter pro každého hráče}{Z}
\btbinfo{Smrt v 1 minutě}{smrt v 1 minutě}{60 vteřin}{2 
}{Prostředí 
}{Z}
\btbinfo{Spoon river}{spoon river}{Neomezený, bývá 5 a více minut, typicky alespoň na 3 promluvy za každého hráče. 
}{4-6 
}{Dvojslovné téma, prostředí 
}{Z}
\btbinfo{Sportovní komentátor}{sportovní komentátor}{3 minuty}{3 
}{Vrcholový sport 
}{Z}
\btbinfo{S rekvizitou}{s rekvizitou}{2 minuty}{2 
}{Reálný předmět , Prostředí 
}{Z}
\btbinfo{Statusy}{statusy}{volný}{4 
}{Téma a lístky s čísly 
}{I}
\btbinfo{Stíhačka}{stíhačka}{1:20, 40s na tým}{2, plus naskakující 
}{Téma 
}{Z}
\btbinfo{Stojí, leží, sedí}{stojí, leží, sedí}{2 min}{3}{Téma}{I}
\btbinfo{Šumlovaná}{šumlovaná}{2 min 
}{2 a více 
}{Téma}{Z}
\btbinfo{Šumlovaná se simultánním překladem}{šumlovaná se simultánním překladem}{2 min 
}{2 a více 
}{Téma}{Z}
\btbinfo{Tam a zpět}{tam a zpět}{neomezený,volný , 3 -5 minut (podle počtu hráčů)}{libovolný (4 a více) 
}{Prostředí, slovo pro každého hráče 
}{Z}
\btbinfo{Televizní komentátor}{televizní komentátor}{3 minuty}{2x3 
}{Téma, pro každé studio 
}{Z}
\btbinfo{Toaster}{toaster}{cca 6 minut  
}{4-7}{Prostředí 
}{Z}
\btbinfo{Trojitá stíhačka}{trojitá stíhačka}{6x40s 
}{Dva týmy nebo dvě dvojice}{Téma 
}{Z}
\btbinfo{Tříhlavý song}{tříhlavý song}{cca 3 min}{3}{první verš básně}{I}
\btbinfo{Tři židle}{tři židle}{5-15 minut}{3 
}{není nutné 
}{I}
\btbinfo{Věty, které by neměly zaznít}{věty, které by neměly zaznít}{neomezený}{libovolný 
}{Situace/činnosti (3 - 5) 
}{I}
\btbinfo{Věty z kapsy}{věty z kapsy}{neomezený,volný}{2 
}{Téma, prostředí 
}{Z}
\btbinfo{Vnitřní hlasy}{vnitřní hlasy}{2-3 minuty}{4 
}{Téma 
}{Z}
\btbinfo{Volná}{volná}{cca 10 minut}{2+libovolně 
}{téma 
}{Z}
\btbinfo{Vtip}{vtip}{cca 1 min 
}{alespoň dva}{nemusí být 
}{I}
\btbinfo{Vyprávěná}{vyprávěná}{2 minuty}{neomezen 
}{téma 
}{Z}
\btbinfo{Vysílání FM}{vysílání fm}{volný, obvykle 5 - 10 minut}{5 a více, obvykle se zapojí všichni 
}{Téma 
}{Z}
\btbinfo{Western}{western}{cca 3 minuty}{2+libovolně 
}{téma  
}{Z}
\btbinfo{Zemři}{zemři}{do vyčerpání}{3 - 20 
}{Téma 
}{I}
\btbinfo{Změna}{změna}{2 min}{2 
}{Téma 
}{Z}
\btbinfo{Zpívaná}{zpívaná}{neomezený,volný , 3 -5 minut (podle počtu hráčů)}{libovolný (3 a více) 
}{Téma, prostředí 
}{Z}
\btbinfo{Zvuky}{zvuky}{2-3 minuty}{2 
}{Téma 
}{I}
\btbinfo{Žánrová vyprávěná}{žánrová vyprávěná}{3 minuty}{4 
}{Téma, 4 žánry 
}{Z}
\btbinfo{Ženy vs muži (kategorie)}{ženy vs muži (kategorie)}{2*2minuty}{skupina mužů, skupina žen 
}{dvouslovné téma/situace 
}{Z}
\btbinfo{Živé obrazy}{živé obrazy}{neomezený, obvykle 2-3 minuty}{neomezený 
}{Prostředí 
}{I}
\btbinfo{Živé rekvizity}{živé rekvizity}{2 minuty}{2 
}{Prostředí, dva dobrovolníci 
}{Z}
\btbinfo{1000 způsobů jak}{1000 způsobů jak}{volný, dokud nedojdou nápady}{libovolný, obvykle se zapojí všichni 
}{Činnosti (3-5) 
}{I}
\btbinfo{4, 2, 1}{4, 2, 1}{1 min}{4 
}{Téma nebo prostředí 
}{I}

\end{document}
\end {multicols}


\pagebreak
\newcommand{\faulinfo}[5]{
#4 & #1 & #3 & #5 \\

}

\begin{tabular}[t]{|p{3cm}|p{4cm}|p{1cm}|p{6cm}    |}
\documentclass[main.tex]{subfiles}
 
\begin{document}
\faulinfo{Bránění ve hře}{bránění ve hře}{1 }{\obrazekmaly{Braneni.png}}{Rozhodčí zkříží ruce ve výši prsou}
\faulinfo{Hrubá chyba}{hrubá chyba}{2}{\obrazekmaly{Hruba-chyba.png}}{pohyb zalomenou paží nahoru a dolů}
\faulinfo{Klišé}{klišé}{1 }{\obrazekmaly{Klise.png}}{Rozhodčí klepne dlaní o kotník}
\faulinfo{Mimo rámec}{mimo rámec}{1 }{\obrazekmaly{Mimo-ramec.png}}{Rozhodčí do vzduchu kreslí rámeček}
\faulinfo{Nepovolená rekvizita}{nepovolená rekvizita}{1 }{\obrazekmaly{Nepovolena-rekvizita.png}}{Rozhodčí se tahá za cíp dresu}
\faulinfo{Nepovolený postup}{nepovolený postup}{1 }{\obrazekmaly{Nepovoleny-postup.png}}{poklepávání na paži}
\faulinfo{Nepozornost}{nepozornost}{1}{\obrazekmaly{Nepozornost.png}}{ruka držící druhou v zápěstí}
\faulinfo{Nepřehlednost hry}{nepřehlednost hry}{1}{\obrazekmaly{Neprehlednost-hry.png}}{Rozhodčí dělá rukama mlýnek}
\faulinfo{Nerespektování postavy}{nerespektování postavy}{1}{\obrazekmaly{Nerespektovani-postavy.png}}{Rozevřenou rukou (dlaní) si zakrejeme obličej}
\faulinfo{Nevhodné chování}{nevhodné chování}{2 }{\obrazekmaly{Nevhodne-chovani.png}}{Ruce v bok}
\faulinfo{Nátlak (faul)}{nátlak (faul)}{1 }{\obrazekmaly{Natlak.png}}{Rozhodčí bouchá pěstí do dlaně}
\faulinfo{Plynulost hry}{plynulost hry}{1 }{\obrazekmaly{Plynulost-hry.png}}{Rozhodčí krouží prstem nad hlavou}
\faulinfo{Přečíslení}{přečíslení}{1 }{\obrazekmaly{Precisleni.png}}{Rozhodčí se plácá dlaní do hlavy}
\faulinfo{Vypadnutí z role}{vypadnutí z role}{1}{\obrazekmaly{Vypadnuti-z-role.png}}{Ruka zaťatá s pěstí, zapumpujeme dvakrát nahoru a dolu   }
\faulinfo{Výstraha}{výstraha}{2}{\obrazekmaly{Vystraha.png}}{poklepávající ruce v bok }
\faulinfo{Šaškování}{šaškování}{1 }{\obrazekmaly{Saskovani.png}}{Rozhodčí ukazuje dlouhý nos}
\end{document}
\end{tabular}
\pagebreak
\setcounter{tocdepth}{1}
\tableofcontents
\end{document}

